
%%%%%%%%%%%%%%%%%%%%%%%%%%%%%%%%%%%%%%%%

\documentclass[12pt,amstex,amssymb]{article}
%\documentclass[oneside,openany,a4paper]{amsart}
\usepackage[margin=0.79in]{geometry}
%\documentclass[oneside, amstex,amssymb, a4paper]{amsart}
\usepackage{amsmath,amssymb,latexsym, amsfonts, amscd, amsthm}
\usepackage{graphicx,epsfig}
\usepackage{xcolor}
\usepackage{
	%dsfont, 
	%soul, 
	stackrel, tikz-cd, graphicx, etoolbox 
	%multirow
}
\usepackage{hyperref}
\usepackage{enumerate}
\usepackage{bm}


\theoremstyle{plain}
\newtheorem{theorem}{Theorem}[subsection]
\newtheorem{mtheorem}{Theorem}
\newtheorem{conjecture}[theorem]{Conjecture}
\newtheorem{mconjecture}[mtheorem]{Conjecture}
\newtheorem{proposition}[theorem]{Proposition}
\newtheorem{corollary}[theorem]{Corollary}
\newtheorem{hypothesis}[theorem]{Hypothesis}
\newtheorem{assumption}[theorem]{Assumption}
\newtheorem{lemma}[theorem]{Lemma}
\newtheorem{question}[theorem]{Question}
\newtheorem{problem}[theorem]{Problem}
\newtheorem{example}[theorem]{Example}
\newcommand{\PP}{\mathbb{P}}
\theoremstyle{definition}
\newtheorem{definition}[theorem]{Definition}
\newtheorem{remark}[theorem]{Remark}
 

\def\image{\text{Im}}
\def\And{\quad\hbox{ and }\quad}
\def\cc{completely continuous\ }
\def\ccly{completely continuously\ }
\def\fg{{\frak g}}
\def\fn{{\frak n}}
\def\ft{{\frak t}}
\def\fb{{\frak b}}
\def\fU{{\frak U}}

\def\lra{\longrightarrow}
\def\g{\gamma}
\def\cA{{\cal A}}
\def\fA{{\frak A}}
\def\B{{\mathbb B}}
\def\cB{{\cal B}}
\def\cC{{\cal C}}
\def\tcC{{\widetilde{\cal C}}}
\def\cD{{\cal D}}
%\def\Dist{\hbox{\Dist}}
\def\cE{{\cal E}}
\def\fD{{\frak D}}
\def\cR{{\cal R}}
\def\tR{{\frak R}}
\def\tF{{\frak F}}
\def\tJ{{\frak J}}
\def\tE{{\frak E}}
\def\tC{{\frak C}}
\def\cF{{\cal F}}
\def\cK{{\cal K}}
\def\tH{{\frak H}}
\def\cH{{\cal H}}
\def\tK{{\widetilde{\cal K}}}
\def\cS{{\cal S}}
\def\dV{{(V^\ast)}}
\def\fI{{\frak I}}
\def\fS{\frak S}
\def\fH{{\frak H}}
\def\cL{{\cal L}}
\def\cM{{\cal M}}
\def\cO{\mathcal {O}}
\def\cX{{\cal X}}
\def\Z{{\mathbb Z}}
\def\R{{\mathbb R}}
\def\bP{{\mathbb P}}
\def\D{{\mathbb D}}
\def\A{{\mathbb A}}
\def\C{{\mathbb C}}
\def\G{{\mathbb G}}
\def\N{{\mathbb N}}
\def\Q{{\mathbb Q}}
\def\F{{\mathbb F}}
\def\inv{^{-1}}
\def\Real{\text{\rm Re}}
\def\Imaginary{\text{\rm Im}}
\def\End{\text{End}}
\def\Hom{\text{Hom}}
\def\Aut{\text{Aut}}
\def\supp{\text{supp}\,}
\def\ord{\text{ord}\,}
\def\det{\text{det}\,}
\def\Step{\text{Step}\,}
\def\ker{\text{ker}}
\def\Ker{\text{Ker}}
\def\Sym{\text{Sym}}
%\def\Res{\text{Res}}
\def\Vect0{\text{Vect}_0}
\def\Ab{\text{Ab}\,}
\def\Rings{\text{Rings}\,}
\def\Alg{\text{Gr-Alg}\,}
\def\Tr{\text{Tr}\,}
\def \oQ {{\overline\Q}}
\def \fQ {{\frak Q}}
\def \OoQ {{\cal O}_{\oQ_p}}
\def\bx {{\bf x}}
\def\by {{\bf y}}
\def\bv {{\bf v}}
\def\ulambda {{\underline{\lambda}}}
\def\sign {{\rm sign}}
\def\ua{{\underline{a}}}
\def\BP{{\mathbf{P}}}
\def\l{\lambda} 
\def\d{\delta}
\def\g{\gamma}
\def\s{\sigma}

  \newcommand{\nc}{\newcommand}
  \newcommand{\be}{\begin{eqnarray*}}
  \newcommand{\ee}{\end{eqnarray*}}
  \newcommand{\bea}{\begin{eqnarray}}
  \newcommand{\eea}{\end{eqnarray}}
  \newcommand{\ueast}{\underline{e^\ast}}
  \newcommand{\Res}{\operatorname{Res}}
  \newcommand{\res}{\operatorname{res}}
  \newcommand{\GL}{\operatorname{GL}}
  \newcommand{\SL}{\operatorname{SL}}
   \newcommand{\M}{\operatorname{M}}
 \newcommand{\ba}{\mathbf{a}}
  \newcommand{\bb}{\mathbf{b}}
\newcommand{\tfH}{\tilde{\mathfrak{H}}}
%\newcommand{\bx}{{\bf x}}
\newcommand{\bc}{{\bf c}}
%\newcommand{\fH}{{\mathfrak{H}}}
 

\newcommand{\Dist}{\operatorname{Dist}}
\newcommand{\Spec}{\operatorname{Spec}}
\definecolor{red}{RGB}{255,0,0}


\newcommand{\Ell}{\operatorname{Ell}}
%\newcommand{\Sch}{\operatorname{Sch}}
\newcommand{\Sets}{\operatorname{Sets}}
\newcommand{\Mor}{\operatorname{Mor}}
\newcommand{\Isom}{\operatorname{Isom}}
\newcommand{\pr}{\operatorname{pr}}


\topmargin -0.4in
\textwidth 7.2in
\textheight 9 in

\hoffset=-0.15 in
\voffset=-0.05in

%\small

%

\begin{document}


%\title{The Milnor $K$-theory and the Beilinson-Kato modular symbol for $\GL_n$}
%\title{The Milnor $K$-group of Siegel units and modular symbol for $\GL_n$}
\title{Modular Symbols with Values in Beilinson-Kato Distributions}

\author{Cecilia Busuioc\footnote{\texttt{celiabusuioc@ucl.ac.uk}, Department of Mathematics, 
University College London, 
25 Gordon Street, London, WC1H 0AY}, Jeehoon Park\footnote{ \texttt{jpark.math@gmail.com}, QSMS, Seoul National University, 1 Gwanak-ro, Gwanak-gu, Seoul 08826, South Korea}, Owen Patashnick\footnote{ \texttt{owen.patashnick@kcl.ac.uk}, Department of Mathematics, King's College London, London, UK; School of Mathematics, University of Bristol, Bristol, UK; The Heilbronn Institute for Mathematical Research, Bristol, UK}, and Glenn Stevens\footnote{ \texttt{ghs@math.bu.edu}, Department of Mathematics, Boston University, Boston, MA 02215}}

\date{\today}
%\author{Jeehoon Park}

\maketitle
%\abstract
\begin{abstract}
		For each integer $n\geq 1$, we construct a $\GL_n(\Q)$-invariant modular symbol $\bm\xi_n$ with coefficients in a space of distributions that takes values in the Milnor $K_n$-group of the modular function field.
	The Siegel distribution $\bm\mu$ on $\Q^2$, with values in the modular function field, serves as the building block for $\bm\xi_n$; we define $\bm\xi_n$ essentially by taking the $n$-Steinberg product of $\bm\mu$.  
	The most non-trivial part of this construction is the cocycle property of $\bm\xi_n$; we prove it by using an induction on $n$ based on the first two cases $\bm\xi_1$ and $\bm\xi_2$; the first case is trivial, and the second case essentially follows from the fact that Beilinson-Kato elements in the Milnor $K_2$-group modulo torsion satisfy the Manin relations. 
	
	
	
	%$K$-theoretic-valued distributions.
	%, and show that the definition is well-defined.
%	For each integer $n\geq 1$, we construct a $\GL_n(\Q)$-invariant modular symbol with coefficients in a space of distributions with values in the Milnor $K_n$-group of the modular function field.
	% of a modular curve with full level structure. 
%	Our $\GL_1(\Q)$-modular symbol is the Siegel distribution on $\Q^2$ with values in the modular function field. %$K$-theoretic-valued distributions.
	%, and show that the definition is well-defined.
\end{abstract}



\tableofcontents










\section{Introduction}

%The successful interplay between classical modular forms and modular symbols ({\color{red} ref.s needed}) suggests that there should be interesting applications given suitable generalizations.  The difficulty with generalizing modular forms is well known ({\color{red} ref.s needed}) but modular symbols are ripe for generalization by just replacing the classical module of coefficients $\Sym^{k-2}(\C^2)$ by any suitable module with an appropriate $\GL_2(\Q)$ action.  

The aim of this note is to 
define modular symbols with coefficients in $K$-theoretic-valued distributions.  Although the various elements that go into this definition appear in various places in the literature, to our knowledge no-one has yet packaged them in this way.  
%make just such a generalization, 
%in particular to propose by example the existence of motivic modular symbols via motivic-valued  distributions. 
In future work we hope to show  that this point of view may be useful for describing new and known Euler systems i.e. ``norm-compatible systems of cohomology classes" 
%({\color{red} ref.s for Beilinson Kato, Loeffler-Lei-Zerbes}) 
(especially that of Kato \cite{Kat}) via the distribution property of our modular symbols, and relating them to computing special values of $L$-functions and zeta functions.
%in describing special values of L-functions and zeta functions.
%More specifically, we are

The Siegel units on the upper half plane studied in \cite{Kat} by K. Kato play a key role in his construction of an Euler system for elliptic curves; see also \cite{Sch} for related Siegel-Kato functions on the universal elliptic curve.
Two key features of the Siegel units are known, namely,

(S-1) the Siegel units satisfy a  distribution property; see Kubert-Lang \cite{KubertLang} for details,

(S-2) the Siegel units satisfy the Manin relations in the Milnor $K_2$-group of a modular curve; %the connections between the Manin relations and $K_2$ were first suggested by 
see Brunault \cite{Bru} and (independently) Goncharov \cite{Gon}.

This distribution property and the Manin relations suggest the  existence of the aforementioned $\GL_2(\Q)$-invariant modular symbol. The goal of this article is to construct such a modular symbol explicitly.


We now recall the ``classical" notion of modular symbol over a subgroup $G$ of $\GL_2(\Q)$ taking values in a contravariant $G$-module $M$.   
For this we let $\Delta_0 :=\operatorname{Div}^0(\BP^1(\Q))$ which we endow with the action induced by the usual covariant action of $\GL_2(\Q)$ on  
$\BP^1(\Q)$.   An additive homomorphism 
$$
\xi : \Delta_0 \lra M
$$ 
 is said to be an $M$-valued modular symbol over $G$ if it satisfies the condition 
 $$
 \gamma^\ast \cdot \xi(\gamma D) = \xi(D)
 $$ 
 for all $D\in \Delta_0$ and all $\gamma\in G$.  We will also call this a $G$-invariant modular symbol, with the understanding that $G$ acts on $\Hom_\Z(\Delta_0, M)$ on the right by $(\xi,\gamma) \longmapsto \xi | \gamma$ where $\xi|\gamma \in \Hom_\Z(\Delta_0, M)$ is given by
$$
(\xi | \gamma)(D) = \gamma^\ast \cdot \xi ( \gamma D).
$$
Thus $\xi$ is said to be a modular symbol over $G$ if and only if it is fixed by $G$ with the above-defined right (i.e.\ contravariant) action.
Finally we will denote the group of all $M$-valued modular symbols over $G$ by $\operatorname{Symb}_G(M)$.


\noindent
%{\color{blue} {\bf Delete the following.}
%For a subgroup $G$ of $\GL_2(\Q)$ and a contravariant $G$-module $M$, a $M$-valued $G$-invariant modular symbol is defined to be a $G$-equivariant homomorphism from $\Delta_0$ to $M$, where $\Delta_0=\operatorname{Div}^0(\BP^1(\Q))$ is the group of degree 0 divisors of the cusps $\BP^1(\Q)=\Q \sqcup \{\infty\}$ on which  $G$ acts on the left. Let $\operatorname{Symb}_G(M)$ be the module of $M$-valued $G$-invariant modular symbols. We use the notation $r \to s = (s)-(r) \in \operatorname{Div}^0(\BP^1(\Q))$ for $r,s \in \BP^1(\Q)$\footnote{For a congruence subgroup $\Gamma \subset \SL_2(\Z)$ and $M=\Sym^{k-2}(\C^2)$, Ash-Stevens \cite{AS} proved that
%	\be
%	\operatorname{Symb}_{\Gamma}(\Sym^{k-2}(\C^2)) \simeq H^1_c(\Gamma, \Sym^{k-2}(\C^2))
%	\ee
%	where $\Sym^{k-2}(\C^2)$ is the space of homogeneous polynomials of degree $k-2$ with $\Gamma$-action and $H^1_c$ means the cohomology with compact support. The  space of modular forms for $\Gamma$ of weight $k$ can be embedded into $H^1_c(\Gamma, \Sym^{k-2}(\C^2))$ by Eichler-Shimura theory.}
%	and 
%	$$
%	\g ((s)-(r))= (\g s) - (\g r), \quad \g s= \frac{a s+b}{c s+d}, \quad \g =\begin{pmatrix}a&b\\c&d\end{pmatrix} \in \GL_2(\Q).
%	$$
%}

For an integer $N\geq 1$, let $\Gamma(N)$ denote the modular congruence subgroup of full level $N$ structure and $\fH/\Gamma(N)$ the associated modular curve. 
It is well known that there is a smooth affine algebraic curve $Y(N)$ over $\C$ such that the associated analytic space $Y(N)(\C)$ is isomorphic to $\fH/\Gamma(N)$; see Appendix \ref{appendix} for details. Let $k_N=\C(Y(N))$ be the field of rational functions on $Y(N)$. 
%Let $k_N$ be the field of meromorphic functions on the upper half-plane $\fH$ which are invariant under $\Gamma(N)$ and which extend to meromorphic functions on the usual compactification of $\fH/\Gamma(N).$  
For each $n \geq 1$, let $K_n^M(k_N)$ be the $n$th Milnor $K$-group of $k_N$. Let $k_\infty$ be the union of the $k_N$'s and let $\tilde K^M_2(k_\infty)$ be the associated direct limit of the groups
\be
\tilde K^M_2(k_N):=K^M_2(k_N) \otimes_\Z \Q.
\ee
We sometimes use the notation $ r \{x,y\}= \{x,y\} \otimes_\Z r$ where $r \in \Q$ and $x,y \in k_\infty^\times$.
%of the torsion free parts of the Milnor $K$-groups $K_2(k_N)$.  

%{\begin{color}{blue}
%First state the theorem we are proving 
%$\GL_n(\Q)$.  Say is=t as, for each n, there is a unique modular symbol such that...
%
%State the main case is
%Theorem \ref{modularsymbol}
%
%ther eis a unique $\GL_2(\Q)$ modualr symbol form there to there that the pathf rom infty to zero is euqal to mu.
%\end{color}}



Using this language, our main result is the construction of an explicit modular symbol
\be
{\bm\xi} \in \operatorname{Symb}_{\GL_2(\Q)} \biggl(\Dist(M_{2\times2}(\Q), \tilde K^M_2(k_\infty))\biggl)
\ee
using Siegel units.
%explicitly a $\Dist(\Q^4, \tilde K_2(k_\infty))$-valued $\GL_2(\Q)$-modular symbol ${\bm\xi}$ using the Siegel units.
Here $\Dist(M_{2\times2}(\Q), \tilde K^M_2(k_\infty))$ is the space of $\tilde K^M_2(k_\infty)$-valued distributions on the space $M_{2\times2}(\Q)$ of $2\times 2$ matrices with coefficients in $\Q$.
%and $\tilde K^M_2(k_\infty)$ is the direct limit of $\tilde K^M_2(k_N)$ which is the torsion free part of the Milnor $K^M_2$-group of the field $k_N$, where $k_N$ is the field of meromorphic functions on the open modular curve $\fH/\Gamma(N)$ with $\fH$ is the upper half plane and full level structure $\Gamma(N)$ for $N\geq 1$.


%{\begin{color}{blue}
%The feature (S-1) enables us to construct a 
%%distribution of 
%$M_{2\times2}(\Q)$-valued modular symbols ${\bm\xi}$. For any triple $r,s,t\in \PP^1(\Q)$ of distinct cusps, we need to show that the Manin relations hold
%\be
%&&{\bm\xi}(r\to s) + {\bm\xi}(s\to r) =0 \\
%&&{\bm\xi}(r\to s) + {\bm\xi}(s\to t) + {\bm\xi}(t\to r) = 0.
%\ee
%We show in this note that the first relation follows from the antisymmetry of the Steinberg symbol and the second one follows from the feature (S-2).
%\end{color}}


There is an important technical point regarding (S-2). For our construction of ${\bm\xi}$, we need the Manin relations for all $N\geq 1$.  Unfortunately \cite[Theorem 1.4]{Bru} only proves them for $N$ not divisible by $3$, and we find that Goncharov's claimed proof for all $N$ in \cite{Gon} has a gap.  (See Section \ref{GonchMethod} for more details)

%{\color{blue} Goncharov's argument is based on the computation of a ``Residue map" of certain elements in $K_3$ of universal elliptic curves, and an application of Suslin's reciprocity law.
%It seems likely that the functions that Goncharov 
%claimed exist and 
%{\color{red} claims to construct} for his proof in \cite{Gon} do not actually exist.}
% may not exist in $K_3$ of universal elliptic curves.
Our key technical contribution is to salvage his argument and  actually construct symbols of theta functions on Milnor $K_3$ of universal elliptic curves whose images under the appropriate residue map give the Manin relations by Suslin's reciprocity.
%here in the introduction we just write down our explicit elements in $K_3$ of universal elliptic curves without further explanation 

We should also point out that \cite{Bru2} contains a proof different from ours which works for $\GL_2(\Q)$ uniformly for all levels $N$.

Finally, we also provide a generalization of ${\bm\xi}$ to a $\GL_n(\Q)$-setting, using the language of Ash-Rudolph \cite{AR} for $\GL_n(\Q)$-invariant modular symbols. More precisely, for each $n \geq 1$, we explicitly construct a $\GL_n(\Q)$-invariant modular symbol ${\bm\xi}_n$ with values in $\Dist (M_{n\times2}(\Q), \tilde K^M_n(k_\infty))$, which equals ${\bm\xi}$ when $n=2$. Here $M_{n\times2}(\Q)$ is the vector space of $n\times 2$ matrices with coefficients in $\Q$.
%based on a simple induction argument.

%{\color{red} 
In the recent paper \cite{SV}, Sharifi and Venkatesh use motivic complexes (with theta functions similar to our article) to construct a $\GL_2(\Z)$-group cocycle with values in $K_2$-group of a modular curve. This construction is  closely related to our $\GL_2(\Q)$-invariant modular symbol $\bm\xi$ as discussed in Theorem \ref{modularsymbol}. Peter Xu generalized the method of \cite{SV} to a $\GL_n(\Z)$-setting: he sketched a construction of a variation of a $\GL_n(\Z)$-version of \cite{SV}, which he calls $\Theta_{\cC}^M$ in \cite[Section 4.2.2]{Xu}. 
It would be a good project to compare this generalization with our work.

%}


%{
% \color{blue} 
 Bergeron, Charollois, and Garcia \cite{BCG} interpreted a construction of Eisenstein cohomology classes of $\GL_n$ as a transgression of the Euler class of a real vector bundle with $\SL_n(\Z)$-action. Peter Xu worked out in \cite{Xu} that ``a regulator" of his $\Theta_{\cC}^M$ is essentially the same as the Eisenstein cohomology class of \cite{BCG}.
% In fact, the Eisenstein cohomology class in \cite{BCG} is expected to be ``a regulator" of our modular symbol with values in $K_n(k_{\infty})$. 
 
% the details of this in his thesis (\cite[Corollary 4.11]{Xu}) under the assumption that $\Theta_{\cC}^M$ in \cite[Section 4.2.2]{Xu} can be identified with our modular symbol; It would be interesting to prove that this assumption is correct.
% }



%{\color{blue}
Also note that  in \cite{FK}, Fukaya and Kato constructed ``zeta maps" (including $p$-adic \'etale regulators) to $K_2$ of modular curves and computed the image of the the Steinberg symbol of Siegel units under these maps,
which was a crucial part in their work on Sharifi's conjecture.

% lead to  a crucial part of Sharifi's conjecture.

%the modular symbols (which Goncharov claimed to have constructed) with values in  
%}

%{\color{blue} We need to mention the work of Sharifi-Venkatesh \cite{SV}: they used motivic complexes (instead of theta functions in our article) to construct a $\GL_2(\Z)$-group cocycle with values in $K_2$-group of a modular curve, which is closely related to our $\GL_2(\Q)$-invariant modular symbol $\bm\xi$ in Theorem \ref{modularsymbol}. It seems to be worthwhile to generalize \cite{SV} to a $\GL_n(\Z)$-setting and compare precisely \cite{SV} with our work.}


%{\begin{color}{blue} These comments need some modification???
%It is possible that what we have defined is not the most general possible modular symbol for the applications we have in mind.  : this statement overlaps with the comment at the end of section 1.1.
%Under the well known isomorphisms $K_n^M(F) \otimes \Q \simeq K_n^{(n)}(F)$ and $K_n^{(n)}(F)\cong \mathrm{CH}^n(F, n)\otimes \Q$ for $F$ a field 
Under the well known isomorphisms $K_n^M(F) \cong \mathrm{CH}^n(F, n)$ for $F$ a field (see for example \cite[p. 177]{Totaro}), we have actually  constructed modular symbols with values in the non-torsion parts of certain motivic cohomology groups. We think it would be 
%Suggests more general definitions.  
interesting to explore the putative existence of, and the putative properties of, modular symbols with values in more general motivic cohomology groups, 
 as well as relating them via  realizations to other well-known modular symbols with values in interesting $\Gamma$-modules (for $\Gamma$ a discrete subgroup of $\GL_n(\Q)$ which fixes a lattice in $\Q^n$).  
%\end{color}
%}
%For the authors, this generalization is foundational for ongoing work to describe new and known Euler systems i.e. "norm-compatible systems of cohomology classes" 
%({\color{red} ref.s for Beilinson Kato, Loeffler-Lei-Zerbes}) 
%(especially that of Kato \cite{Kat}) via the distribution property of our modular symbols, and relate this to computing formulae relating special values of $L$-functions and zeta functions.
% of a suitable motive (\`a la the Beilinson Conjectures ({\color{red} maybe the book rather than the primary articles })) to global invariants of the motive in question.  Despite the preliminary nature of this note, we hope that it will fill a gap in the literature and suggest new avenues of exploration to other readers as well.


%\begin{remark}
%{\color{red} compare with the beginning.  Merge these two discussions?}
%There is a reason why we are interested in constructing a $\GL_2(\Q)$-modular symbol ${\bm\xi}$ with values in $\Dist(\Q^4, \tilde K_2(k_\infty))$: in the future works, we like to construct certain interesting $G$-module $M$ (where $\Gamma$ is a discrete subgroup of $\GL_2(\Q)$) which admits a specialization map $\tilde K_2(k_\infty) \to A$ so that the distribution property of ${\bm\xi}$ would enable us to construct ``an Euler system" for
%$\operatorname{Symb}_\Gamma (A)$. %In other words, 

%(1) since the $K$-theory is motivic, there would be a chance to find a $\Gamma$-module homomorphism to an interesting module $A$,

%(2) since the value of ${\bm\xi}$ is the distribution, the distribution property for all test functions on $\Q^4$ would help to construct ``a norm-compatible system of cohomology classes", so called ``an Euler system".
% $\operatorname{Symb}_\Gamma (\Dist(\Q^4,A))$.
%\end{remark}








%\subsection{Proving the Manin relations}
\subsection{Statement of theorems and the Manin relations}
%For each natural number $N \geq 1$, let $Y(N)$ be the open modular curve with full level $N$ structure 
%%$\Gamma(N)$ 
%defined over $\Q$.
%%over the cyclotomic field $\Q(\zeta_N)$ where $\zeta_N$ is a primitive root of unity.
%Let $k_N$ be the field of meromorphic functions on $\fH/\Gamma(N) =Y(N)(\C)$ where $\fH$ is the upper half plane. 
%Let
%$$
%k_\infty := \bigcup_{N} k_N.
%$$

%A meromorphic modular function $f$ is said to be rational over a subfield $k$ of $\C$ if its Fourier coefficients (at $\infty$) all lie in the field $k$.   We define $k_N$ to be the field of all meromorphic modular functions rational over the cyclotomic field $\Q(\zeta_N)$.   In Chapter 6 of his book, Shimura gives an explicit finite set of generators over $\Q$ of $k_N$ (see section 6.2 and Proposition 6.9 of \cite{shimura}).  More precisely, he shows $k_1 = \Q(j)$ and for $N>1$, that $k_N$ is generated over $k_1$ by the Fricke functions $\left\{\,f_a \bigm| a\in \frac{1}{N}\Z^2/\Z^2\,\right\}$.    Furthermore, he defines the field of all modular functions over cyclotomic fields to be the field 
%$$
%k_\infty := \bigcup_{N} k_N.
%$$



%In section {\color{red} blah} we provide a proof for the convenience of the reader of the following well known yet key result for Siegel units. %although well known, does not have an explicit proof in the l is the following.
We start by defining Siegel units of level $N$ following Kato (\cite[subsection 1.9]{Kat}). Let $A(N):=\frac{1}{N}\Z^2/\Z^2$.

\begin{definition}\label{Siegelunits}\cite[Definition 1.1]{Bru}
For $\ba = (a_1 + \Z, a_2 +\Z)\in A(N)$ with $0\le a_1,a_2 < 1$, we define $g_\ba = 1$ if $\ba = {\bf 0}$ and otherwise define
$$
g_\ba(\tau) := q^{\frac{1}{2}B_2(a_1)} \prod_{n\ge 0} \left(1 - q^{n+a_1} e^{2\pi i a_2} \right) \prod_{n\ge 1} \left(1 - q^{n-a_1} e^{-2\pi i a_2} \right)
$$
where $q=e^{2 \pi i \tau}$ and $B_2(x)=x^2-x+\frac{1}{6}.$ We also use the notation $g_{a_1,a_2}=g_\ba$ later in Definition \ref{glnm}.
\end{definition}

%Then $g_\ba$ deter well-defined as an element of $\tilde K_1(k_N)\simeq \Q \otimes K_1(k_N)$.
Let $\cR_N$ be the subring of $k_N$ consisting of functions that are holomorphic on the upper half-plane and let
$$
\cR_\infty := \bigcup_N{\cR_N}.
$$
Then $g_\ba^{12N}$ is an element of $\cR_N^\times$ (cf. \cite{KubertLang} p. 182).
%Looking at the $q$-expansions of $g_\ba(\tau)$, we see that for any $(a_1,a_2), (b_1,b_2)$ $ \in \Q^2$, we have $(a_1,a_2) \equiv (b_1,b_2) \mod \Z^2$ if and only if $g_\ba/g_\bb$ is a a root of unity.  Thus $g_\ba$ is a well-defined element of $\cR_\infty^\times/\mu_\infty$ where $\mu_\infty$ is the torsion subgroup of $\cR_\infty^\times$.
Therefore we can regard $g_\ba$ as a well-defined element ``$g_\ba^{12N} \otimes_\Z \frac{1}{12N}$" of $\tilde K^M_1(k_N):= K^M_1(k_N) \otimes_\Z \Q \simeq k_N^\times \otimes_\Z \Q$.
%is well-defined as an element of $\cR_N^\times/\mathfrak{m}_N \subset \tilde K_1(k_N)$ (where $\mathfrak{m}_N$ is the torsion subgroup of $\cR_N^\times$).

Let $\cS(\Q^2)$ be the abelian group of $\Z$-valued locally constant functions on $\Q^2$ with bounded support; see \eqref{testftn} for details.
We let $\Q^\times$ act on $\cS(\Q^2)$ and $k_\infty$ through the action of $\GL_2(\Q)$ via the diagonal embedding of $\Q^\times$ in $\GL_2(\Q)$. 
Thus $\Q^\times$ acts trivially on $k_\infty$ and it acts on $\cS(\Q^2)$ as follows:
\be
(t^* f)(x)=f(tx), \quad t \in \Q^\times, x \in \Q^2.
\ee

The key feature of the Siegel units is the distribution property:
for all $0\neq t \in \Z$ and all $\ba \in A(N)$ we have
\be
g_\ba= \sum_{\substack{\bb \in A(N)\\ t \bb=\ba}} g_\bb,
% \quad t\in \Z\setminus \{0\}.
\ee
which can be reformulated as the following proposition.
\begin{proposition}[The Siegel distribution]\cite[Theorem 1.8]{Colmez}\label{Sieg}
There is a unique $\Q^\times$-invariant (i.e. $t \in \Q^\times$, $\bm\mu(t^* f)=\bm\mu(f)$) homomorphism\footnote{This $\Q^\times$-invariance plays a role in defining $\GL_2(\Q)$-invariant modular symbols: see the first sentence of the proof of Proposition \ref{wdefined}.}
$$
{\bm\mu} : \cS(\Q^2) \lra \tilde K^M_1(k_\infty) 
$$
for which 
$$
{\bm\mu}([\ba+\Z^2]) = g_{\ba}\ \hbox{\rm for all ${\bf a}\in \Q^2/\Z^2$}
$$
where $[\ba+\Z^2]$ is the characteristic function on $\bf a +\Z^2$.
% and $\mu_\infty$ is the torsion subgroup of $\cR_\infty^\times$.
%Moreover, ${\bm\mu}$ is invariant under the action of $\GL_2^+(\Q)$, i.e.\ for any $\gamma \in \GL_2^+(\Q)$ and any $\phi\in \cS(\Q^2)$ we have
%\footnote{This invariant property of ${\bm\mu}$ seems important, though we do not use it in our construction of $\GL_n(\Q)$-invariant modular symbols. For us, $\Q^\times$-invariance plays a role in defining $\GL_n(\Q)$-invariant modular symbols.}
%$$
%{\bm\mu}(\gamma^*\phi)= {\bm\mu}(\phi)\big |_{\gamma} 
%$$ 
%where $(f\big |_\gamma )(\tau) = f (\gamma \tau)$ for $f \in \cR_\infty^\times$ and $\tau \in \fH$.
\end{proposition}

%$\GL_2(\A_f)$, i.e.\ for any $\alpha \in \GL_2(\A_f)$ and any $\phi\in \cS(\Q^2)$ we have
%$$
%{\bm\mu}(\alpha\phi)|\alpha = {\bm\mu}(\phi).
%$$ 
%\end{proposition}

%In the language of distibutions we say that ${\bm\mu}\in \Dist(V,U)$ is a $\GL_2(\A_f)$-invariant $U$-valued distribution on $V$, i.e.
%${\bm\mu}|\alpha = {\bm\mu}$ for all $\alpha\in \GL_2(\A_f)$ with the obvious conventions.


%If $a,b\in \Q^2$ we let $A(a,b) \in M_{2\times2}(\Q)$ be the matrix whose top row is $a$ and whose bottom row is $b$. In the following we will identify $\cS(\Q^2)\otimes\cS(\Q^2)$ with  $\cS(M_{2\times2}(\Q))$ by identifying $\phi_1\otimes \phi_2\in \cS(\Q^2)\otimes\cS(\Q^2)$ with the test function on $M_{2\times2}(\Q)$ defined by
%\bea \label{identify}
%(\phi_1\otimes \phi_2) : A(a,b)\longmapsto \phi_1(a)\cdot \phi_2(b)\quad\hbox{\rm for all $a,b\in \Q^2$}.
%\eea
%For $\delta\in \Dist(M_{2\times2}(\Q),\tilde K_2(k_\infty))$ and $\alpha\in \GL_2(\Q)$, we define the group action $\alpha^*\delta$ in $\Dist(M_{2\times2}(\Q),\tilde K_2(k_\infty))$ by 
%\bea
%(\alpha^* \delta)(\phi) = \delta((\alpha^{-1})^{*}\phi)
%\eea
%for an arbitrary element $\phi\in \cS(M_{2\times2}(\Q))$, where $\alpha^*(\phi)(A) = \phi(\alpha A)$ 
%%and $(\alpha(\phi))(A) = \phi(A\alpha)$ 
%for any $A\in M_{2\times2}(\Q)$.

%If $a,b\in \Q^n$ we let $A(a,b) \in M_{n\times 2}(\Q)$ be the $n\times 2$ matrix whose left column is $a$ and whose right column is $b$. In the following we will identify $\cS(\Q^n)\otimes\cS(\Q^n)$ with  $\cS(M_{n\times 2}(\Q))$ by identifying $\phi_1\otimes \phi_2\in \cS(\Q^n)\otimes\cS(\Q^n)$ with the test function on $M_{n\times 2}(\Q)$ defined by
%\bea \label{identify}
%(\phi_1\otimes \phi_2) : A(a,b)\longmapsto \phi_1(a)\cdot \phi_2(b)\quad\hbox{\rm for all $a,b\in \Q^n$}.
%\eea


If $a_1, \cdots, a_n \in \Q^2$, let $A(a_1, \cdots, a_n) \in M_{n\times 2}(\Q)$ be the $n\times2$ matrix whose rows are $a_1, \cdots, a_n$.
In the following we will identify $\cS(\Q^2)\otimes\cdots \otimes \cS(\Q^2)$ with  $\cS(M_{n\times 2}(\Q))$ by identifying $\phi_1\otimes \cdots \otimes \phi_n\in \cS(\Q^2)\otimes \cdots \otimes\cS(\Q^2)$ with the test function on $M_{n\times 2}(\Q)$ defined by
\bea \label{identify}
(\phi_1\otimes \cdots \otimes \phi_n) : A(a_1,\cdots, a_n)\longmapsto \phi_1(a_1)\cdots \phi_n(a_n)\quad\hbox{\rm for all $a_i\in \Q^2$}.
\eea
For $\delta\in \Dist(M_{n\times 2}(\Q),\tilde K^M_n(k_\infty))$ and $\alpha\in \GL_n(\Q)$, we define the group action 
\be
\alpha^*\delta \in \Dist(M_{n \times 2}(\Q),\tilde K^M_n(k_\infty)), \quad 
(\alpha^* \delta)(\phi) = \delta((\alpha^{-1})^{*}\phi)
\ee
for an arbitrary element $\phi\in \cS(M_{n\times 2}(\Q))$, where $\alpha^*(\phi)(A) = \phi(\alpha A)$ 
%and $(\alpha(\phi))(A) = \phi(A\alpha)$ 
for any $A\in M_{n\times 2}(\Q)$.

\begin{definition}\cite[1.4.2]{Colmez}\label{BKdist}
We define the Beilinson-Kato distribution 
$$
{{\bm\mu}^2} \in \Dist(M_{2\times2}(\Q),\tilde K^M_2(k_\infty)), \quad
{{\bm\mu}^2}(\phi_1\otimes \phi_2) := \left\{\,{\bm\mu}(\phi_1),{\bm\mu}(\phi_2)\,\right\}.
$$
\end{definition}


\begin{theorem}\label{modularsymbol}
There is a unique $\GL_2(\Q)$-invariant modular symbol
$$
{\bm\xi} : \Delta_0 \lra \Dist(M_{2\times2}(\Q),\tilde K^M_2(k_\infty))
$$
with the property that 
$${\bm\xi}\biggl((0)-(\infty)\biggr) = {{\bm\mu}^2}. $$
%for which ${\bm\xi}\biggl((s)-(r)\biggr) = {\bm\xi}(r\to s)$ for every pair of distinct cusps $r,s\in \PP^1(\Q)$.
In other words, $\bm\xi$ is a group homomorphism satisfying
\be
\bm\xi (\g D) = (\g^{-1})^* \bm\xi(D), \quad D \in \Delta_0, \g \in \GL_2(\Q).
\ee
\end{theorem}


%Using the Beilinson-Kato distribution ${{\bm\mu}^2}$, we will construct a modular symbol
%$$
%{\bm\xi} : \Delta_0=\operatorname{Div}^0(\BP^1(\Q)) \lra \Dist(M_{2\times2}(\Q),\tilde K_2(k_\infty))
%$$
%such that ${\bm\xi}\biggl((0)-(\infty)\biggr)={{\bm\mu}^2}$.
%See Theorem \ref{modularsymbol} for details.
%For any pair of distinct rational cusps $r,s\in \PP^1(\Q)$ of the upper half-plane we chose $\alpha\in \GL_2(\Q)$ so that $\alpha(\infty) = r$ and $\alpha(0) = s$ and define ${\bm\xi}(r\to s) \in \Dist(M_{2\times2}(\Q), \tilde K_2(k_\infty))$ by 
%$$
%{\bm\xi}(r\to s) :=  \alpha{{\bm\mu}^2}.
%$$
%We will show that ${\bm\xi}$ is well-defined and, in fact, is the desired modular symbol; see Theorem \ref{modularsymbol} for details.




%Our goal is to construct, using Siegel units, an explicit $\GL_2(\Q)$-invariant modular symbol with values in a Milnor $K_2$-group-valued distribution space $\Dist(\Q^4, \tilde K_2(k_\infty))$. 
For any triple $r,s,t\in \PP^1(\Q)$ of distinct cusps, we need to show that the Manin relations hold
\be
&&{\bm\xi}(r\to s) + {\bm\xi}(s\to r) =0 \\
&&{\bm\xi}(r\to s) + {\bm\xi}(s\to t) + {\bm\xi}(t\to r) = 0.
\ee
The first relation follows immediately from the antisymmetry of the Steinberg symbol. %{\color{red} 
However, the proof of the second Manin relation is more involved.
%}
%{\color{blue} and the second Manin relation needs a serious work.}
We prove the second Manin relation \eqref{Maninrelation} for Siegel units by applying Suslin's  reciprocity law to certain elements in $K_3$ of the function field of an elliptic surface.
%a universal elliptic curve.
%a field consisting of theta functions.


We define the complex manifold $\cE_N:=(\C \times \fH)/ \tilde \Gamma(N)$ and a natural quotient map
\bea \label{univell}
\pi_N: \cE_N:=(\C \times \fH)/ \tilde \Gamma(N) \to Y_N:=\fH/\Gamma(N)
\eea
where the action of $\tilde \Gamma(N)=\Z^2 \rtimes \Gamma(N)$ on $\C \times \fH$ is the one described in the proof of Proposition \ref{Theta:transform} (see displayed equation \eqref{saction})
and $\pi_N$ is induced from the natural projection from $\C \times \fH$ to $\fH$. 
For $N\geq 3$, it is known that there is an algebraic variety $\cE(N)$ over $Y(N)$ of relative dimension $1$ such that the associated analytic space $\cE(N)(\C)$ is isomorphic to $\cE_N$ as complex manifolds; in fact, there is an elliptic surface $\overline \cE(N)$ over $X(N)$, which is the smooth compactification of $Y(N)$, in the sense of \cite[Definition on p. 202]{Silverman}; see Appendix \ref{appendix} for details.
Let $E(N)$ be the generic fiber of $\overline \cE(N)$, i.e. $E(N)=\Spec(k_N)\times_{X(N)}\overline \cE(N)$:
\[\begin{tikzcd}
	{E(N)} & {\overline \cE(N)} \\
	{\Spec(k_N)} & {X(N).}
	\arrow[" ", from=1-1, to=1-2]
	\arrow["{}"', from=1-1, to=2-1]
	\arrow[from=1-2, to=2-2]
	\arrow[" ", from=2-1, to=2-2]
\end{tikzcd}\]
Note that $E(N)$ is an elliptic curve over $k_N$ and the group of $E(N)(k_N)$ of $k_N$-rational points of $E(N)$ is exactly the group of points of order $N$ of $E(N)$ for $N\geq 3$; see \cite[Remark 5.5, Theorem 5.5]{Shioda} for more details.
\begin{definition} \label{fn}
For $N\geq 3$, let $F_N=k_N(E(N))$ be the field of rational functions on $E(N)$ over the field $k_N$. 
\end{definition}

Let $F_\infty:= \bigcup_{N} F_N$. Then it is easy to see that $F_\infty$ is a field extension of $k_\infty$ of transcendence degree 1, since $\cE(N)$ has relative dimension 1 over $Y(N)$.
We will use Theorem \eqref{SRL} (Suslin's reciprocity law for Milnor $K$-groups) applied to $$\Res: \tilde K^M_3(F_\infty) \to \tilde K^M_2(k_\infty)$$ 
%{\color{red} 
as follows:
%}
%Let $A(N) := {1\over N} \Z^2/\Z^2$ and $A(N)^\prime := A(N)\setminus\{0\}$.  

%
%Then $g_\ba$ is a holomorphic function on $\fH$ and $g_\ba^{12N}$ is known to be modular with respect to $\Gamma(N)$ and lies in ${\cO}(Y(N))^\times$; \cite[Chapter 19  \S 2]{Lang}.
%Thus $g_\ba$ is well-defined as an element of  ${\cO}(Y(N))^\times\otimes \Q$.
For each integer $N \geq 3$, we construct an element $\Phi_N \in K^M_3(F_N)$ such that
\bea
\Res(\Phi_N)=N^2(12N^2)^3\left( \{g_{\ba-\bb},g_{\bc-\ba}\} +  \{g_{\bc-\ba},g_{\bb-\bc}\} +  \{g_{\bb-\bc},g_{\ba-\bb}\} \right)
\eea
where $\ba,\bb,\bc \in A(N)=\frac{1}{N}\Z^2/\Z^2$;  
$$
\Phi_N:=\sum_{\bx \in A(N)} \left\{\frac{_{N}\!\Theta_{\ba}(u,\tau)^{12}}{_{N}\!\Theta_{\bx}(u,\tau)^{12}},\frac{_{N}\!\Theta_{\bb}(u,\tau)^{12}}{_{N}\!\Theta_{\bx}(u,\tau)^{12}},\frac{_{N}\!\Theta_{\bc}(u,\tau)^{12}}{_{N}\!\Theta_{\bx}(u,\tau)^{12}}\right\}.
%\in \tilde K_3(F_N)
$$
where $_{N}\!\Theta_{*}(u,\tau)^{12}$ are specific theta functions on universal elliptic curves which will be defined in the body of the paper;
see Lemma \ref{bcomputation} and Proposition \ref{klem} for further details.
By applying Suslin's reciprocity law to $\Res(\Phi_N)$, we get the following identity for all $N\geq 1$ (cf. Theorem \ref{sres}):
%(so called the Manin relation)
\bea\label{Maninrelation}
\{g_\ba, g_\bb \}+\{g_\bb, g_\bc\}+\{g_\bc, g_\ba\}=0\quad\text{in}\quad  \tilde K^M_2(k_N),
\eea
for $\ba,\bb,\bc \in A(N)$ satisfying $ \ba+\bb+\bc=0$.
%and gives a beautiful explicit description of the group $\Aut(k_\infty)$ of all automorphisms of $\cF$.  

For the $\GL_n(\Q)$-invariant modular symbol for $n\geq 1$, we use the language of Ash-Rudolph \cite{AS}: see Definition \ref{AR} for a precise definition.
We define the distribution ${\bm\mu}^n$ in the $\GL_n(\Q)$-setting.
%If $a_1, \cdots, a_n \in \Q^2$, let $A(a_1, \cdots, a_n) \in M_{n\times 2}(\Q)$ be the $n\times2$ matrix whose rows are $a_1, \cdots, a_n$. Then, using the identification analogous to  equation \eqref{identify}, we identify $\cS(\Q^2)\otimes \cdots \otimes \cS(\Q^2)$ with $\cS(M_{n\times 2}(\Q))$.
\begin{definition}\label{glnSD}
%For any $\ulambda = (\lambda_1,\lambda_2,\ldots,\lambda_n) \in (V^\ast)^n$, 
%We define
The distribution ${\bm\mu}^n \in \Dist(M_{n\times2}(\Q),\tilde K^M_n(k_\infty))$ is defined by
\be
{\bm\mu}^n(\phi_1 \otimes \cdots \otimes \phi_n) = \{{\bm\mu}(\phi_1), \ldots, {\bm\mu}(\phi_n) \},
\quad \phi_i \in \cS(\Q^2),
\ee
where $\phi_1 \otimes \cdots \otimes \phi_n \in \cS(\Q^2) \otimes \cdots \otimes \cS(\Q^2) \simeq \cS(M_{n\times 2}(\Q))$.
%Here we use the notation $\bold g_f = g_{a_1,a_2}=g_{\ba}$ where $f$ is the characteristic function on $\ba=(a_1,a_2)+ \Z^2$ and $g_\ba$ is the Siegel units in Definition \ref{Siegelunits}.
%we use $\cS(M_{n\times 2}(\Q)) \simeq \cS(\Q\times \Q) \otimes \cdots \otimes \cS(\Q\times \Q)$.
\end{definition}
Note that ${\bm\mu}^1$ is exactly same as the Siegel distribution ${\bm\mu}$ in Proposition \ref{Sieg} and ${\bm\mu}^2$ is the Beilinson-Kato distribution in Definition \ref{BKdist}.
%Let $V=\Q^n$ and $V^*$ be the dual vector space of $V$.
Let $\{e_1, \cdots, e_n\}$ be the standard basis of $\Q^n$ and $\{e_1^*, \cdots, e_n^*\}$ denote the dual basis of $(\Q^n)^*:=\Hom(\Q^n,\Q)$.

\begin{theorem}\label{glnmodularsymbol}
For each $n \geq 1$ there is a unique $\GL_n(\Q)$-invariant modular symbol ${\bm\xi}_n$
\bea
{\bm\xi}_n: (\operatorname{Hom}(\Q^n,\Q))^n \to \Dist(M_{n\times 2}(\Q), \tilde K^M_n(k_\infty))
\eea
such that ${\bm\xi}_n (e_1^*, \cdots, e_n^*) = {\bm\mu}^n$.
\end{theorem}
%When $n=1$, one can check that 
Note that ${\bm\xi}_1$ reduces to the Siegel distribution ${\bm\mu}$, and 
%When $n=2$, one can check that 
${\bm\xi}_2$ reduces to the Beilinson-Kato modular symbol ${\bm\xi}: \operatorname{Div}^0(\BP^1(\Q)) \to \Dist(M_{2\times2}(\Q), \tilde K^M_2(k_\infty))$: see Proposition \ref{cgcom} for details.
In Section \ref{glngeneral}, we will actually prove a basis-free version of Theorem \ref{glnmodularsymbol}, as we work with a $\Q$-vector space of dimension $n$ instead of $\Q^n$.




\vspace{1em}
We are aware that our result is not the most general possible.  We expect that our results may extend to curves over more general rings.
Note in particular that (almost) integrality is clear from our explicit computations of the residue map on Steinberg symbols of theta functions; for example, Lemma \ref{bcomputation} shows the precise denominator $\frac{1}{(12N^2)^3}$ in the computation. 
It seems possible to show that the computation generalizes to arbitrary moduli spaces of elliptic curves over a base scheme $S$ and their associated modular curves.
A potential sticking point is a precise generalization of Suslin's reciprocity result, though we feel confident that more recent advances in the field should be applicable (see for example \cite{Kriz}).\
It would certainly be a worthwhile project to pursue this abstraction.  However, we felt that a paper that explains the complex analytic case clearly and correctly would be a valuable prior contribution to the literature.
%The main reason we have not yet pursued this abstraction was for ease of exposition; feeling that a paper that explains the complex analytic case clearly and correctly would be the more valuable first contribution to the literature.
%\subsection{$\GL_2(\Q)$-invariant modular symbols}
%
%Shimura gave a beautiful explicit description of the group $Aut(k_\infty)$ of all automorphisms of $k_\infty$.  The result is as follows.
%
%\begin{theorem}
%There is a canonical exact sequence
%$$
%1 \lra \Q^\times \lra \GL_2(\A_f) \overset{s}\lra Aut(k_\infty)\lra 1
%$$
%where $s: \GL_2(\A_f) \lra Aut(k_\infty)$ is characterized by the following three properties, for all $f\in k_\infty$, and we write
%$f|\alpha := f^{s(\alpha)}$.
%\begin{itemize}
%\item[(a)] For any $\alpha \in \GL_2^+(\Q)$ we have $f|\alpha = f\circ\alpha$;
%\item[(b)] For any $\kappa \in \GL_2(\hat\Z)$, and $N\in \N$, if $f\in k_N$ and $\kappa \equiv I \pmod{N}$, then $f|\kappa = f$;
%\item[(c)] For any $d\in \hat{\Z}^\times$, if $\sigma_d := \begin{pmatrix} 1&0\\ 0&d\end{pmatrix}$, then 
%the $q$-expansion of $f|\sigma_d$ is obtained by applying ${\rm Art}(d)$ to the fourier coefficients of $f$, where ${\rm Art}(d)$ is the automorphism of the cyclotomic field $\Q^{ab}$ given by $\zeta\mapsto \zeta^d$ for every root of unity $\zeta$.
%\end{itemize}
%\end{theorem}
%

%Then both $\cR$ and its unit group $\cR^\times$ are preserved by the action of $Aut(k_\infty)$ on $k_\infty$, thus $\cR^\times$ inherits a natural action of $\GL_2(\A_f)/\Q^\times$.
%
%
%There is a natural action of the group $\GL_2(\A_f)$ on $\cS(\Q^2)$ characterized by the following two properties:
%
%\begin{itemize}
%\item[(a)] For all $\alpha\in \GL_2(\Q)$ and $\phi\in \cS(\Q^2)$ we have 
%$$
%(\alpha\phi)(x) = \phi(x\alpha)\ \hbox{\rm for all $x\in \Q^2$}.
%$$
%\item[(b)] For all $\kappa \in \GL_2(\hat{\Z})$, and $\ba = a +\Z^2\in \Q^2/\Z^2$,  if $\chi_{\ba}$ is the characteristic function of $\ba$ on $\Q^2$, then
%$$
%\kappa \chi_\ba = \chi_{\ba\kappa^{-1}}.
%$$
%\end{itemize}
%Indeed, note that if we let 
%$
%\bigotimes_p{} \cS(\Q_p^2)  
%$
%denote the {\it restricted} tensor product with respect to the distinguished elements $\chi_{\Z_p^2} \in \cS(\Q_p)$ for each prime $p$, then the canonical map
%$$
%\bigotimes_p \cS(\Q_p^2) \lra \cS(\Q^2)  
%$$
%is an isomorphism.   Clearly $\GL_2(\A_f)$ acts on the restricted tensor product.   By transporting this action via the above isomorphism we obtain an action on $\cS(\Q^2)$ that is easily seen to satisfy properties (a) and (b) above.%\subsection{Structure of the paper}

%{\color{red} moved the summary here}


\vspace{1em}
The structure of the rest of the paper is as follows: 
%{\color{red} 
Following a brief digression (Section \ref{GonchMethod}) on the gap in the proof in \cite{Gon}, we begin in Section \ref{background} by recalling the concepts of distributions, Milnor $K$-theory, and Suslin reciprocity.  In Section \ref{BeilinsonKato}, we continue by recalling theta functions and how to construct elliptic functions out of theta functions.  We then prove that the Manin relations are satisfied and complete the construction of the $\GL_2(\Q)$-invariant modular symbol.  In Section \ref{glngeneral}, we generalize our construction of modular symbols to those with coefficients in $\tilde K^M_n(k_\infty)$.  
%after recalling the facts about Siegel units that we will need, 









\subsection{Remark on Goncharov's method}\label{GonchMethod}
%{\color{red} 
In this section we briefly explain the argument of Goncharov in \cite{Gon}, we point out where we believe a gap exists in his argument, and we explain how we clarify it.

Let $E$ be an elliptic curve over an algebraically closed field $k$.
Let $k(E)$ be a function field of $E/k$.
For any field $F$, Goncharov considers the Bloch complex (\cite[section 2.1]{Gon})
\be
\cB_2(F) \xrightarrow{\delta_2} \wedge^2 F^\times
\ee
and the commutative diagram
\[\begin{tikzcd}
	{\cB_2(k(E))\otimes k(E)^\times} & {\wedge^3 k(E)^\times} \\
	{\cB_2(k)} & {\wedge^2k^\times.}
	\arrow["{\delta_2\wedge \mathrm{Id}}", from=1-1, to=1-2]
	\arrow["{\mathrm{Res}}", from=1-1, to=2-1]
	\arrow["{\mathrm{Res}}", from=1-2, to=2-2]
	\arrow["{\delta_2}", from=2-1, to=2-2]
\end{tikzcd}\]

Goncharov \cite[Theorem 2.4]{Gon} proved (in a previous paper) that there exists a map $h$ such that the following diagram commutes:
\[\begin{tikzcd}
	{\cB_2(k(E))\otimes k(E)^\times} & {\wedge^3k(E)^\times} \\
	{\cB_2(k)} & {\wedge^2k^\times.}
	\arrow["{\delta_2\wedge \mathrm{Id}}", from=1-1, to=1-2]
	\arrow["{\mathrm{Res}}", from=1-1, to=2-1]
	\arrow["h", from=1-2, to=2-1]
	\arrow["{\mathrm{Res}}", from=1-2, to=2-2]
	\arrow["{\delta_2}", from=2-1, to=2-2]
\end{tikzcd}\]
For each $N$-torsion point $a$ in $E(k)$ Goncharov defines an element (\cite[Corollary 2.6]{Gon})
\be
\theta_E(a) \in k^\times \otimes \Q.
\ee
Then, for any $N$-torsion points $a,b$ on $E(k)$, he claims that the function
\be
t \mapsto \frac{\theta_E(t-a)^N}{\theta_E(t-b)^N}
\ee
is an element of $k(E)^\times$ satisfying $\mathrm{Div} (\frac{\theta_E(t-a)^N}{\theta_E(t-b)^N}) = N (\{a\} - \{b\})$.
(see \cite[section 2.4, page 1889-1890]{Gon}).  
%{\color{red} 
Note that Goncharov is claiming something quite strong here, namely that he has constructed a function on $E$ with a very specific form.  Although it is true that the divisor $N (\{a\} - \{b\})$ on an elliptic curve over a field is principal, there is no sense in which such a function is canonical or has a canonical form (see for example the opening sentence of \cite[section 1.2]{Sch}).  Furthermore, it is no longer true in general that $N (\{a\} - \{b\})$ is principal on an elliptic curve over a base scheme $S$ different from a field. (For an explicit example see the last remark in \cite[section 1.2]{Sch}.)
%}

Goncharov then computes the image of the wedge product of theta functions $\frac{\theta_E(t-a)^N}{\theta_E(t-b)^N}$ in $k(E)^\times$ under the map $h$ (in \cite[Theorem 2.4]{Gon}) in order to prove the Manin relations for Siegel units.  %{\color{red} 
So for Goncharov, the form of his alleged function is crucial to his proof.
The problem is that he does not actually define $\theta_E(t)$ for $t \in E(k)$ (Note that $\theta_E(t)$ is defined only for each $N$-torsion point $t$) in \cite{Gon} and thus the notation $\frac{\theta_E(t-a)^N}{\theta_E(t-b)^N}$ does not make sense as an element of $k(E)^\times$.
%{\color{blue} In order to guess what Goncharov had in mind when he was writing the notation $\frac{\theta_E(t-a)^N}{\theta_E(t-b)^N}$, }
%{\color{red} 
	
There is a hint as to what Goncharov may have had in mind here regarding the construction of $\theta_E(t)$ in a reference \cite{GL} to a previous paper with A. Levin.  However, this does not seem to help. In \cite[Remark, p 400]{GL} the authors specify the need to view $t\mapsto \theta_E(t)$ as valued in the Bloch group $B_2(E)$, and it is clear from its properties that the group $B_2(E)$ does not contain $k(E)^\times$ as a subgroup.
%}
%{\color{blue} we searched his previouspaper with Levin regarding the construction of $\theta_E$ in \cite{GL} but this does not help either: Goncharov-Levin claims in \cite[Remark, p 400]{GL} that one needs to view $t\mapsto \theta_E(t)$ valued in the Bloch group $B_2(E)$, which is still not an element of $k(E)^\times$.} 

%So there is a gap in the argument of Goncharov.


%what a possible definition for $\theta_E(t), t\in E(k)$ is available but we were not able to succeed: if one looks at the original construction of $\theta_E$ of Levine-Goncharov in \cite[Remark, p 400]{GL}, one needs to view $t\mapsto \theta_E(t)$ valued in the Bloch group $B_2(E)$, which is still not an element of $k(E)$.
Our clarification has two important features that differentiate it from \cite{Gon}.  
The first is that we do not use the map $h$ (in \cite[Theorem 2.4]{Gon}), but rather a related result that follows from the Suslin reciprocity law for $\wedge^3 k(E)^\times \xrightarrow{\mathrm{Res}} \wedge^2 k^\times$, and the second is that we  explicitly (Proposition \ref{dcomputation}) find a function $f_{a,b}$ which satisfies $\mathrm{Div}(f_{a,b}) = 12N^2 (\{a\}-\{b\})$ using the Kato-Siegel functions (see for example \cite[Theorem 1.2.]{Sch} for details about how such functions are constructed over a more general base). 
%{\color{blue}- we refer to Proposition \ref{dcomputation}.}



%{\color{blue}The following remark of T. Scholl was our guiding venue to find $f_{a,b}$: if one looks at \cite[section 1.2]{Sch}, there is a sentence: "If $\cD$ is a principal divisor of an elliptic curve over a field, there is in general no `canonical' function with divisor $\cD$". So finding $f_{a,b}$ satisfying $\mathrm{Div}(f_{a,b}) = 12N^2 (\{a\}-\{b\})$ in a `canonical form' is not a trivial matter. After some efforts, we ended up to consider}

%{\color{blue} In particular, we define  $f_{a,b}=(_{N}\!\Theta_{\ba}/_{N}\!\Theta_{\bb})^{12}$, a `canonical' meromorphic function on $\cE_N$ which has divisor $\mathrm{Div}(f_{a,b}) = 12N^2 (\{a\}-\{b\})$.} 

%{\color{red} 
Note that we do not see how to recover the function Goncharov wants from our construction; if one defines $\Theta_{\ba}$ in a similar way without $N$,
	%(see displayed equation \eqref{cdiv} for details)
	 then $\Theta_{\ba}^{12}/\Theta_{\bb}^{12}$ would not be invariant under $\tilde \Gamma(N)$ (see the footnote after displayed equation \eqref{univell})  and so is not a meromorphic function.  In particular, $(\Theta_{\ba}^{12}/\Theta_{\bb}^{12})^N$ does not define a function with divisor $N (\{a\} - \{b\})$.%}

%{\color{blue} We suspect that  $\Theta_{\ba}^{12}/\Theta_{\bb}^{12}$ `somehow corresponds' to Goncharov's wrong $\frac{\theta_E(t-a)^{12N^2}}{\theta_E(t-b)^{12N^2}}$.}

To be clear, our proof in Section \ref{ManinRelsProof} is similar in outline to Goncharov's plan in \cite{Gon} for the proof of the Manin relations.  In particular, we use (a corrected version of) Goncharov's beautiful trick (Lemma \ref{gonchvariantlemma}; see also \cite{Gon}, Lemma 2.11) of summing over all $N$-torsion points to get zero.


%\vspace{1em}
%
%\subsection{On Generalizations}
%We are aware that our result is not the most general possible.  We fully expect that our results should extend to curves over more general rings, in particular, allowing for a $\GL_n(\Z)$ action (rather than $\GL_n(\Q)$), distributions over $\Z^n$, and modular curves over integral or almost integral rings.  
%Note in particular that (almost) integrality is clear from our explicit computations of the residue map of Steinberg symbols of theta functions; for example, Lemma \ref{bcomputation} shows the precise denominator $\frac{1}{(12N^2)^3}$ in the computation. It is fairly straightforward to show that the computation generalizes to arbitrary moduli spaces of elliptic curves over a base scheme $S$ and their associated modular curves.
%
%A potential sticking point is a precise generalization of Suslin's reciprocity result, though we feel confident that more recent advances in the field should be applicable (see for example \cite{Kriz}).
%
%The main reason we have not yet pursued this abstraction was for ease of exposition; feeling that a paper that explains the base case clearly and correctly would be the more valuable first contribution to the literature.







\subsection{Acknowledgements}
All four authors warmly thank three anonymous referees for their close reading of our paper and for their  substantive, perceptive, and constructive suggestions for improving it. 
 
The work of 
%Jeehoon Park 
the second author was supported by the National Research Foundation of Korea (NRF-2021R1A2C1006696) and the National Research Foundation of Korea (NRF) grant funded by the Korea government (MSIT) (No.2020R1A5A1016126).  The work of the third author was supported by the Heilbronn Institute for Mathematical Research, Bristol, UK.

\section{Background and notation}\label{background}

\subsection{Test functions and distributions}
Throughout this paper, $V$ will be a finite dimensional $\Q$-vector space of dimension $n$ and $V^\ast := \Hom(V,\Q)$ will denote the dual space.
We let $G := \Aut_\Q(V)$ be the group of $\Q$-linear automorphisms of $V$.
% and $\tilde G$ will denote the group of affine linear transformations of $V$.  We identify $\tilde G = V \rtimes G$, the semidirect product with respect to the action of $G$ on $V$.  A typical element $(x,\gamma) \in \tilde G$ acts on $V$ by
%$$
%(x,\gamma)\cdot v = x + \gamma v.
%$$

By a lattice in $V$ we mean a finitely generated additive subgroup of $V$ that spans $V$ as a $\Q$-vector space.  A subset
$X\subseteq V$ is said to be {\it uniform} if there is a lattice $L$ in $V$ for which $X+\ell = X$ for every $\ell \in L$, in which case we say
$X$ is uniform with respect to $L$.  
We say $X$ is {\it bounded} if $X$ is contained in a lattice $M$ in $V$, in which case we say $X$ is bounded by $M$.   The collection of bounded uniform subsets of $V$ generates a topology on $V$ which we call the {\it lattice topology} on $V$.   %The collection of all bounded uniform subsets of $V$ is clearly preserved by the action of $\tilde G$.

For a function $f$ on 
$V$ with bounded support and a vector $\ell\in V$ we say $\ell$ is a period of $f$ if $f(x+\ell) = f(x)$ for all $x\in V$.   The set of all periods of $f$ will be denoted 
$L_f$.   We say $f$ is locally constant if $L_f$ contains a lattice in $V$.  If $f\ne 0$, then both $L_f$ and the subgroup $M_f$ of $V$ generated 
by the support of $f$ are lattices in $V$.  We then call $L_f$ the period lattice of $f$ and $M_f$ the support lattice of $f$ and note that $L_f \subseteq M_f$.
We define
%If $R$ is an arbitrary commutative ring, we define the module of $R$-valued {\it test functions} on $V$ to be the $R$-module
\bea \label{testftn}
\cS(V) := \{\, f : V \lra \Z\,|\, f \hbox{ has bounded support and is locally constant}\}.
\eea
%We note that $\cS(V,R)$ is a ring under pointwise multiplication of functions.
%If $R = \Z$ we write simply $\cS(V):= \cS(V,\Z)$.
We note that $G$ acts contravariantly on $\cS(V)$ by pull-back:  $(\g^\ast f) (x) = f( \g x)$.
For an arbitrary abelian group $A$, we define the group of $A$-valued distributions on
$V$ to be the group
$$
\Dist(V,A) := \Hom(\cS(V),A).
$$
For $\mu \in \Dist(V,A)$ and $f\in \cS(V)$ %{\color{blue} we will also use any of the the suggestive notations} %{\color{red} 
we sometimes write $\mu(f)$ in one or more of the following (suggestive) forms:
$$
\mu(f) =: \int f\cdot d\mu = \int f(x)d\mu(x) = \int_V f(x)d\mu(x).
$$
 In the applications $A$ will always be endowed with a contravariant action by a subgroup $H \subseteq G$.   In that case,  $\Dist(V,A)$ inherits a natural contravariant action by $H$, satisfying
\bea \label{distributionaction}
\int \gamma^\ast f \cdot d\gamma^\ast\mu = \gamma^\ast \left(\int f(x)d\mu(x)\right),
\eea
for any $\mu\in \Dist(V,A)$, $f\in \cS(V), x \in V$ and $ \gamma\in H$.

%
%{ \color{blue}  {\bf Please delete this paragraph -- it is redundant.  Too much notation. }
%In our presentation, the distribution $\Dist(V\times V, A)$ on $V\times V$ will play a key role. The group $G$ acts on $V\times V$ component-wise on the left.
%Accordingly $G$ acts contravariantly on $\cS(V\times V)$ by pull-back.
%We note that if $A$ is endowed with a contravariant action of a subgroup $H \subseteq G$, then $\Dist(V\times V,A)$ inherits a canonical contravariant action of $H$
%satisfying
%\bea \label{distributionaction}
%\int \gamma^\ast f \cdot d\gamma^\ast\mu = \gamma^\ast \left(\int f(x)d\mu(x)\right),
%\quad \g^\ast f (x) = f( \g x)
%\eea
%for any $\mu\in \Dist(V\times V,A)$, $f\in \cS(V\times V), x \in V\times V$ and $ \gamma\in H$.
%}





\subsection{Milnor $K$-groups and Suslin reciprocity}

We briefly recall the Milnor ring of a field.
Let $F$ be a field. The Milnor $K$-groups are defined as follows:
\begin{eqnarray*}
K^M_0(F) &=& \Z, \quad K^M_1(F) = F^\times \\
K^M_2(F) &=& F^\times \otimes F^\times / \langle a\otimes b :  a+b= 1 \rangle \\
K^M_n(F) &=& F^\times \otimes \cdots \otimes F^\times / \langle a_1 \otimes \cdots \otimes a_n
 : \text{ there exists } i \text { such that } a_i+a_{i+1}= 1 \rangle \\
 K^M_\ast(F) &=& \bigoplus_{n\geq 0} K^M_n(F).
\end{eqnarray*} 
Then $K^M_\ast(F)$ is a graded $\Z$-algebra which is a quotient of the tensor algebra on copies of $F^\times$ modulo the homogenous ideal in degree 2 generated by the Steinberg relation ($a_i+a_{i+1}=1$)\footnote{we will hereafter abuse notation and refer to this ideal as ``the Steinberg relations''}.
%is a graded $\Z$-algebra generated by $K^M_1(F)=F^\times$ modulo the Steinberg relation ($a_i+a_{i+1}=1$).
  We use the notation $\{x_1, \cdots, x_n\} \in K^M_n(F)$ to denote the element $x_1\otimes \cdots \otimes x_n$ modulo the Steinberg relations.  
The abelian group structure of $K^M_n(F)$ is determined by the following multi-linearity property:
$$
\{x_1, x_2, \ldots, x_n\} +\{x_1', x_2, \ldots, x_n  \}=\{x_1 \cdot x_1', x_2, \ldots, x_n\}.
$$
where similar multi-linearity holds for the other components. 

%Thus, $K^M_n(A)$ is an abelian group and 
%{\color{blue}
%Then $K^M_*(F)$ has a ring structure, whose product law is induced from concatenation:
%$$
%K^M_n(F) \otimes K^M_m(F) \to K^M_{n+m}(F), \quad \{x_1, \ldots, x_n\} \cdot \{y_1, \ldots, y_m\} := \{x_1, \ldots, x_n, y_1, \ldots, y_m\}.
%$$
%}
A straightforward computation (\cite[Lemma 1.1]{Sus}) shows that 
\begin{eqnarray*}
&&\{x,y\} = -\{y, x\}, \quad x,y \in A^\times, \\
&&\xi \cdot\eta = (-1)^{nm} \eta \cdot \xi, \quad \xi \in K^M_n(A), \eta \in K^M_m(A)\\
&&\{x_1, \ldots, x_n\}=0, \quad \text{if} \quad x_1+\cdots+x_n=0 \text{ or } 1. 
%\\
%&& \{x,x\}= \{x, -1\} =\{-1, x\}, \quad  \{-x^{-1}, x\} =0,  \quad x \in A^\times.
\end{eqnarray*}

We use the notation
\be
\tilde K^M_n(F):=K^M_n(F) \otimes_\Z \Q.
\ee
%be the quotient group of $K_n(A)$ by its torsion subgroup.

%Let $F$ be a field. 
Let $v$ be a discrete valuation on $F$ and $\cO_v=\{x \in F : v(x) \geq 0\}$ be the corresponding valuation ring whose maximal ideal is denoted by $\frak m_v=\{x \in F : v(x) > 0\}$.
Let $k(v)=\cO_v/\frak m_v$ be the residue field.
For an element $a \in \cO_v$ we denote by $\bar a$ the canonical image of $a$ in $k(v)$.

\begin{lemma}\cite[Lemma 1.3]{Sus} \label{boundary}
There is an additive homomorphism $\partial_v$ from the additive group of the ring $K^M_\ast(F)$ to the additive group of $K^M_{\ast-1}(k(v))$ such that
\begin{enumerate}[(1)]
\item $\partial_v (\{u_1, \ldots, u_{n-1}, a\} = v(a) \cdot \{ \bar u_1, \ldots, \bar u_{n-1}\}$ where $u_1, \ldots, u_{n-1} \in \cO^\times$ and $a \in F^\times$.
\item $\partial_v$ is an epimorphism of graded groups of degree -1 whose kernel contains $\{1 + \frak m_v\}\cdot K^M_\ast (F)$.
\item the following diagrams are commutative:
\[\begin{tikzcd}
	{K^M_1(F)} & {K^M_0(k(v))} && {K^M_2(F)} & {K^M_1(k(v))} \\
	{F^\times} & \Z && {(\wedge^2 F^\times)/\sim} & {k(v)^\times}
	\arrow["{\partial_v}", from=1-1, to=1-2]
	\arrow["{\partial_v}", from=1-4, to=1-5]
	\arrow["v", from=2-1, to=2-2]
	\arrow["{( \ ,\ )_v}", from=2-4, to=2-5]
	\arrow["\simeq", from=2-1, to=1-1]
	\arrow["\simeq", from=1-2, to=2-2]
	\arrow["\simeq", from=2-4, to=1-4]
	\arrow["\simeq", from=1-5, to=2-5]
\end{tikzcd}\]
where $\sim$ denotes the Steinberg relations and $(a,b)_v =\bar c$, where $c=(-1)^{v(a)v(b)}\frac{a^{v(b)}}{b^{v(a)}}$.  
\end{enumerate}
\end{lemma}
(Note that we choose to define the tame symbol with a sign matching class field theory rather than the usual definition in classical descriptions of algebraic $K$-theory.  This choice of course has no bearing on the validity of subsequent constructions in this paper.)

Suppose that $F$ is a finite extension of a field $k$.
%, and let $w$ be an extension of $v$ to $E$. We denote by $e=e(w/v)$ the ramification index of $w$ with respect to $F$ and by $j_{w/v}$ the canonical inclusion of the residue field $k(v)$ in $k(w)$.
%By \cite[Lemma 1.4]{Sus}, we have the commutative diagram:
%\[\begin{tikzcd}
%	{K^M_\ast(F)} && {K^M_\ast(E)} \\
%	{K^M_*(k(v))} && {K^M_\ast(k(w))}
%	\arrow["{\partial_v}"', from=1-1, to=2-1]
%	\arrow["{\partial_w}", from=1-3, to=2-3]
%	\arrow[from=1-1, to=1-3]
%	\arrow["{e\cdot (j_{w/v})_\ast}", from=2-1, to=2-3]
%\end{tikzcd}\]
Then there is a well-defined norm map (see \cite[Lemma 1.8]{Sus} for details)
\be
N_{F/k}:K^M_n(F) \to K^M_n(k), \quad n=0,1,2.
\ee
In general, there is a well-defined norm map
\be
N_{F/k}: \tilde K^M_n(F) \to \tilde K^M_n(k), \quad n\geq 0.
\ee
%where $\tilde K^M_\ast(E)$ is the quotient group of $K^M_\ast(E)$ by its torsion subgroup.

\begin{theorem}[Suslin's reciprocity law]\cite[Theorem 2.6]{Sus}\label{SRL}
If $F$ is a field extension of a field $k$ with transcendence degree 1, then for each $x\in \tilde K^M_\ast(F)$
\be
\Res(x):=  \sum_{w \in \Sigma(F/k)} N_{k(w)/k}(\partial_w(x))=0 \ \ \text{(this is a finite sum),}\quad  \Res:\tilde K^M_n(F) \to \tilde K^M_{n-1}(k),
\ee
where $\Sigma(F/k)$ is the set of all discrete valuations (including archimedean valuations) of $F$ which are trivial on $k$. This holds for $x \in K^M_n(F)$ for $n=1,2,3$, instead of $\tilde K^M_n(F)$.
\end{theorem}


%%\begin{defn}\label{Resdef}
%%Let $E$ be a field extension of a field $k$ with the transcendence degree 1.
%%	Define the map $\Res:K_3(E) \rightarrow K_2(k)$ to be 
%%	$$\Res(x):=\sum_{w\in \Sigma(E)}N_{k(w)/k}(\partial_w(x)).$$
%%	% \quad\quad \res_x:K_3(E) \to K_2(k(v)),$$
%%%	where $\res_x$ is the group homomorphism uniquely defined by 
%%%	\be
%%%	\res_x(\pi^m_x\wedge u_1\wedge u_2) &=& m\cdot (\overline{u_1}\wedge\overline{u_2}),  \quad u_i\in \Theta(\overline \Q)^\times_x,\\
%%%	\res_x(u_0 \wedge u_1\wedge u_2) &=& 0, \quad u_i\in \Theta(\overline \Q)^\times_x
%%%	\ee
%%	%the map induced from the connecting homomorphism in the localization sequence 
%%\end{defn} 
%\begin{definition}\label{Resdef}
%Let $E$ be a field extension of a field $k$ with the transcendence degree 1.
%	Define the map $\Res:K_3(E) \rightarrow \bigoplus_{w \in \Sigma(E)} K_2(k(w))$ to be 
%	$$\Res(x):=\sum_{w\in \Sigma(E)}\partial_w(x).$$
%	% \quad\quad \res_x:K_3(E) \to K_2(k(v)),$$
%%	where $\res_x$ is the group homomorphism uniquely defined by 
%%	\be
%%	\res_x(\pi^m_x\wedge u_1\wedge u_2) &=& m\cdot (\overline{u_1}\wedge\overline{u_2}),  \quad u_i\in \Theta(\overline \Q)^\times_x,\\
%%	\res_x(u_0 \wedge u_1\wedge u_2) &=& 0, \quad u_i\in \Theta(\overline \Q)^\times_x
%%	\ee
%	%the map induced from the connecting homomorphism in the localization sequence 
%\end{definition} 
%%Note that $\res$ does not depend on the choice of uniformizer $\pi$.
% If $k$ is algebraically closed, then $k(w)=k$ and so the Weil reciprocity law implies the following corollary.
%\begin{corollary}\label{ac}
%Let $E$ be a field extension of a field $k$ with the transcendence degree 1. If $k$ is algebraically closed, then the image of $\Res$ lands in $K_2(k)$ and
%\be
%\Res(x)=0, \quad x \in K_3(E)
%\ee
%\end{corollary}







\section{The Beilinson-Kato modular symbol}\label{BeilinsonKato}

\subsection{Siegel units and theta functions}

%\begin{lemma}\label{Theta:siegel} For $\ba \in A(N)^\prime$ and $(a_1,a_2) \in \frac{1}{N}\Z^2$ the least non-negative representative of $\ba$, we have
%\begin{itemize}
%\item[\rm (a)]  $\displaystyle -e^{\pi i a_2} q^{\frac{a_1^2}{2}}\cdot \Theta\left(a_1\tau + a_2,\tau \right) = g_\ba(\tau);$ and
%\item[\rm (b)]  $\displaystyle \frac{\Theta(Nu + Na_1\tau+Na_2, \tau)}{ \Theta(Nu,\tau)} = (-1)^{Na_2} (-t^{N})^{-Na_1}q^{-\frac{N^2a_1^2}{2}}$, where $t=e^{2\pi iu}$ and $q=e^{2\pi i \tau}$.
%\end{itemize}
%\end{lemma}




We begin with the basic theta function $\Theta$ defined on the space $\tfH := \C\times \fH$ by
$$
\Theta(u,\tau) := q^{\frac{1}{12}}(t^{1/2} - t^{-1/2}) \prod_{n=1}^\infty \left(1 - q^nt\right) \left(1 - q^nt^{-1}\right)
$$
with $t = e^{2\pi i u}$ and $q=e^{2\pi i \tau}$.  The product converges uniformly on compact subsets of $\tfH$, hence defines a holomorphic function on that space.  One checks easily that, as a function of $u\in \C$, $\Theta(u,\tau)$ has simple zero at each $\omega\in \Lambda_\tau$ and no zeroes outside of $\Lambda_\tau$.  In particular, for fixed $\tau$, the divisor of $\Theta(u,\tau)$ as function of $u$ is the characteristic function $[\Lambda_\tau]$ of the lattice $\Lambda_\tau=\Z\tau+\Z$ and is therefore invariant under translation by elements of $\Lambda_\tau$.   

\begin{proposition}\label{Theta:transform}
	For all $\omega = r\tau + s \in \Lambda_\tau$ we have
	$$
	\frac{\Theta(u+\omega,\tau)}{\Theta(u,\tau)} = \psi(\omega)e^{-2\pi i r (u + \omega/2)} = (-1)^s(-t)^{-r}q^{-r^2/2}
	$$
	where 
	$$
	\psi(\omega) = \begin{cases}
		1&\hbox{\rm if $\omega\in 2\Lambda_\tau$;}\\
		\\
		-1&\hbox{\rm otherwise.}\\
	\end{cases}.
	$$ 
\end{proposition}

\begin{proof} 
	This is well-known \cite{Lang}.
\end{proof}

Let $\R^2$ denote the space of two-dimensional real row vectors and $\tilde G := \R^2 \rtimes \GL^+_2(\R)$ be the semidirect product group in which the operation $\ast$ is given by
$$
(\bx_1,\gamma_1)\ast (\bx_2,\gamma_2) = (\bx_1 + \bx_2\gamma_1^{-1}, \gamma_1\gamma_2)
$$
for any $\bx_i\in \R^2$ and $\gamma_i \in \GL^+_2(\R)$.   Setting $\tfH := \C \times \fH$, one easily checks\footnote{One needs the cocycle condition $j(\g,\tau)\cdot j(\g', \g\tau)=j(\g'\g, \tau)$ where $j(\g,\tau)=c\tau +d$ and $\g=\begin{pmatrix}a&b\\c&d\end{pmatrix}$.} that the map
\bea \label{saction}
\begin{array}{ccccc}
	\tilde G &\times& \tfH &\lra&\tfH \\
	(\bx,\gamma)&,& (u,\tau)&\longmapsto& \biggl(\displaystyle\frac{u+\rho_{\bx\gamma}(\tau)}{c\tau + d}, \gamma\tau\biggr)\\
\end{array}
\eea
defines a left action of $\tilde G$ on $\tfH$ for $\bx = (x_1,x_2)\in \R^2$, where $\rho_\bx(\tau)$ is defined by
$$
\rho_\bx(\tau) := x_1\tau + x_2.
$$ 
The group $\tilde G$ acts on the right by pullback of functions on $\tfH$.
For any function $f(u,\tau)$ on $\tfH$ and any $\tilde\gamma := (\bx,\gamma)\in \tilde G$ we have
$$
\tilde\gamma^\ast f(u,\tau) = f\left(\frac{u + \rho_{\bx\gamma}(\tau)}{c\tau +d},\gamma\tau\right)
$$
Finally, let $\Gamma := \SL_2(\Z)$ and define $\tilde\Gamma := \Z^2 \rtimes \Gamma \subseteq \tilde G$.   Then $\tilde\Gamma$
is easily seen to be a subgroup of $\tilde G$.   

%For any element $\ba\in \Q^2$ we define $\Theta_\ba : \tfH \lra \C$ by 
%$$
%\Theta_\ba(u,\tau) := \Theta(u-\rho_\ba(\tau),\tau).
%$$
%Clearly, the divisor of $\Theta_\ba(u,\tau)$ is $[\rho_\ba(\tau) + \Lambda_\tau]$.   Moreover, for $\gamma\in \Gamma$ we have 
%$$
%\gamma^\ast(\Theta_\ba)(u,\tau) = \Theta_{\ba\gamma}\left(\frac{u}{c\tau+d},\gamma\tau\right)
%$$
%
%For a pair of elements $\ba,\bb\in \Q^2$ we define
%$$
%_N\Theta_{\ba,\bb} := \frac{\Theta_\ba}{\Theta_\bb}.
%$$

\begin{definition} \label{mero}
	We define 
	$_{N}\!\Theta :\tilde \fH=\C \times \fH \lra \C$ by
	\bea \label{thetad}
	_{N}\!\Theta(u,\tau) :=  \Theta(u,\tau)^{N^2}/\Theta(Nu,\tau).
	\eea
Note that $_{N}\!\Theta$ is a meromorphic function on $\tilde \fH$, since the denominator $\Theta(Nu,\tau)$ does not introduce essential singularities; $\Theta$ is a holomorphic function on $\tilde \fH$ and $\Theta(u,\tau)$ has simple zero at each $\omega\in \Lambda_\tau$ and no zeroes outside of $\Lambda_\tau$. 

\end{definition}
%\begin{theorem}  The function $_{N}\!\Theta^{12}$.  Moreover, for $\ba,\bb\in \frac{1}{N}\Z^2$ the functions $_{N}\!\Theta_{\ba,\bb}^{12N^2}$ are elliptic and
%the finite set 
%$$
%\left\{_{N}\!\Theta_{\ba,\bb}^{12N^2} \,\bigm|\, \ba,\bb \in A(N) \right\}
%$$
%is preserved by the action of $\Gamma$.  More precisely, for $\gamma\in \Gamma$ we have 
%$$
%\gamma^\ast(_{N}\!\Theta_{\ba,\bb}^{12N^2}) = _{N}\!\Theta_{\ba\gamma,\bb\gamma}^{12N^2}.
%$$
%\end{theorem}

\noindent
The main result of this section is the following theorem.

\begin{theorem}\label{vartheta:invariance}
	There is a multiplicative character $\epsilon : \tilde\Gamma \lra {\mathbb \mu}_{12}\subseteq \C^\times$ such that the function $_{N}\!\Theta$ satisfies the relation
	$$
	\tilde\gamma^\ast (_{N}\!\Theta) =\epsilon(\tilde\gamma)\cdot_{N}\!\Theta
	$$
	for all $\tilde\gamma \in \tilde\Gamma$.  More precisely, for $\bx\in \Z^2$,  $T=\begin{pmatrix}1&1\\0&1\end{pmatrix}$, and $S=\begin{pmatrix}0&1\\-1&0\end{pmatrix}$
	we have 
	$$
	\epsilon(\bx, 1) = 1,\ \  \epsilon(0, T) = e^{2\pi i(N^2-1)/12}\ \ \hbox{\rm and}\ \ \epsilon(0,S) = i^{N^2-1}.
	$$
\end{theorem}


\begin{proof}
	Since $(\bx, 1)$, $(0,T)$ and $(0,S)$ generate $\tilde\Gamma$ we need only prove the following four identities:
	\begin{itemize}
		\item[\rm (1)]  $_{N}\!\Theta\left(u + 1,\tau\right)= _{N}\!\Theta(u,\tau)$;
		\item[\rm (2)]  $_{N}\!\Theta\left(u + \tau ,\tau\right)= _{N}\!\Theta(u,\tau)$;
		\item[\rm (3)]  $_{N}\!\Theta\left(u,\tau+1\right)=e^{2\pi i (N^2-1)/12}\cdot_{N}\!\Theta(u,\tau)$; and
		\item[\rm (4)]  $_{N}\!\Theta\left(-\frac{u}{\tau},-\frac{1}{\tau} \right)= i^{N^2-1}\cdot_{N}\!\Theta(u,\tau)$.
	\end{itemize}
	
	To prove (1) and (2) we use Proposition {\ref{Theta:transform}}
	$$
	\Theta(u+r \tau + s,\tau) = (-1)^{r+s+rs}e^{-2\pi i r(u+\frac{r\tau+s}{2})}\cdot \Theta(u,\tau)
	$$
	from which we deduce 
	\begin{itemize}
		\item[\rm (i)] $\Theta(u+1,\tau) = -\Theta(u,\tau)$ and $\Theta(N(u+1),\tau) = (-1)^N\Theta(Nu,\tau)$; 
		\item[\rm (ii)] $\Theta(u+\tau,\tau) = -t^{-1}q^{-1/2}\Theta(u,\tau)$  and  $\Theta(N(u+\tau),\tau) = (-t)^{-N^2}q^{-N^2/2}\Theta(Nu,\tau)$.
	\end{itemize}
	Clearly (i) $\Longrightarrow$ (1), and (ii) $\Longrightarrow$ (2) and (3) follows from inspection of the $q$-expansions.  So only (4) remains to be proved.
	
	
	To prove (4) we consider the function
	$$
	\psi(u,\tau) = e^{\pi i u^2/\tau}\cdot\frac{\Theta(u,\tau)}{\Theta\left(-\frac{u}{\tau}, -\frac{1}{\tau}\right)}.
	$$
	A simple calculation using (1) and (2) above shows that $\psi(u+1,\tau) = \psi(u+\tau,\tau)=\psi(u,\tau)$.  Thus $\psi(u,\tau)$ is an elliptic function for the lattice $\Lambda_\tau$.
	But since the numerator and denominator of the above expression for $\psi(u,\tau)$ vanish nowhere outside of $\Lambda_\tau$ and they both have a simple zero at the origin, it follows that
	$\psi(u,\tau)$ has no zeroes or poles in the complex plane. Thus $\psi(u,\tau)$ is a non-zero constant (depending on $\tau$).   To evaluate this constant we calculate $\psi(u,\tau)$ at $u=0$.
	We have
	$$
	\psi(0,\tau) = \displaystyle  \lim_{u\to 0}\frac{\Theta(u,\tau)}{\Theta(-\tfrac{u}{\tau},-\tfrac{1}{\tau})} =\displaystyle -\tau \frac{\Theta^\prime(0,\tau)}{\Theta^\prime(0,-\tfrac{1}{\tau})}
	$$
	Writing 
	$\Theta(u,\tau)	= 2i\sin(\pi u)T(u,\tau)$ where $T(u,\tau)=q^{1/12}\prod_{n>0}(1-q^nt)(1-q^nt^{-1})$, 
	we note that 
	$$
	\Theta^\prime(u,\tau) = 2\pi i\cos(\pi u) \cdot T(u,\tau)  + 2i\sin(\pi u)\cdot T^\prime(u,\tau). 
	$$
	Setting $u=0$ into this identity gives us $\Theta^\prime(0,\tau) = 2\pi i T(0,\tau)$.
	But clearly
	$$
	T(0,\tau) = \eta(\tau)^2
	$$
	where $\eta(\tau)$ is the famous Dedekind $\eta$-function.   Thus
	$$
	\Theta^\prime(0,\tau) = 2\pi i \cdot \eta(\tau)^2  
	$$
	and replacing $\tau$ in this identity by $-1/\tau$ we obtain
	$$
	\Theta^\prime(0,-1/\tau) = 2\pi i \cdot \eta(-1/\tau)^2.
	$$
	It follows now that
	$$
	\psi(0,\tau) = -\tau\cdot\frac{\eta(\tau)^2}{\eta(-1/\tau)^2}
	$$
	and from the well-known identity
	$\eta(\tau)^2 = \frac{i}{\tau} \cdot\eta(-1/\tau)^2$
	we conclude 
	$$
	\psi(0,\tau) = -i.
	$$
	
	We therefore have
	$$
	\Theta\left(-\frac{u}{\tau}, -\frac{1}{\tau}\right) = i e^{2\pi i u^2/\tau} \cdot \Theta(u,\tau)
	$$
	from which it follows at once that
	$$
	_{N}\!\Theta\left(-\frac{u}{\tau}, -\frac{1}{\tau}\right) = \frac{\Theta\left(-\frac{u}{\tau}, -\frac{1}{\tau}\right)^{N^2}} {\Theta\left(-\frac{Nu}{\tau}, -\frac{1}{\tau}\right)} 
	= i^{N^2-1} \frac{\Theta(u,\tau)^{N^2}}{\Theta(Nu,\tau)}  = {i^{N^2-1}} _{N}\!\Theta(u,\tau)
	$$
	and (4) is established.  This completes the proof.
	
\end{proof}



The following corollary follows immediately from Theorem \ref{vartheta:invariance}. Recall that $A(N):=\frac{1}{N}\Z^2/\Z^2$.  
\medskip
\begin{corollary}\label{periodic}
	For $\ba=(a_1+\Z,a_2+\Z) \in A(N)-(0,0)$, define $_{N}\!\Theta_{\ba}: \tfH := \C\times \fH \to \C$ as follows:
	\bea\label{cdiv}
	_{N}\!\Theta_{\ba}(u,\tau):=_{N}\!\Theta (u-a_1\tau-a_2, \tau):=\frac{\Theta(u-a_1\tau-a_2,\tau)^{N^2}}{\Theta(Nu-Na_1\tau-Na_2,\tau)}.
	\eea
	
	(1)
	If $N >0$ is prime to 6, then $_N\Theta(u,\tau)$ is invariant under $\tilde\Gamma=\Z^2 \rtimes \SL_2(\Z)$. In particular $_{N}\!\Theta$ is periodic in $u$ with period $\Lambda_\tau$.
	
	(2) For any positive integer $N$, the function $_N\Theta(u,\tau)^{12}$ is invariant under $\tilde\Gamma=\Z^2 \rtimes \SL_2(\Z)$.
	
	(3) 
	The function $_{N}\!\Theta_{\ba}(u,\tau)^{12}$ is invariant
	%\footnote{We have \be _{N}\!\Theta_{\ba}(\frac{u}{c\tau+d},\frac{a\tau+b}{c\tau+d})^{12} = _{N}\!\Theta_{\ba}(u,\tau)^{12} \ee for $\gamma=\begin{psmallmatrix} a & b \\ c & d \end{psmallmatrix} \in \Gamma(N)$.} 
	under $\tilde \Gamma(N)=\Z^2 \rtimes \Gamma(N)$.
	
\end{corollary}
\begin{proof}
	(1) Since $N >0$ is prime to 6, $N=6k+1$ or $N=6k+5$ for some integer $k$. Then
	\be
	\epsilon(0, T) = e^{2\pi i({N}^2-1)/12}=1\ \ \hbox{\rm and}\ \ \epsilon(0,S) = i^{{N}^2-1}=1.
	\ee
	Thus the result follows.
	
	(2) It is clear, since $\epsilon$ has values in the multiplicative group %$\mu_{12}$ 
	of $12th$ roots of unity.
	
	
	%\begin{cor}\label{periodic}
	%(1)
	% If $c >0$ is prime to 6, then $_c\Theta(u,\tau)$ is invariant under $\tilde\Gamma$. In particular $_{c}\!\Theta$ is periodic in $u$ with period $\Lambda_\tau$.
	% 
	% (2) For any positive integer $c$, the function $_c\Theta(u,\tau)^{12}$ is invariant under $\tilde\Gamma$.
	%\end{cor}
	%\begin{proof}
	%(1) Since $c >0$ is prime to 6, $c=6k+1$ or $c=6k+5$ for some integer $k$. Then
	%\be
	%\epsilon(0, T) = e^{2\pi i({c}^2-1)/12}=1\ \ \hbox{\rm and}\ \ \epsilon(0,S) = i^{{c}^2-1}=1.
	%\ee
	%Thus the result follows.
	%
	%(2) It is clear, since $\epsilon$ has values in $\mu_{12}$.
	%\end{proof}
	
	
	
	
	(3) First, we have
	\be
	_{N}\!\Theta_{\ba} (u + r\tau +s ,\tau)^{12} = _{N}\!\Theta (u-a_1\tau-a_2+ r\tau +s , \tau)^{12}=_{N}\!\Theta_{\ba} (u,\tau)^{12}, \quad r\tau +s \in \Lambda_\tau=\Z\tau+\Z.
	\ee
	Second, we have
	\be
	&&_{N}\!\Theta_\ba \left(\frac{u}{c\tau+d}, \frac{a\tau+b}{c\tau+d}\right)^{12} \\
		&=& _{N}\!\Theta \left(\frac{u}{c\tau+d}-a_1\cdot \frac{a\tau+b}{c\tau+d}-a_2, \frac{a\tau+b}{c\tau+d}\right)^{12}
	\quad (a=1+Nk_1, b=Nk_2, c=Nk_3, d=1+Nk_4)\\
%	&=&  _{N}\!\Theta (\frac{u-a_1\tau-a_2}{c\tau+d}, \frac{a\tau+b}{c\tau+d})^{12}\\
	&=&   _{N}\!\Theta \left(\frac{u-a_1\tau-a_2}{c\tau+d}-\frac{a_1Nk_1\tau+a_2Nk_4}{c\tau+d}-\frac{a_2Nk_3 \tau + a_1Nk_2}{c\tau+d}, \frac{a\tau+b}{c\tau+d}\right)^{12} \\
         &=&  _{N}\!\Theta (u-a_1\tau-a_2 - a_1Nk_1\tau-a_2Nk_4-a_2Nk_3 \tau - a_1Nk_2, \tau)^{12} \\
	&=&   	_{N}\!\Theta_{\ba} (u,\tau)^{12} \quad (a_1Nk_1, a_2Nk_4,a_2Nk_3,a_1Nk_2 \in \Z)
	\ee
	for $\gamma=\begin{pmatrix} a & b \\ c & d \end{pmatrix} \in \Gamma(N)$.
	This implies that $ _{N}\!\Theta_{\ba} (u,\tau)^{12}$ is invariant under $\tilde \Gamma(N)=\Z^2 \rtimes \Gamma(N)$, since the action of $\tilde \Gamma(N)=\Z^2 \rtimes \Gamma(N)$ is generated by the above two actions of $\Gamma(N)$ and $\Lambda_\tau$.
\end{proof}

%We need to show that $_{N}\!\Theta^{12}_\ba$ is a meromorphic function on $\cE_N$.
The functions $_{N}\!\Theta_\ba$ (for $N$ prime to 6) and $_{N}\!\Theta^{12}_\ba$ are meromorphic on $\tilde \fH$ (as explained in Definition \ref{mero}) and invariant under $\tilde \Gamma$ (Corollary \ref{periodic}).  Thus, they are meromorphic functions on the quotient complex manifold $\tilde \fH/\tilde \Gamma$. Similarly,  $_{N}\!\Theta^{12}_\ba$ is a meromorphic function on $\cE_N:=(\C \times \fH)/ \tilde \Gamma(N)$. 
We recall the natural quotient map \eqref{univell}
\be
\pi_N: \cE_N:=(\C \times \fH)/ \tilde \Gamma(N) \to Y_N:=\fH/\Gamma(N)
\ee
and the elliptic surface $\overline \cE(N)$ over $X(N)$ for $N\geq 3$ in the introduction; also recall $\cE(N):=Y(N)\times_{X(N)}\overline \cE(N)$, which is an elliptic curve over $Y(N)$.
We show that $_{N}\!\Theta^{12}_\ba:\cE_N\simeq \cE(N)(\C) \to \C$ extends to a meromorphic function on the compact manifold $\overline \cE(N)(\C)$; see Proposition \ref{rationality} in Appendix \ref{appendix}.
By GAGA,
%\footnote{technically speaking, GAGA is the tool that allows us to turn a global comparison into a local comparison, and hence show that the relevant analytic and algebraic spectral sequences are isomorphic; see \cite{delignesing}, chapter II, section 6 for a clear discussion of the relevant computations (which is a generalization of the very readable article \cite{grothanalg})}
for $N\geq 3$, the function $_{N}\!\Theta^{12}_\ba$ is a rational function on $\overline \cE(N)$, i.e. $_{N}\!\Theta^{12}_\ba$ belongs to $F_N=k_N(E(N))$.   
%; see Appendix \ref{appendix}. Using the isomorphism $\cE(N)(\C) \simeq \cE_N$, we can view $_{N}\!\Theta^{12}_\ba$ as a function on $E(N):=\Spec(k_N)\times_{X(N)}\overline \cE(N)$ over $k_N$. 
%Moreover, the divisor of $_{N}\!\Theta^{12}_\ba$ given in \eqref{divisors} is a finite linear combination of $\C$-points of irreducible algebraic subvarieties of $\cE(N)$ and $_{N}\!\Theta^{12}_\ba$ can be viewed as an element of $F_N=k_N(E(N))$. 
%\begin{remark}
%    Let $Y$ be an affine smooth curve over $\C$ with its compactification $X$ that is a complete smooth curve over $\C$. Then the field $\C(X)$ of rational functions on $X$ is isomorphic to the field of meromorphic functions on $X(\C)$ but the field $\C(Y)$ of rational functions on $Y$ is \textit{not} isomorphic to the field of meromorphic functions on $Y(\C)$. By the way, $\C(X)$ is isomorphic to $\C(Y)$.
%\end{remark}
%Recall the natural quotient map in the introduction (from the universal elliptic curve with full level $N$-structure to the open modular curve)
%\be
%\pi_N: \cE_N:=(\C \times \fH)/ \tilde \Gamma(N) \to Y_N:=\fH/\Gamma(N).
%\ee
%{\color{blue}

Let $N\geq 3$. In order to study the divisors of $_{N}\!\Theta^{12}$ and $_{N}\!\Theta^{12}_\ba$, we define
\bea \label{divs}
D_\bx:=\{[(x_1\tau+x_2, \tau)]_{\tilde \Gamma(N)} \in \cE_N: \tau \in \fH \} \subset \cE_N.
\eea
for $\bx \in A(N)$. Note that $D_\bx$ is the image of the $N$-torsion section $s_\bx$ for the map $\pi_N:\cE_N \to Y_N$ defined by $[\tau]_{\Gamma(N)} \mapsto s_\bx([\tau]_{\Gamma(N)}):=[(x_1\tau+x_2, \tau)]_{\tilde \Gamma(N)}$. Here $[\cdot ]_G$ means the $G$-coset.
 A crucial fact is that, according to the discussion above, \eqref{thetad}, and Corollary \ref{periodic}, $_{N}\!\Theta^{12}=\left( \Theta(u,\tau)^{N^2}/\Theta(Nu,\tau) \right)^{12}$ is a meromorphic function on $\cE_N$ with divisor given by
$$
12 \cdot \left( N^2\cdot D_{\mathbf{0}}  - \sum_{\bx\in A(N)} D_\bx\right)
$$
where $\mathbf{0} =(0+\Z,0+\Z) \in A(N)$, and thus $_{N}\!\Theta_{\ba}^{12}$  in \eqref{cdiv} is a meromorphic function on $\cE_N$ whose divisor is given by
\bea \label{divisors}
12 \cdot \left( N^2\cdot  D_\ba  - \sum_{\bx\in A(N)} {D_\bx}\right).
\eea.
 
%Instead of the above notation, we could also choose to use Scholl?s more suggestive notation:  
%$\ker (\times N) := \sum_{\bx\in A(N)} {D_\bx}$.
 
%}






\begin{proposition}\label{dcomputation}
	%The function $\left(\frac{_{N}\!\Theta_{\ba}}{_{N}\!\Theta_{\bb}}\right)^{12}$ has divisor:
	The function $(_{N}\!\Theta_{\ba}/_{N}\!\Theta_{\bb})^{12}$ defines a meromorphic function\footnote{If one defined $\Theta_{\ba}$ in a similar way without $N$, then $\Theta_{\ba}^{12}/\Theta_{\bb}^{12}$ would not be invariant under $\tilde \Gamma(N)$: for example, by Proposition \ref{Theta:transform}, we have
\be
\frac{\Theta_{\ba}(u+\tau,\tau)}{\Theta_{\bb}(u+\tau,\tau)}= \frac{e^{2\pi i (a_1 \tau +a_2 -\tau/2)}}{e^{2\pi i (b_1 \tau +b_2 -\tau/2)} }\frac{\Theta_{\ba}(u,\tau)}{\Theta_{\bb}(u,\tau)}\neq \frac{\Theta_{\ba}(u,\tau)}{\Theta_{\bb}(u,\tau)}.
\ee
}
on $\cE_N$ whose divisor is given by
	\bea\label{tdiv}
	(12N^2)\cdot \left( D_\ba - D_\bb \right)
	\eea
	where
	$D_\ba=\{[(a_1\tau+a_2, \tau)]_{\tilde \Gamma(N)} \in \cE_N: \tau \in \fH \} \subset \cE_N$ is defined in \eqref{divs}.
	
\end{proposition}


\begin{proof}
This is clear from Corollary \ref{periodic} and \eqref{divisors}.
%	The divisor of $_{N}\!\Theta_{\ba}(u,\tau)^{12}$ is
%	\be
%	12N^2\cdot D_\ba -12 \sum_{N\cdot x\in  \Lambda_\tau} \{[(a_1\tau+a_2+x, \tau)]_{\tilde \Gamma(N)} : \tau \in \fH\}.
%	\ee
%	Since, for $\ba, \bb \in \frac{1}{N}\Z^2/\Z^2-(0,0)$, we have
%	\be
%	\sum_{N\cdot x\in  \Lambda_\tau} \{[(a_1\tau+a_2+x, \tau)]_{\tilde \Gamma(N)} : \tau \in \fH\}= \sum_{N\cdot x\in  \Lambda_\tau} \{[(b_1\tau+b_2+x, \tau)]_{\tilde \Gamma(N)} : \tau \in \fH\},
%	\ee
%	the result follows.
\end{proof}





%\subsection{Siegel units}
%Let $A(N) := {1\over N} \Z^2/\Z^2$ and $A(N)^\prime := A(N)\setminus\{0\}$.  Following Kato (\cite[subsection 1.9]{Kat}), we define the Siegel units of level $N$.
%
%\begin{definition}
%For $\ba = (a_1 + \Z, a_2 +\Z)\in A(N)$ with $0\le a_1,a_2 < 1$ we define $g_\ba = 1$ if $\ba = {\bf 0}$ and otherwise define
%$$
%g_\ba(\tau) := q^{\frac{1}{2}B_2(a_1)} \prod_{n\ge 0} \left(1 - q^{n+a_1} e^{2\pi i a_2} \right) \prod_{n\ge 1} \left(1 - q^{n-a_1} e^{-2\pi i a_2} \right)
%$$
%where $B_2(x)=x^2-x+\frac{1}{6}.$
%\end{definition}


\begin{lemma}\label{Theta:siegel} For $\ba \in A(N)-(0,0)$ and $(a_1,a_2) \in \frac{1}{N}\Z^2$ the least non-negative representative of $\ba$, we have
	\begin{itemize}
		\item[\rm (a)]  $\displaystyle -e^{\pi i a_2} q^{\frac{a_1^2}{2}}\cdot \Theta\left(a_1\tau + a_2,\tau \right) = g_\ba(\tau);$ and
		\item[\rm (b)]  $\displaystyle \frac{\Theta(Nu + Na_1\tau+Na_2, \tau)}{ \Theta(Nu,\tau)} = (-1)^{Na_2} (-t^{N})^{-Na_1}q^{-\frac{N^2a_1^2}{2}}$, where $t=e^{2\pi iu}$ and $q=e^{2\pi i \tau}$.
	\end{itemize}
	
\end{lemma}

\begin{proof}  Both of these are straightforward calculations.   For (a) we have  
	$$
	\begin{array}{rcl}
		%e^{\pi i a_2} q^{\frac{a_1^2}{2}}\cdot 
		\Theta\left(a_1\tau + a_2,\tau \right) &=&\displaystyle q^{\frac{1}{12}}\left( e^{\pi i a_2}q^{\frac{a_1}{2}} - e^{-\pi i a_2}q^{\frac{-a_1}{2}}\right)
		\prod_{n\ge 1}\left( 1-q^{n+a_1}e^{2\pi i a_2}\right)\left( 1-q^{n-a_1}e^{-2\pi i a_2}\right)\\
		&=&\displaystyle -q^{\frac{1}{12}}e^{-\pi i a_2}q^{\frac{-a_1}{2}} \prod_{n\ge 0}\left( 1-q^{n+a_1}e^{2\pi i a_2}\right)\prod_{n\ge 1} \left( 1-q^{n-a_1}e^{-2\pi i a_2}\right)\\
		&=&\displaystyle -e^{-\pi i a_2}q^{-\frac{a_1^2}{2}}q^{\frac{1}{2}(a_1^2-a_1+\frac{1}{6})} \prod_{n\ge 0}\left( 1-q^{n+a_1}e^{2\pi i a_2}\right)\prod_{n\ge 1} \left( 1-q^{n-a_1}e^{-2\pi i a_2}\right)\\
		&=&-e^{-\pi i a_2}q^{-\frac{a_1^2}{2}}g_\ba
	\end{array}
	$$
	The proof of (b) is a simple substitution in Proposition \ref{Theta:transform} using the fact that $Na_1, Na_2\in \Z$. 
\end{proof}




\begin{proposition}\label{vartheta:siegelspecial}
%{\color{blue}
For any non-zero $\ba,\bb \in A(N)$ the meromorphic function  $(_{N}\!\Theta_{\ba}/_{N}\!\Theta_{\bb})^{12}$ on $\cE_N$ is holomorphic on the $\mathbf{0}$-section $s_{\mathbf{0}}$, and its pullback to $Y_N$ along $s_{\mathbf{0}}$ is
$
 \left(g_\ba / g_\bb\right)^{12N^2}.
$
%}
%	Let $\ba,\bb\in A(N)-(0,0)$ with $\ba = (a_1+\Z, a_2+\Z)$ and $\bb = (b_1+\Z, b_2+\Z)$ where $(a_1,a_2)$ and $(b_1,b_2)$ are the least non-negative representatives for $\ba$ and $\bb$ respectively.  Then $\displaystyle {_{N}\!\Theta(u+ a_1\tau + a_2,\tau)}\slash {_{N}\!\Theta(u+b_1\tau+b_2,\tau)}$ is holomorphic at $u=0$ and 
%	$$
%	\frac{_{N}\!\Theta(u+ a_1\tau + a_2,\tau)}{_{N}\!\Theta(u+b_1\tau+b_2,\tau)}\biggm|_{u=0} = \pm \left(\frac{g_\ba(\tau)}{g_\bb(\tau)}\right)^{N^2}
%	$$
%	where the sign $\pm$ is given by $\pm 1 = (-1)^{N(a_1+a_2 + b_1+ b_2) + N^2(a_2+b_2)}$.
\end{proposition}

\begin{proof}  
	Recall the definition $_{N}\!\Theta(u,\tau) :=  \Theta(u,\tau)^{N^2}/\Theta(Nu,\tau)$.
	This proposition follows from the following computation: from lemma \ref{Theta:siegel} we have
	
	$$
	\begin{array}{rcl}
		\displaystyle \frac{_{N}\!\Theta(u+ a_1\tau + a_2,\tau)}{_{N}\!\Theta(u+b_1\tau+b_2,\tau)}\biggm|_{u=0} 
		&=& \displaystyle \frac{\Theta(a_1\tau + a_2,\tau)^{N^2}}{\Theta(b_1\tau +b_2,\tau)^{N^2}} \cdot \frac{\Theta(Nu+Nb_1\tau + Nb_2,\tau)}{\Theta(Nu + Na_1\tau +Na_2,\tau)}\biggm|_{u=0}\\
		&=& \displaystyle \left(\frac{e^{-\pi i a_2} q^{-\frac{a_1^2}{2}}g_\ba}{e^{-\pi i b_2} q^{-\frac{b_1^2}{2}}g_\bb}\right)^{N^2}\cdot \frac{\Theta(Nu+Nb_1\tau + Nb_2,\tau)}{\Theta(Nu + Na_1\tau +Na_2,\tau)}\biggm|_{u=0}\\
		&=& \displaystyle \left(\frac{e^{-\pi i a_2} q^{-\frac{a_1^2}{2}}g_\ba}{e^{-\pi i b_2} q^{-\frac{b_1^2}{2}}g_\bb}\right)^{N^2}\cdot 
		\frac{(-1)^{Nb_2} (-t^{N})^{-Nb_1}q^{-\frac{N^2b_1^2}{2}}}{(-1)^{Na_2} (-t^{N})^{-Na_1}q^{-\frac{N^2a_1^2}{2}}}\biggm|_{t=1}\\
		&=& \displaystyle (-1)^{N(a_1+a_2 + b_1+ b_2) + N^2(a_2+b_2)}\left(\frac{g_\ba}{g_\bb}\right)^{N^2}.\\
%		&=& \displaystyle \pm\left(\frac{g_\ba}{g_\bb}\right)^{N^2}.\\
	\end{array}
	$$
\end{proof}
\begin{corollary}\label{TS}
	For $\ba=(a_1+\Z,a_2+\Z), \bb=(b_1+\Z,b_2+\Z) \in \frac{1}{N}\Z^2/\Z^2-(0,0)$ and $\bx=(x_1+\Z,x_2+\Z) \in \frac{1}{N}\Z^2/\Z^2-(0,0)$ which is different from $\ba,\bb$, we have
	\be
	\frac{_{N}\!\Theta ((x_1-a_1)\tau+(x_2-a_2), \tau)^{12}}{_{N}\!\Theta ((x_1-b_1)\tau+(x_2-b_2), \tau)^{12}} =\left(\frac{g_{\bx-\ba}(\tau)}{g_{\bx-\bb}(\tau)}\right)^{12N^2}.
	\ee
\end{corollary}

%Using Corollary \ref{periodic}, we define, for a positive integer $c$ and $\ba=(a_1+\Z, a_2+\Z) \in \Q^2/\Z^2-(0,0)$,
%\be
%_{c}\!g_{\ba}(\tau):=_{c}\!\Theta(a_1\tau+a_2, \tau)=\frac{\Theta(a_1\tau+a_2, \tau)^{c^2}}{\Theta(ca_1\tau+ca_2,\tau)}.
%\ee









\subsection{The Manin relations}\label{ManinRelsProof}
For simplicity of notation, we define:
\be
\theta_{\ba}(u, \tau):=_{N}\!\Theta_{\ba}(u,\tau)^{12},\quad
_{N}\!\Theta_{\ba}(u,\tau)=_{N}\!\Theta (u-a_1\tau-a_2, \tau)=\frac{\Theta(u-a_1\tau-a_2,\tau)^{N^2}}{\Theta(Nu-Na_1\tau-Na_2,\tau)},
\ee
where $_{N}\!\Theta_{\ba}(u,\tau)^{12}$ is an element of $F_N$ ($N\geq 3$) by Corollary \ref{periodic}.  
% a function on $\cE_N=\tfH /\tilde\Gamma(N)$ given in \eqref{cdiv}.
%
%
%For simplicity of notations, we suppress $c$ from the notation:
%\be
%\theta_{\ba}(u, \tau):=_{c}\!\Theta_{\ba}(u,\tau)
%\ee
%where $_{c}\!\Theta_{\ba}(u,\tau)$ is a function on $\tfH /\tilde \Gamma$ given in \eqref{cdiv}.
%A meromorphic modular function $f$ is said to be rational over a subfield $k$ of $\C$ if its fourier coefficients (at $\infty$) all lie in the field $k$.   We define $\cF_N$ to be the field of all meromorphic modular functions rational over $\Q(\mu_N)$
%and defines the field of all modular functions over cyclotomic fields to be the field 
%$$
%\cF := \bigcup_{N} \cF_N.
%$$
%%   In Chapter 6 of his book, Shimura gives an explicit finite set of generators over $\Q$ of $\cF_N$ (see section 6.2 and Proposition 6.9 of \cite{shimura}).  More precisely, he shows $\cF_1 = \Q(j)$ and for $N>1$, that $\cF_N$ is generated over $\cF_1$ by the Fricke functions $\left\{\,f_a \bigm| a\in \frac{1}{N}\Z^2/\Z^2\,\right\}$.    Furthermore, he defines the field of all modular functions over cyclotomic fields to be the field 
%%$$
%%\cF := \bigcup_{N} \cF_N.
%%$$
%For each natural number $N$ we let $\cR_N$ be the subring of $\cF_N$ consisting of functions that are holomorphic on the upper half-plane and let
%$$
%\cR := \bigcup_N{\cR_N}.
%$$
%%Then both $\cR$ and its unit group $\cR^\times$ are preserved by the action of $Aut(\cF)$ on $\cF$, thus $\cR^\times$ inherits a natural action of $\GL_2(\A_f)/\Q^\times$.
%For $\ba \in \frac{1}{N}\Z^2/\Z^2-(0,0)$, we have
%\be
%_{c}\!g_{\ba} \in \cR_N^\times \subset \left(\cM(\Gamma(N), \Q(\mu_N))[\frac{1}{\Delta}]\right)^\times, \quad
%g_{\ba} \in \Q \otimes \cR_N^\times \subset  \Q \otimes (\cM(\overline \Q)[\frac{1}{\Delta}])^\times.
%\ee
%Recall the natural map $\pi_N: \cE_N \to Y_N$ given in \eqref{univell}.
%Let $\C(\cE)$ be the field of meromorphic functions on $\cE_N$ for all $N \geq 1$ and $\C(Y)$ the field of meromorphic functions on $Y_N$ for all $N \geq 1$. 
%Let $\cE_N(\C) \to Y(N)(\C)$ be the universal elliptic curve where $\cE_N(\C)=\tfH /\tilde \Gamma(N)$ and $Y(N)(\C)=\frak H/\Gamma(N)$.
%Then for $\ba, \bb \in \frac{1}{N}\Z^2/\Z^2-(0,0)$, the function $\frac{\theta_{\ba}}{\theta_\bb}=\left(\frac{_{N}\!\Theta_{\ba}}{_{N}\!\Theta_{\bb}}\right)^{12}$ in \eqref{tdiv} is an element of $\C(\cE_N)$.
%where $\C(\cE_N)$ is the field of meromorphic functions on $\cE_N(\C)$.
%Let $\C(\cE)$ be the field of meromorphic functions on $\cE_N(\C)$ for all $N \geq 1$ and $\C(Y)$ the field of meromorphic functions on $Y(N)(\C)$ for all $N \geq 1$. 
For $\ba, \bb \in A(N)$, we have
\be
\frac{\theta_{\ba}}{\theta_\bb} \in F_N^\times,  \quad
g_{\ba} \in \tilde K^M_1(k_N).
%, \quad _{c}\!g_{\ba} \in \C(Y)^\times \subset \Frac(\cM(\C))^\times.
\ee
%Let $\overline{\C(\cE)}$ (resp. $\overline{\C(Y)}$) be the algebraic closure of $\C(\cE)$ (resp. $\C(Y)$).
%Then it is crucial to study the map $\Res: K_3(\overline{\C(\cE)}) \to K_2(\overline{\C(Y)})$ in order to understand our modular symbol. 


%We introduce
%\be
%\cK_*(A) := K_*(A)/\langle \{-1\} \rangle
%\ee
%where $\langle \{-1\} \rangle$ is the homogeneous ideal of $K_*(A)$ generated by $\{-1\}$. Then we have the decomposition $\cK_*(A) =\oplus_{n\geq 0} \cK_n(A)$.

Since $g_{-\ba}=-e^{2 \pi i a_2} g_{\ba}$, we have $\{g_{-\ba}, g_{\bb}\}\equiv \{g_{\ba}, g_{\bb}\}$ in $\tilde K^M_2(k_N)$.
For $\ba \in A(N)$, the divisor
$D_\ba$
%:=\{[(a_1\tau+a_2, \tau)]_{\tilde \Gamma(N)} \in \cE_N: \tau \in Y_N \} \subset \cE_N$
provides a discrete valuation on $F_N$ which is trivial on $k_N$. So we use the notation $\partial_\ba$ instead of $\partial_{v_\ba}$ where $v_\ba$ is the associated discrete valuation to $\ba$.
%Let $\Frac(\cM(\overline\Q))$ be the fraction field of $\cM(\overline\Q)$.


\begin{lemma} \label{bcomputation}
	For $N\geq 3$, let $\ba,\bb,\bc,\bx \in A(N)-(0,0)$ be distinct non-zero elements. We have the following identities in $\tilde K^M_2(k_N)$:
	$$
	\begin{array}{rclcl} \displaystyle
		\frac{1}{(12N^2)^3} {\partial}_\ba \left( \left\{\frac{\theta_{\ba}}{\theta_\bx},\frac{\theta_{\bb}}{\theta_\bx},\frac{\theta_{\bc}}{\theta_\bx}\right\}\right)
		&=& \displaystyle \left\{ \frac{g_{\ba-\bb}}{g_{\ba-\bx}}, \frac{g_{\bc-\ba}}{g_{\bx-\ba}} \right\} &=& \{g_{\ba-\bb},g_{\bc-\ba}\} - \{g_{\ba - \bb}, g_{\bx-\ba}\}-\{g_{\ba-\bx},g_{\bc-\ba}\} ;\\
		\\
		\displaystyle \frac{1}{(12N^2)^3} {\partial}_\bb \left( \left\{\frac{\theta_{\ba}}{\theta_\bx},\frac{\theta_{\bb}}{\theta_\bx},\frac{\theta_{\bc}}{\theta_\bx}\right\}\right)
		&=& \displaystyle \left\{ \frac{g_{\bb-\bc}}{g_{\bb-\bx}}, \frac{g_{\ba-\bb}}{g_{\bx-\bb}} \right\}
		&=&  \{g_{\bb-\bc}, g_{\ba-\bb}\} - \{g_{\bb-\bc}, g_{\bx-\bb}\} - \{g_{\bb-\bx},g_{\ba-\bb}\} ;\\
		\\
		\displaystyle \frac{1}{(12N^2)^3} {\partial}_\bc \left( \left\{\frac{\theta_{\ba}}{\theta_\bx},\frac{\theta_{\bb}}{\theta_\bx},\frac{\theta_{\bc}}{\theta_\bx}\right\}\right)
		&=&\displaystyle\left\{ \frac{g_{\bc-\ba}}{g_{\bc-\bx}}, \frac{g_{\bb-\bc}}{g_{\bx-\bc}} \right\} 
		&=& \{g_{\bc-\ba},g_{\bb-\bc} \} - \{g_{\bc -\ba},g_{\bx-\bc}\} - \{g_{\bc - \bx}, g_{\bb-\bc}\};\\
		\\
		\displaystyle \frac{1}{(12N^2)^3} {\partial}_\bx \left( \left\{\frac{\theta_{\ba}}{\theta_\bx},\frac{\theta_{\bb}}{\theta_\bx},\frac{\theta_{\bc}}{\theta_\bx}\right\}\right) 
		&=&\displaystyle \left\{ \frac{g_{\bx-\bc}} {g_{\bx-\ba}}, \frac{g_{\bx-\bb}}{g_{\bx-\ba}}\right\}
		&=&-\left\{ g_{\bx-\bc},g_{\ba- \bx}\right\} - \{g_{\bx -\ba}, g_{\bb-\bx}\} - \{ g_{\bx -\bb}, g_{\bc - \bx}\} \\
	\end{array}
	$$
	where $\partial_*: \tilde K^M_3(F_N) \to \tilde K^M_2(k(*))$ is the boundary map in Lemma \ref{boundary}.
	Consequently, we have
	%\footnote{The left hand side ${\Res}\left( \left\{\frac{\theta_{\ba}}{\theta_\bx},\frac{\theta_{\bb}}{\theta_\bx},\frac{\theta_{\bc}}{\theta_\bx}\right\}\right) $ is zero by the Weil reciprocity law by Corollary \ref{ac}, since $\C$ is algebraically closed.}
	$$
	\frac{1}{(12N^2)^3} {\Res}\left( \left\{\frac{\theta_{\ba}}{\theta_\bx},\frac{\theta_{\bb}}{\theta_\bx},\frac{\theta_{\bc}}{\theta_\bx}\right\}\right) = 
	\theta(\ba:\bb:\bc) - \theta(\ba:\bb:\bx) - \theta(\ba:\bx:\bc) - \theta(\bx:\bb:\bc)
	$$
	where $\Res:\tilde K^M_3(F_N) \to \tilde K^M_2(k_N)$ is the residue map in Theorem \ref{SRL} and we put 
	$$
	\theta(\ba:\bb:\bc) := \{g_{\ba-\bb},g_{\bc-\ba}\} +  \{g_{\bc-\ba},g_{\bb-\bc}\} +  \{g_{\bb-\bc},g_{\ba-\bb}\} \in \tilde K^M_2(k_N).
	$$
	
\end{lemma}


\begin{proof} 
	This follows from calculations based on the definition of the residue map and the properties of the Steinberg symbol.
	More explicitly, we use the following computation
	\be
	\partial_v (\{\pi_v^{r_1} u_1, \pi_v^{r_2}u_2, \pi_v^{r_3}u_3 \})=
	r_1 \{\overline u_2, \overline u_3\}-r_2\{\overline u_1, \overline u_3 \} + r_3 \{ \overline u_1, \overline u_2\}
	\ee
	where $v$ is a discrete valuation on $F_N$ which is trivial on $k_N$ and $\pi_v$ is a uniformizer with $v$-adic units $u_1, u_2, u_3$ (with notations from Lemma \ref{boundary}).
    Note that discrete valuations on $F_N$ which are trivial on $k_N$ are in bijection with prime divisors of $E(N)$,
	 which is the elliptic curve (in particular, the algebraic normal irreducible variety) over $k_N$ in the introduction:
for an irreducible normal variety $X$ and a prime divisor $D \subset X$, $\cO_{X ,D}$ is the subring of the function field $\C(X)$ defined by
\be
\cO_{X ,D}:= \{\varphi \in \C(X ) \ | \ \varphi \text{ is defined on } U \subset X \text{ open with } U \cap D\neq \phi \}	
\ee	
which is a discrete valuation of $\C(X)$ corresponding to $D$.
	The first four identities follow from Proposition \ref{dcomputation} and Corollary \ref{TS}. For the readers' convenience, we include the proof of the first identity:
	{\small{
	\be
	  {\partial}_\ba \left( \left\{\frac{\theta_{\ba}}{\theta_\bx},\frac{\theta_{\bb}}{\theta_\bx},\frac{\theta_{\bc}}{\theta_\bx}\right\}\right) 
	&=& 12N^2 \left\{ \overline{\frac{\theta_{\bb}}{\theta_{\bx}}},  \overline{\frac{\theta_{\bc}}{\theta_{\bx}}} \right\}   \\
	&=& 12N^2 \left\{ \frac{_{N}\!\Theta(a_1\tau + a_2-b_1\tau-b_2,\tau)^{12}}{_{N}\!\Theta(a_1\tau+a_2-x_1\tau-x_2,\tau)^{12}} ,  \frac{_{N}\!\Theta( a_1\tau + a_2-c_1\tau-c_2,\tau)^{12}}{_{N}\!\Theta(a_1\tau+a_2-x_1\tau-x_2,\tau)^{12}} \right\}     \\
		&=&(12N^2)^3 \displaystyle \left\{ \frac{g_{\ba-\bb}}{g_{\ba-\bx}}, \frac{g_{\bc-\ba}}{g_{\bx-\ba}} \right\}.
	\ee
	}}
	Then we compute $\Res$ using these first four identities:
	\be
	&&\Res\left( \left\{\frac{\theta_{\ba}}{\theta_\bx},\frac{\theta_{\bb}}{\theta_\bx},\frac{\theta_{\bc}}{\theta_\bx}\right\}\right)\\
	&:=&\sum_{v \in \Sigma(F_N/k_N)} \partial_v \left(\left\{\frac{\theta_{\ba}}{\theta_\bx},\frac{\theta_{\bb}}{\theta_\bx},\frac{\theta_{\bc}}{\theta_\bx}\right\} \right) \\
	&=&{\partial}_\ba \left( \left\{\frac{\theta_{\ba}}{\theta_\bx},\frac{\theta_{\bb}}{\theta_\bx},\frac{\theta_{\bc}}{\theta_\bx}\right\}\right)
	+{\partial}_\bb \left( \left\{\frac{\theta_{\ba}}{\theta_\bx},\frac{\theta_{\bb}}{\theta_\bx},\frac{\theta_{\bc}}{\theta_\bx}\right\}\right)
	+{\partial}_\bc \left( \left\{\frac{\theta_{\ba}}{\theta_\bx},\frac{\theta_{\bb}}{\theta_\bx},\frac{\theta_{\bc}}{\theta_\bx}\right\}\right)+
	{\partial}_\bx \left( \left\{\frac{\theta_{\ba}}{\theta_\bx},\frac{\theta_{\bb}}{\theta_\bx},\frac{\theta_{\bc}}{\theta_\bx}\right\}\right) \\
	&=&
	(12N^2)^3 \left(\displaystyle \left\{ \frac{g_{\ba-\bb}}{g_{\ba-\bx}}, \frac{g_{\bc-\ba}}{g_{\bx-\ba}} \right\} +
	\displaystyle \left\{ \frac{g_{\bb-\bc}}{g_{\bb-\bx}}, \frac{g_{\ba-\bb}}{g_{\bx-\bb}} \right\} +
	\displaystyle\left\{ \frac{g_{\bc-\ba}}{g_{\bc-\bx}}, \frac{g_{\bb-\bc}}{g_{\bx-\bc}} \right\} +
	\displaystyle \left\{ \frac{g_{\bx-\bc}} {g_{\bx-\ba}}, \frac{g_{\bx-\bb}}{g_{\bx-\ba}}\right\}  \right) \\
	&=&
	(12N^2)^3 \biggl(\theta(\ba:\bb:\bc) - \theta(\ba:\bb:\bx) - \theta(\ba:\bx:\bc) - \theta(\bx:\bb:\bc)\biggr).
	\ee
\end{proof}

Next we state a lemma closely related to, but with slightly different conditions from, \cite[Lemma 2.11]{Gon} which we will apply in Proposition \ref{klem} to $\theta(\ba:\bb:\bx)$ and other terms.\footnote{The first condition in \cite[Lemma 2.11]{Gon} does not apply to our version of $\theta$ and does not appear to apply to the version of $\theta$ defined in \cite{Gon} either}

Let $A$ and $B$ be abelian groups. Let \begin{align*}
\psi: \wedge^3\Z[A]\to B, \>\>\> {\ba}\wedge{\bb}\wedge{\bc}\mapsto \psi(\ba:\bb:\bc)
\end{align*} 
be a group homomorphism.

\begin{lemma}\label{gonchvariantlemma}
Let $\ba,\bb,$ and $\bc$ be elements of $A$ (and $-\ba, -\bb,-\bc \in A$ are their inverses in $A$).
Assume that the group $A$ is finite and that the map $\psi$ satisfies
\begin{itemize}
        %\item $\psi(\ba: \bb: \bc)=-\psi(\bb: \ba: \bc)$
	\item $\psi(-\ba:-\bb:-\bc)=\psi(\ba:\bb:\bc)$
	\item $\psi(\ba + \bx, \bb + \bx, \bc + \bx) = \psi(\ba , \bb , \bc )$ for any $\bx\in A$.
\end{itemize}
Then modulo 2-torsion one has \be
\sum_{\bx \in A} \psi(\ba:\bb:\bx)=\sum_{\bx \in A}\psi(\ba:\bx:\bc)=\sum_{\bx \in A}\psi(\bx:\bb:\bc) =0.
\ee
\end{lemma}
\begin{proof}
The proof (if not the statement) is identical to the proof of \cite[Lemma 2.11]{Gon}.  Indeed, by the group homomorphism structure of $\psi$, it suffices to prove only 
$\sum_{\bx \in A} \psi(\ba:\bb:\bx)=0$.  But 
\be
\sum_{\bx \in A} \psi(\ba:\bb:\bx)
&=&\sum_{\bx \in A}\psi(\ba-(\ba+\bb):\bb-(\ba+\bb):\bx-(\ba+\bb))\\
&=&\sum_{\by \in A}\psi(-\bb:-\ba:-\by)=-\sum_{\by \in A}\psi(-\ba:-\bb:-\by)=-\sum_{\by \in A}\psi(\ba:\bb:\by).
\ee
\end{proof}
We apply the lemma when $A=A(N):= \frac{1}{N}\Z^2/\Z^2$.
\begin{proposition}\label{klem}
	Let $N\ge 3$ an integer, and let $\ba,\bb,\bc \in A(N)-(0,0)$ be distinct non-zero elements. We have
	\be
	\{g_{\ba-\bb},g_{\bc-\ba}\} +  \{g_{\bc-\ba},g_{\bb-\bc}\} +  \{g_{\bb-\bc},g_{\ba-\bb}\} =0 \quad \text{in} \quad \tilde K^M_2(k_N).
	\ee
\end{proposition}

\begin{proof}  
	The proof follows from Lemma \ref{bcomputation}, and the following observation:
	Summing the terms in Lemma \ref{bcomputation} over $\bx \in A(N)$ and applying Lemma \ref{gonchvariantlemma} 
	%\cite[Lemma 2.11]{Gon}\footnote{
		%In other words,
	%	\be
	%	\sum_{\bx \in A(N)} \theta(\ba:\bb:\bx)=\sum_{\bx \in A(N)}\theta(\ba:\bx:\bc)=\sum_{\bx \in A(N)}\theta(\bx:\bb:\bc) =0.
	%	\ee} 
	to $\theta(\ba:\bb:\bx), \theta(\ba:\bx:\bc),$ and $\theta(\bx:\bb:\bc)$, we obtain
	$$
	\frac{1}{N^2(12N^2)^3} \sum_{\bx \in A(N)} {\Res}\left( \left\{\frac{\theta_{\ba}}{\theta_\bx},\frac{\theta_{\bb}}{\theta_\bx},\frac{\theta_{\bc}}{\theta_\bx}\right\}\right)=\{g_{\ba-\bb},g_{\bc-\ba}\} +  \{g_{\bc-\ba},g_{\bb-\bc}\} +  \{g_{\bb-\bc},g_{\ba-\bb}\}.
	$$
	Then, by Suslin's reciprocity law it follows that $\{g_{\ba-\bb},g_{\bc-\ba}\} +  \{g_{\bc-\ba},g_{\bb-\bc}\} +  \{g_{\bb-\bc},g_{\ba-\bb}\}= 0$.
	%, which is a straightforward calculation from the definition of the residue map and the properties of the Steinberg symbol.
%	\begin{remark}\label{srescomment}




\end{proof}






\begin{theorem}\label{sres}
	Let $\ba, \bb, \bc \in \Q^2/\Z^2-(0,0)$ and suppose $\ba + \bb + \bc = 0$.   Then
	$$
	\{g_\ba,g_\bb\} + \{g_\bb,g_\bc\} + \{g_\bc,g_\ba\} = 0\ \ \hbox{\rm in $\tilde K^M_2(k_\infty) $}.
	$$
\end{theorem}
	\begin{remark}\label{srescomment}
	
	Note that Theorem \ref{sres} is still true for $\ba, \bb, \bc$  %holds in $\tilde K^M_2(k_\infty)$ 
	associated to $N=1,2$: the $N=1$ case holds trivially (as $g_{(0,0)}=1$) and the $N=2$ case has only one Manin relation, which follows directly from the $N=4$ case by observing that
	$g_{(0,1/2)}, g_{(1/2,0)}, g_{(1/2,1/2)}$
	have the same $q$-expansions as $g_{(0,2/4)}, g_{(2/4,0)}, g_{(2/4,2/4)}$ respectively.  
\end{remark}

%{\color{red} 
%{\color{blue} 
\begin{remark} 
We remark that Theorem \ref{sres} is equivalent to 
\cite[Theorem 4.1]{Bru2} modulo torsion, which strengthens\footnote{It removes the assumption that $3$ does not divide $N$ in \cite{Bru}.} \cite[(7) Theorem 1.4, (5) Definition 1.2]{Bru}, the main result of \cite{Bru}.
%\cite[(7) Theorem 1.4, (5) Definition 1.2]{Bru}, the main result of \cite{Bru}.  
Indeed:
\begin{itemize}
	\item Let $Y(N)^{(1)}$ denote the set of points of $Y(N)$ of codimension $1$. It follows e.g. from \cite[Section 2.1, especially (2.5)]{SV} that the kernel of the tame symbol map
	$K^M_2(k_N) \xrightarrow{\partial} \bigoplus_{D\in Y(N)^{(1)}} K_1^M(\C(D))$ is isomorphic to the motivic cohomology group $H^2(Y(N), 2)$.
	\item We identify motivic cohomology groups with the higher Chow groups 
	\be
	\mathrm{CH}^p(Y(N), n) = H^{2p-n} (Y(N), p)
        \ee
         as per standard practice in the literature, and we observe (e.g. \cite[p. 177]{Totaro} or \cite{Bloch}) that the generalized Grothendieck Riemann Roch theorem 
	\be
	K_n(Y(N)) \otimes \Q \simeq \bigoplus_p \mathrm{CH}^p(Y(N), n) \otimes \Q
	\ee
	induces an embedding 
		$H^2(Y(N),2)\otimes \Q \hookrightarrow K_2(Y(N))\otimes \Q$.
%	$H^2(\cM(N),2)\otimes \Q \hookrightarrow K_2^M(\Q(\cM(N)))\otimes \Q$.
\end{itemize}
Since the $g_\ba$ are units (and hence have trivial divisor) for $\ba$ as above,  the tame symbol of $\{g_\ba, g_\bb\}$ is zero by definition, and so $\{g_\ba, g_\bb\}$ belongs to the Quillen $K_2$ of the modular curve of the appropriate level  (modulo tensoring with $\Q$).
%Note that since the $g_\ba$ are units (and hence have trivial divisor) for $\ba$ as above,  the tame symbol of $\{g_\ba, g_\bb\}$ is zero by definition, and it belongs to the Quillen $K_2$ of the modular curve of the appropriate level (modulo tensoring with $\Q$); \cite[(2.5)]{SV} says that the kernel of the tame symbol map $K_2(k_N) \xrightarrow{\partial} \bigoplus_{D\in Y(N)_1} K_1^M(\C(D))$ is isomorphic to the motivic cohomology group $H^2(Y(N), 2)$ recalled in \cite[Section 2.1]{SV} (here $Y(N)_1$ is the set of points of $Y(N)$ of codimension $1$) and $H^2(Y(N), 2) \otimes \Q$ is embedded  in the Quillen $K$-group $K_2(Y(N))\otimes \Q$ by \cite[p. 177]{Totaro}
%\be
%K_n(Y(N)) \otimes \Q \simeq \bigoplus_p \mathrm{CH}^p(Y(N), n) \otimes \Q, \quad \mathrm{CH}^p(Y(N), n) = H^{2p-n} (Y(N), p).
%\ee
%; we refer to \cite[Theorem 2.2.1]{SV} for details.
In this paper we use the Quillen $K$-groups of the irreducible modular curve $Y(N)$ over $\C$, whereas  Brunault \cite{Bru} uses that of the irreducible modular curve $\cM(N)$ over $\Q$ (see \eqref{cm} and \cite[p. 2]{Bru}). There is a natural embedding from the function field $\Q(\cM(N))$ to $k_N=\C(Y(N))$, which induces a map $K_2^M(\Q(\cM(N))) \to K_2^M(k_N)$. According to \cite[Lemma 6.1.3, Chap 3]{Weibel}, for every field extension $F \subset E$, the kernel of $K_2^M(F) \to K_2^M(E)$
is a torsion subgroup; this implies that $K_2^M(\Q(\cM(N)))\otimes \Q \to K_2^M(k_N)\otimes \Q$ is injective. Through the composite map $$H^2(\cM(N),2)\otimes \Q \hookrightarrow K_2^M(\Q(\cM(N)))\otimes \Q \hookrightarrow K_2^M(k_N)\otimes \Q,$$ it follows that our Theorem \ref{sres} is equivalent to \cite[Theorem 4.1]{Bru2} modulo torsion.\footnote{Brunault's result is stated in $K_2^M(\Q(\cM(N)))\otimes \Z[\frac{1}{6N}]$.}

%which strengthens \cite[(7) Theorem 1.4, (5) Definition 1.2]{Bru}, the main result of \cite{Bru}.
%it further induces a map $K_2(\cM(N))\otimes \Q \to K_2(Y(N)) \otimes \Q$.
%Since the natural map $K_2(\cM(N)) \to K_2^M(\Q(\cM(N)))$ is injective, there is a map $K_2(\cM(N))$ to $K_2(Y(N))$.
%Hence this theorem is equivalent to \cite[(7) Theorem 1.4, (5) Definition 1.2]{Bru}, the main result of \cite{Bru}, through this map. %except that we use the irreducible modular curve $Y(N)$ over $\C$ and \cite{Bru} uses the irreducible modular curve $\cM(N)$ over $\Q$ (see \eqref{cm}).
%Furthermore, given a compatible system of symbols, the system belongs to the appropriate tower of modular curves (again, modulo tensoring with $\Q$). 
%Hence this theorem is equivalent to \cite[(7)]{Bru}, the main result of \cite{Bru}.
\end{remark}
%, as the author of \cite{Bru} defines his units $g_\ba$ the same way as they are defined in this paper.
%}
%Note that this corollary is equivalent to \cite[(7)]{Bru}, the main result of \cite{Bru}. 
%}

%{\color{blue} This corollary is equivalent to \cite[(7)]{Bru}, the main result of Brunault; we have given a new proof of his result using Suslin's reciprocity law (Theorem \ref{SRL}) based on Goncharov's idea.}




\subsection{The construction and the proof}
%\subsection{The Beilinson-Kato moduar symbol}


%We note that the group $\GL_2(\Q)$ acts by matrix multiplication both on the left and on the right on $M_{2\times2}(\Q)$.
%% and that this induces a two-sided action of $\GL_2(\A_f)$ on $\cS(M_{2\times2}(\Q))$.  
%Correspondingly, we 

%We have an action of $\GL_2(\Q)$ on $\Dist(M_{2\times2}(\Q),\tilde K_2(k_\infty))$.  To make this precise, if $\delta\in \Dist(M_{2\times2}(\Q),\tilde K_2(k_\infty))$ and $\alpha\in \GL_2(\Q)$ we define $\alpha^*\delta$ in $\Dist(M_{2\times2}(\Q),\tilde K_2(k_\infty))$ by 
%\bea
%(\alpha^* \delta)(\phi) = \delta((\alpha^{-1})^{*}\phi)
%\eea
%for an arbitrary element $\phi\in \cS(M_{2\times2}(\Q))$, where $(\alpha^*(\phi))(A) = \phi(\alpha A)$ 
%%and $(\alpha(\phi))(A) = \phi(A\alpha)$ 
%for any $A\in M_{2\times2}(\Q)$.




We provide a proof for Theorem \ref{modularsymbol}.
We start from the following important result on the Beilinson-Kato distribution ${{\bm\mu}^2}$.
\begin{proposition}\label{wdefined}
The Beilinson-Kato distribution ${{\bm\mu}^2} \in \Dist(M_{2\times2}(\Q),\tilde K^M_2(k_\infty))$ has the following properties.
\begin{itemize}
\item[(a)] (Invariance under $T(\Q)$)  For any $t$ in the standard torus $T(\Q)$ of diagonal matrices of $\GL_2(\Q)$, we have $t^*({{\bm\mu}^2}) = {{\bm\mu}^2}$;

\item[(b)] (The Manin Relations)   For $S = \begin{pmatrix}0&-1\\ 1&0\end{pmatrix}$ and $R = \begin{pmatrix}0&-1\\ 1&-1\end{pmatrix}$ we have
$$
(1 + (S^{-1})^*){{\bm\mu}^2} = (1+(R^{-1})^* +(R^{-2})^*){{\bm\mu}^2} = 0.
$$
\end{itemize}

\end{proposition}


\begin{proof}  Invariance under $T(\Q)$ is an immediate consequence of the fact that the Siegel distribution is invariant under the action of $\Q^\times$.   The first Manin relation follows from the skew symmetry of the Steinberg symbol.  The second is equivalent to Theorem  \ref{sres}. 
\end{proof}




For any pair of distinct rational cusps $r,s\in \PP^1(\Q)$ of the upper half-plane we choose $\alpha\in \GL_2(\Q)$ so that $\alpha(\infty) = r$ and $\alpha(0) = s$ and define ${\bm\xi}(r\to s) \in \Dist(M_{2\times2}(\Q), \tilde K^M_2(k_\infty))$ by 
\bea \label{BKD}
{\bm\xi}(r\to s)={\bm\xi}( \alpha(\infty)\to \alpha(0)) := (\alpha^{-1})^*{{\bm\mu}^2} =(\alpha^{-1})^* \left( {\bm\xi}(\infty \to 0)\right)
\eea
where ${{\bm\mu}^2}$ is the Beilinson-Kato distribution. We note first of all that this definition does not depend on the choice of $\alpha$.  Indeed, if $\beta \in \GL_2(\Q)$ is any other element satisfying $\beta(\infty) = r$ and $\beta(0) = s$
then $\beta =  \alpha t$ for some $t\in T(\Q)$.  So by \textit{(a)} of Proposition \ref{wdefined} we have $(\beta^{-1})^* {{\bm\mu}^2} = (t^{-1} \alpha^{-1} )^*{{\bm\mu}^2}= (\alpha^{-1})^* ({t^{-1}})^*{{\bm\mu}^2} = (\alpha^{-1})^* {{\bm\mu}^2}$.   


We have to prove that there is a unique $\GL_2(\Q)$-invariant modular symbol
$$
{\bm\xi} : \Delta_0 \lra \Dist(M_{2\times2}(\Q),\tilde K^M_2(k_\infty))
$$
with the property that 
$${\bm\xi}\biggl((0)-(\infty)\biggr) = {{\bm\mu}^2}. $$
%for which ${\bm\xi}\biggl((s)-(r)\biggr) = {\bm\xi}(r\to s)$ for every pair of distinct cusps $r,s\in \PP^1(\Q)$.


\begin{proof} [The proof of Theorem \ref{modularsymbol}] The uniqueness is clear and the $\GL_2(\Q)$-invariance of ${\bm\xi}$ is an immediate consequence of the fact (due to \eqref{BKD}) that for any $\gamma\in \GL_2(\Q)$ and any pair of distinct cusps $r,s \in  \PP^1(\Q)$ we have 
$$
{\bm\xi}(\gamma r\to \gamma s) = (\gamma^{-1})^*\left( {\bm\xi}(r\to s) \right). 
$$
Consider the following exact sequence of $\Z$-modules
\be
0\to M\to \Z[\GL_2(\Q)] &\to& \triangle_0 \to 0\\
\gamma &\mapsto& \gamma\cdot (\infty \to 0),
\ee
where $M$ is the module of Manin relations; $\triangle_0$ is generated by $(\infty\to 0)$ as a $\Z[\GL_2(\Q)]$-module and 
the annihilator of $(\infty \to 0)$ is the left ideal generated by the set
$$
\left\{ \d -1 : \d\in T(\Q) \right\} \cup \left\{ 1+S, 1+ R +R^2 \right\}.
$$
%where 
%$$
%\sigma:= \begin{pmatrix} 0 & -1 \\ 1 & 0 \end{pmatrix} \quad \text{ and } \quad \rho:= \begin{pmatrix} 0 & 1 \\ -1 & 1 \end{pmatrix}.
%$$
%
%If one can show 
%\be
%\tilde{\Phi}_{MK}(m)=0\  \text{ for every } m\in M, 
%\ee
%then $\tilde{\Phi}_{MK}$ induces a homomorphism $\triangle_0 \to D(V,\cK_2(V))$. 
So we only need to show that for any triple $r,s,t\in \PP^1(\Q)$ of distinct cusps we have 
$$
{\bm\xi}(r\to s) + {\bm\xi}(s\to r) = {\bm\xi}(r\to s) + {\bm\xi}(s\to t) + {\bm\xi}(t\to r) = 0.
$$
But there is a unique $\alpha \in \GL_2(\Q)$ that sends the ordered triple $(\infty, 0, 1)$ to $(r,s,t)$ and with that choice of $\alpha$ we have, by the Manin relations (\textit{(b)} of Proposition \ref{wdefined}),
$$
{\bm\xi}(r\to s) + {\bm\xi}(s\to r) = (\alpha^{-1})^*\bigl({\bm\xi}(\infty \to 0) + {\bm\xi}(0 \to \infty)\bigr) = (\alpha^{-1})^*(1+(S^{-1})^*){{\bm\mu}^2} = 0,
$$
and also
\be
{\bm\xi}(r\to s) + {\bm\xi}(s\to t) + {\bm\xi}(t\to r) 
&=& (\alpha^{-1})^*\bigl({\bm\xi}(\infty \to 0) + {\bm\xi}(0 \to 1) + {\bm\xi}(1 \to \infty)\bigr)\\
&=& (\alpha^{-1})^*(1+(R^{-1})^*+(R^{-2})^*){{\bm\mu}^2} = 0.
\ee
The theorem follows at once.
\end{proof}












\section{A generalization to $\GL_n(\Q)$-invariant modular symbols}\label{glngeneral}




\subsection{The statement for $\GL_n(\Q)$-invariant modular symbols}

%If $L$ is an affine lattice in $V$ and $\lambda \in V^\ast$, then $\lambda(L)$ is an affine lattice in $\Q$.
%\begin{definition} (wrong definition)
%For $\lambda\in V^\ast$ and $f \in \cS(V\times V)$, we define $\lambda_\ast(f) \in \cS(\Q^2)$ as follows:
%\bea \label{kd}
%\lambda_\ast(\sum_i[L_{1,i} \times L_{2,i}]) = \sum_i[\lambda(L_{1,i}) \times \lambda(L_{2,i})].
%\eea
%where $L_{1,i}$ and $L_{2,i}$ are affine lattices in $V$. Since each $f \in \cS(V\times V)$ is a finite linear combination of characteristic functions of affine lattices of the form $L_{1,i}\times L_{2,i}$ and $\lambda:V \to \Q$ is a linear map, the formula \eqref{kd} provides a well-defined (nope!) element in $\cS(\Q^2)$.
%\end{definition}



Let $\operatorname{Vect}$ be the category whose objects are finite dimensional vector spaces over $\Q$ and whose morphisms are $\Q$-linear isomorphisms.
Let $\operatorname{Ab}$ be the category of abelian groups.
Given any linear isomorphism $\varphi:V_1 \xrightarrow{\simeq} V_2$ of two vector spaces $V_1, V_2$ over $\Q$, we define
\bea \label{test}
\varphi^*: \cS(V_2) \xrightarrow{\simeq} \cS(V_1), \quad \varphi^* (f) (x) := f(\varphi(x)).
\eea
This implies that $\cS: \operatorname{Vect} \to \operatorname{Ab}$ is a contravariant functor.\footnote{Note that if one uses $\Q$-linear maps instead of $\Q$-linear isomorphisms in the category $\operatorname{Vect}$, then $\cS$ does not form a functor.} This functor has a nice additive property: if $V\simeq V_1\oplus V_2$, then $\cS(V)\simeq \cS(V_1) \otimes_\Z \cS(V_2)$.
If $M$ is a covariant functor from $\operatorname{Vect}$ to $\operatorname{Ab}$, then any linear isomorphism $\varphi:V_1\xrightarrow{\simeq} V_2$ induces an isomorphism of abelian groups:
\bea \label{indcov}
\varphi_*: \operatorname{\Dist}(V_1, M) \to \operatorname{\Dist}(V_2, M), \quad \varphi_*(\mu) (f_2)
:=\left( \mu(\varphi^*(f_2)) \right),
\eea
where $f_2 \in \cS(V_2)$, $\mu\in \operatorname{\Dist}(V_1, M)$.
Therefore, the assignments $V \mapsto \operatorname{\Dist}(V, M)$ and $\varphi \mapsto \varphi_*$ give us a covariant functor.





Let $V$ be an $n$-dimensional vector space and $V^*=\Hom(V,\Q)$ the dual vector space. Denote $\Aut_\Q(V)$ by $G$. Let $\cB$ be the set of ordered bases of $V^*$.
Now we give a definition of $G$-invariant modular symbols following \cite{AR}.
Let $G$ act on $V$ on the left. We define a right action of $G$ on $V^*$ as follows:
\be
(\g^* \l)(u ) = \l (\g u), \quad \g \in G, \l \in V^*, u \in V,
\ee
and $G$ acts on $(V^*)^n$ diagonally.
%Let $\cB$ be the set of ordered bases of $V^\ast$. 
Recall the distribution ${\bm\mu}^n$ in Definition \ref{glnSD}.
\begin{definition}\label{glnm}
For any $\ulambda = (\lambda_1,\lambda_2,\ldots,\lambda_n) \in (V^\ast)^n$, we define a distribution (using \eqref{indcov})
\bea \label{keydef}
\xi_\ulambda \in \Dist(V \times  V, \tilde K^M_n(k_\infty)), \quad 
%to be the distribution
% given by defining $\xi_\ulambda : \cS(V \times  V)\lra  \tilde K_n(k_\infty)$ by (using \eqref{indcov})
\xi_\ulambda := \begin{cases}
\phi(\ulambda)^{-1}_\ast ({\bm\mu}^n)  & \hbox{ \rm if\  }\ulambda\in \cB\\
\\
			 0 &\hbox{\rm otherwise,}
\end{cases}
\eea
where $\phi(\ulambda): V\times V\xrightarrow{\simeq} M_{n\times 2}(\Q)$ is the $\Q$-linear isomorphism induced from $\ulambda$:
\be
 V\times V \ni (v_1, v_2) \mapsto 
\begin{pmatrix}
\l_1(v_1) & \l_1(v_2) \\
\l_2(v_1) & \l_2(v_2) \\
\vdots & \vdots \\
\l_n(v_1) & \l_n(v_2) 
\end{pmatrix} \in M_{n\times 2}(\Q).
\ee
In particular, we have 
$$
\xi_\ulambda ([a+ rL_{\ulambda}]\otimes [b+rL_{\ulambda}]) := \begin{cases}
\{g_{\lambda_1(a)/r, \lambda_1(b)/r}, \ldots,g_{\lambda_n(a)/r, \lambda_n(b)/r} \}  & \hbox{ \rm if\  }\ulambda\in \cB\\
\\
			 0 &\hbox{\rm otherwise,}
\end{cases}
$$
where $a, b \in V$, $r \in \N$, and $L_\ulambda =\{ v \in V : \ulambda (v)=(\lambda_1(v), \cdots, \lambda_n(v)) \in \Z^n\}$. 
\end{definition}



\begin{definition}\cite{AR}\label{AR}
For any (contravariant) $G$-module $M$, the map 
\be
\xi:  (V^*)^n \to M
\ee
 is called a $G$-invariant modular symbol with values in $M$, if $\xi$ satisfies the following properties:
\begin{itemize}
\item $\xi$ is $G$-equivariant, i.e. for all $\gamma\in G$ and $\ulambda\in \cB$ we have 
\bea \label{one}
\xi({\gamma^\ast \ulambda}) = \gamma^\ast \left( \xi(\ulambda)\right).
\eea
%\item[\rm (2)] Letting $L_\ulambda \subseteq V$ be the lattice in $V$ whose dual lattice in $V^\ast$ is spanned by $\ulambda$, we have
%$$
%\int_{L(\ulambda)} d\xi_\ulambda = s(\lambda_1)\ast \cdots \ast s(\lambda_n) \in \cK_n(\ulambda).
%$$
%\medskip
\item For all $\ulambda \in \left(V^*\right)^n$ and $(t_1,\ldots, t_n)\in \left(\Q^\times\right)^n$ we have
%\footnote{In other words, $\xi(\gamma^\ast \ulambda)=\xi(\ulambda)$ for any $\gamma \in G(\ulambda)$; see \eqref{torus}.}
\bea \label{two}
\xi(t_1\lambda_1\ldots, t_n\lambda_n) = \xi(\lambda_1,\ldots,\lambda_n).
\eea
\item For any permutation $\sigma \in S_n$ and any $\ulambda = (\lambda_1,\ldots,\lambda_n) \in \left(V^*\right)^n$ we have
\bea \label{three}
\xi(\lambda_{\sigma(1)},\lambda_{\sigma(2)},\ldots, \lambda_{\sigma(n)}) = \sign(\sigma)\cdot 
\xi(\lambda_1, \lambda_2,\ldots,\lambda_n).
\eea
\item If $\ulambda\notin \cB$ (i.e. the $\lambda_1,\ldots,\lambda_n$ are linearly dependent in $V^*$), 
then 
\bea \label{five}
\xi(\ulambda)=0.
\eea
\item  For any $(n+1)$-tuple $(\lambda_0,\lambda_1,\ldots,\lambda_n) \in \left(V^*\right)^{n+1}$ we have
\bea \label{four}
\sum_{i=0}^n (-1)^i\xi(\lambda_0,\ldots,\hat \lambda_i,\ldots,\lambda_n) = 0.
\eea
\end{itemize}
\end{definition}




Recall that $G=\Aut_\Q(V)$. We equip $\tilde K^M_n(k_\infty)$ with a trivial contravariant $ G$-module structure, i.e. $G$ acts on $\Dist(V \times  V, \tilde K^M_n(k_\infty))$ as follows (see \eqref{distributionaction}):
\bea \label{taction}
\int \gamma^\ast f \cdot d\gamma^\ast\mu = \int f(x)d\mu(x), \quad (\g^\ast f)(v_1, v_2) = f(\g v_1, \g v_2), \quad x, (v_1,v_2) \in V \times V,
\eea
for any $\mu\in \Dist(V\times V,\tilde K^M_n(k_\infty))$, $f\in \cS(V\times V)$ and $ \gamma\in  G$.



\begin{theorem} \label{mainthm}
Let $\xi_V  : (V^\ast)^n  \lra  \Dist(V \times  V, \tilde K^M_n(k_\infty))$ be defined by
$$
\begin{array}{ccccc}
\xi_V  &: &(V^\ast)^n  &\lra&  \Dist(V \times  V, \tilde K^M_n(k_\infty))\\
&&\ulambda &\longmapsto& \xi_V(\ulambda)=\xi_\ulambda.\\
\end{array} 
$$
Then $\xi$ is a $G$-invariant modular symbol with values in $\Dist(V \times  V, \tilde K^M_n(k_\infty))$.
%Then $\xi$ is a modular symbol in the sense of Ash and Rudolph (\cite{AR}), i.e. we have
%\begin{enumerate}
%\item[\rm (1)] $\xi_V$ is $G$-equivariant, i.e. for all $\gamma\in G$ and $\ulambda\in \cB$ we have 
%\bea \label{one}
%\xi_V({\gamma^\ast \ulambda}) = \gamma^\ast \left( \xi_V(\ulambda)\right).
%\eea
%%\item[\rm (2)] Letting $L_\ulambda \subseteq V$ be the lattice in $V$ whose dual lattice in $V^\ast$ is spanned by $\ulambda$, we have
%%$$
%%\int_{L(\ulambda)} d\xi_\ulambda = s(\lambda_1)\ast \cdots \ast s(\lambda_n) \in \cK^M_n(\ulambda).
%%$$
%%\medskip
%\item[\rm (2)] For all $\ulambda \in \left(V^*\right)^n$ and $(t_1,\ldots, t_n)\in \left(\Q^\times\right)^n$ we have
%\bea \label{two}
%\xi_V(t_1\lambda_1\ldots, t_n\lambda_n) = \xi_V(\lambda_1,\ldots,\lambda_n).
%\eea
%\item[\rm (3)] For any permutation $\sigma \in S_n$ and any $\ulambda = (\lambda_1,\ldots,\lambda_n) \in \cB\subset\left(V^*\right)^n$ we have
%\bea \label{three}
%\xi_V(\lambda_{\sigma(1)},\lambda_{\sigma(2)},\ldots, \lambda_{\sigma(n)}) = \sign(\sigma)\cdot 
%\xi_V(\lambda_1, \lambda_2,\ldots,\lambda_n).
%\eea
%\item[\rm (4)] If $\ulambda\notin \cB$ (i.e. the $\lambda_1,\ldots,\lambda_n$ are linearly dependent in $V^*$), 
%then 
%\bea \label{five}
%\xi_V(\ulambda)=0.
%\eea
%\item[\rm (5)]  For any $(n+1)$-tuple $(\lambda_0,\lambda_1,\ldots,\lambda_n) \in \left(V^*\right)^{n+1}$ we have
%\bea \label{four}
%\sum_{i=0}^n (-1)^i\xi_V(\lambda_0,\ldots,\hat \lambda_i,\ldots,\lambda_n) = 0.
%\eea
%\end{enumerate}
\end{theorem}
The above theorem reduces to Theorem \ref{glnmodularsymbol} once one fixes an isomorphism between $(V^*, \ulambda)$ and $((\Q^n)^*, \{e_1^*, \cdots, e_n^*\})$:
\be
\phi(\ulambda)_\ast \left( \xi_V(\ulambda)\right)  = \bm\xi_n (e_1^*, \cdots, e_n^*)=\bm\mu^n, \quad \ulambda \in \cB,
\ee
using the $\Q$-linear isomorphism $\phi(\ulambda): V\times V\xrightarrow{\simeq} M_{n\times 2}(\Q)$ induced from $\ulambda$.

\begin{proposition}\label{cgcom}
The statement for $n=2$ of Theorem \ref{glnmodularsymbol} is equivalent to Theorem \ref{modularsymbol}.
\end{proposition}
\begin{proof}
Define the slope of a line (one-dimensional subspace) $\ell$ of $\Q^2$ to be
\be
\mathrm{slope}(\ell):=\left\{\frac{x}{y}\in \Q \sqcup \{\infty\} : (x,y) \in \ell \subset \Q^2 \right\}.
\ee
We identify
\be
\BP^1(\Q):=\{ \ell \subset \Q^2 :  \text{ $\ell$ is a line of $\Q^2$} \}
 &\xrightarrow{\sim}&   \Q \sqcup \{\infty\}\\
\ell &\mapsto& \mathrm{slope}(\ell).
\ee
Note that $\GL_2(\Q)$ acts on $\Q \sqcup \{\infty\}$:
\be
\begin{pmatrix}
a & b\\
c& d
\end{pmatrix} 
\cdot z = \frac{az+b}{cz+d}, \quad z \in \Q \sqcup \{\infty\}, \begin{pmatrix}
a & b\\
c& d
\end{pmatrix}  \in \GL_2(\Q).
\ee
Let $\cB_2$ be the set of ordered basis of $\Hom(\Q^2,\Q)$.
Then $(\Q^\times)^2$ acts on $\cB_2$:
\be
(x_1,x_2)^*(\lambda_1,\lambda_2)= (x_1^*\lambda_1, x_2^*\lambda_2),
\quad (x^* \lambda)(v)=\lambda(x\cdot v)
\ee
where $x_1, x_2,x \in \Q^\times$ and $(\lambda_1,\lambda_2) \in \cB_2$.
Here $[(\lambda_1,\lambda_2)]$ denotes the equivalence class of $(\lambda_1,\lambda_2)$ under the $(\Q^\times)^2$-action.
Then the assignment
\bea \label{assingn}
(\lambda_1,\lambda_2) \mapsto (\mathrm{slope}(\ker(\lambda_1))-\mathrm{slope}(\ker(\lambda_2)))
\eea
induces a map 
$$
\nu: \cB_2/(\Q^\times)^2 \to \Delta_0=\mathrm{Div}^0(\Q\sqcup \{\infty\})) 
$$
such that  $\nu (\g^* [(\lambda_1, \lambda_2)]) = \g^{-1}\cdot \nu([(\lambda_1,\lambda_2)])$ for $\g \in \GL_2(\Q)$. Note that if $(\lambda_1, \lambda_2)$ is not in $\cB_2$, $(\lambda_1, \lambda_2)$ goes to zero under $\nu$.
Unravelling the definitions, we have
\bea \label{rln}
 \bm\xi_2([\lambda_1, \lambda_2]) = (\phi(\ulambda)_\ast)^{-1} \bm\xi(\nu([\lambda_1, \lambda_2])).
\eea
The Manin relations for $\bm\xi$ correspond to \eqref{three} and \eqref{four} of $\bm\xi_2$. Thus we can construct $\bm\xi_2$ from $\bm\xi$.
Conversely, we can reverse the process to construct $\bm\xi$ from $\bm\xi_2$ by using the fact that $\Delta_0$ is generated by $(\infty \to 0)$ as a left $\Z[\GL_2(\Q)]$-module modulo the Manin relations.
%, which also satisfies \eqref{rln}.
%{\color{blue}Then it is clear that} {\color{red} 
Thus it follows from the uniqueness of both constructions that $\bm\xi_2$ is equivalent to $\bm\xi$.
\end{proof}


%\begin{corollary}
%For a non-zero covector $\l \in V^*$, the assignment 
%$$
%\Phi_\l (\xi_{\ulambda}) (\g_0, \cdots, \g_{n-1}) := \xi_{\ulambda}([\g_0^\ast\l, \cdots, \g_{n-1}^\ast \l])
%$$
%defines a homogeneous group cocycle of degree $n$, i.e. 
%\be
%[\Phi_\l (\xi_{\ulambda})]\in H^{n-1}\left(G, \Dist(V \times  V, \tilde K_n(k_\infty))\right).
%\ee
% If we change $\l$ by
%$\l'=\sigma^\ast \l$ for some %$\sigma\in G$, then 
%\be
%\Phi_\l (\xi_{\ulambda})(\g_0, \g_1, \cdots, \g_{n-1}) - \Phi_{\l'}(\xi_{\ulambda})(\g_0, \g_1,
%\cdots, \g_{n-1}) = (\d C_\sigma) (\g_0, \cdots, \g_{n-1}), \quad n \geq 2,
%\ee
%where
%\be
%C_\sigma (\g_1, \cdots, \g_{n-1})
%&:=&\sum_{i=2}^n(-1)^i \xi ([\g_1^\ast\l, \cdots, \g_{i-1}^\ast\l,  (\g_{i}\s )^\ast \l, (\g_{i+1}\sigma)^\ast  \l,\cdots,  (\g_{n-1}\sigma)^\ast \l ])\\
%(\d C_\sigma)  (\g_0, \cdots, \g_{n-1})&:=&\sum_{i=0}^{n-1} (-1)^i C_\sigma (\g_0,\cdots,
%\hat{\g}_i, \cdots, \g_{n-1}).
%\ee
%For $n=1$, we have $\Phi_\l(\xi_\ulambda) (\g_0)-\Phi_{\l'} (\xi_\ulambda)(\g_0)=0$.
%%For $k\geq 1$, $\Phi(\xi)_\l (\g_0)-\Phi(\xi)_{\s \l} (\g_0)= m - m|_\sigma$ (?)for some $m \in M$.
%\end{corollary}

%\begin{proof}
%This can be easily checked using \eqref{one} and \eqref{four}.
%The cochain property $\Phi_\l (\partial (\xi))= \d (\Phi_{\l}(\xi))$ follows from a straightforward computation.
%
%If we define $C_\sigma \in C^{k-2}(\GL(V), M), k \geq 2$, by the formula
%\be
%C_\sigma (\g_1, \cdots, \g_{k-1})
%:=\sum_{i=2}^k(-1)^i \xi ([\g_1 \cdot\l, \cdots, \g_{i-1}\cdot\l,  (\g_{i}\s )\cdot \l, (\g_{i+1}\sigma)\cdot  \l,\cdots,  (\g_{k-1}\sigma) \cdot \l ]),
%\ee
%a direct computation shows \ref{want}. In the case of $k=1$, we have
%$$
%\Phi_{\l}(\xi)(\g_0)-\Phi_{\sigma \l}(\xi)(\g_0) = \xi([\g_0 \l]) - \xi([\g_0 \sigma \l]) =0,
%%\xi[\l]|_{\g_0^{-1}} - \xi[\l]|_{\s^{-1} \g_0^{-1}} =0,
%$$
%since $\xi$ is a cocycle, i.e. $\xi \in Z^0(V,M):=\{f \in C^0(V,M) : f \text{ is a constant function on } C_0(V)  \}$.
%\end{proof}











\subsection{$G$-equivariance}

We will prove \eqref{one}.
Since $\phi(\gamma^\ast \ulambda) = \phi(\ulambda \circ \gamma)$,
we have 
\bea \label{middle}
\phi(\gamma^\ast \ulambda)^{-1} = \gamma^{-1} \circ \phi(\ulambda)^{-1}: M_{n\times 2}(\Q) \to V \times V.
\eea
If $\ulambda \notin\cB$, the $G$-equivariance is obvious.
For $\ulambda \in \cB$, the following computation gives \eqref{one}:
\be 
\xi_V({\gamma^\ast \ulambda}) (f)  &=& \left( (\phi(\gamma^\ast \ulambda)^{-1} )_\ast ({\bm\mu}^n) \right) (f) \quad \text{by} \quad \eqref{keydef}\\
&=& {\bm\mu}^n (f \circ \phi(\gamma^\ast \ulambda)^{-1})  \quad \text{by} \quad \eqref{test}, \eqref{indcov}\\
&=& {\bm\mu}^n ( \left(f \circ \gamma^{-1}\right) \circ \phi(\ulambda)^{-1})  \quad \text{by} \quad \eqref{middle}  \\
&=& (\phi(\ulambda)^{-1}_\ast({\bm\mu}^n)) (f \circ \gamma^{-1}) \quad \text{by} \quad  \eqref{test}, \eqref{indcov},\eqref{keydef} \\
 &=&\left( \gamma^\ast \left( \xi_V(\ulambda)\right)\right)(f) \quad \text{by} \quad \eqref{taction}.
\ee


\subsection{Other basic properties of $\xi_V$}
%Using the distribution relation of the Siegel units, 
For the proof of \eqref{two}, observe that \eqref{keydef} implies, for $a,b \in V$, $t_1, \ldots, t_n \in \Q^\times$ and $\ulambda \in \cB$, 
\be
\xi_\ulambda ([a+rL_{t_1\lambda_1, \ldots, t_n \lambda_n}]\times [b+rL_{t_1\lambda_1, \ldots, t_n \lambda_n}])=\{g_{\lambda_1(at_1/r), \lambda_1(bt_1/r)}, \ldots,g_{\lambda_n(at_n/r), \lambda_n(bt_n/r)}  \} 
\ee
where $
L_{t_1\lambda_1, \ldots, t_n \lambda_n}=\{x \in V : t_i\lambda_i(x) \in \Z, \quad i=1, \ldots,n \}.$
Since \eqref{keydef} also implies that
\be
\xi_{t_1\lambda_1, \ldots, t_n \lambda_n} ([a+rL_{t_1\lambda_1, \ldots, t_n \lambda_n}]\times [b+rL_{t_1\lambda_1, \ldots, t_n \lambda_n}])=\{g_{\lambda_1(at_1/r), \lambda_1(bt_1/r)}, \ldots,g_{\lambda_n(at_n/r), \lambda_n(bt_n/r)}  \},
\ee
\eqref{two} is proved.

\noindent Equation \eqref{three} follows from the antisymmetry of the Milnor $K_n$-group, and  \eqref{five} follows from the  definition of $\xi_V$.



\subsection{The cocycle relation}

%
%\begin{proposition}\label{cocyclep}
%Let $u_j, v_j \in \Q/\Z-\{0\}$ for $j=0, \ldots, n$.
%If $(u_0,v_0) = (\sum_{i=1}^n u_i, \sum_{i=1}^n v_i)$, then
%\bea \label{cocycle}
%\sum_{i=0}^n(-1)^i \{g_{u_0,v_0}, \ldots, \widehat{g_{u_i,v_i}}, \ldots, g_{u_n,v_n}  \} =0
%\eea
%in $\tilde K^M_n(k_\infty)$. 
%%$K^M_n(\cM(\overline{\Q}[\frac{1}{\Delta}])^\times)\otimes \Q$.
%\end{proposition}
%
%\begin{proof}
%The proof goes by induction on $n$. The case $n=1$ is obvious.
%%The case $n=2$ follows from \cite[Theorem 1.4]{Bru} or \cite{Gon}: 
%The case $n=2$ follows from the Manin relation (Theorem \ref{klem}): if $(-u_0 + u_1 +u_2, -v_0+v_1+v_2) =(0,0)$, then one can easily check that \eqref{cocycle} is equivalent to
%Corollary \ref{sres} 
%% \cite[(7)]{Bru} (or Corollary \ref{sres}).
%
%Now suppose $m\geq 2$ and the result is true for $n = m$. 
%We will prove the result for $n=m+1$. 
%
%Set $u_0=u_1+\cdots+u_{m+1}$ and $v_0=v_1+\cdots+v_{m+1}$. 
%Also, let $u_0'=u_1'+\cdots+ u_m'$ with $u_0'=u_0-u_{m+1}$, $u_1'=u_1, u_2'=u_2, \cdots, u_m'=u_m$ and $v_0'=v_1'+\cdots+ v_m'$ with $v_0'=v_0-v_{m+1}$, $v_1'=v_1, v_2'=v_2, \cdots, v_m'=v_m$. Then by the induction hypothesis
%\begin{eqnarray}\label{ih}
%\sum_{i=0}^m (-1)^i \left\{g_{u_0',v_0'}, \ldots, \widehat{g_{u_i',v_i'}}, \ldots, g_{u_m',v_m'} \right\} =0.
%\end{eqnarray}
%Multiplying \eqref{ih} on the left by $\{g_{u_{m+1}, v_{m+1}}\}$ and $\{g_{u_0,v_0}\}$ we get
%\begin{eqnarray*}
%A:=\sum_{i=0}^m (-1)^i \left\{g_{u_{m+1}, v_{m+1}}, g_{u_0',v_0'}, \ldots, \widehat{g_{u_i',v_i'}}, \ldots, g_{u_m',v_m'}  \right\} &=& 0,\\
%B:=\sum_{i=0}^m (-1)^i \left\{g_{u_0,v_0}, g_{u_0',v_0'}, \ldots, \widehat{g_{u_i',v_i'}}, \ldots, g_{u_m',v_m'}  \right\} &=& 0.
%\end{eqnarray*}
%By using the fact $u_0=u_0'+u_{m+1}$ and $v_0=v_0'+v_{m+1}$ we obtain
%$$
%\{g_{u_0',v_0'}, g_{u_{m+1}, v_{m+1}}\} - \{g_{u_0,v_0},g_{u_{m+1}, v_{m+1} }\} + \{g_{u_0,v_0},g_{ u_0', v_0'}\} =0.
%$$
%Therefore $A-B$ is equal to
%\be
%&&\left\{g_{u_{m+1},v_{m+1}}, g_{u_1,v_1}, \ldots, g_{u_m,v_m} \right\} + \sum_{i=1}^m (-1)^i \left\{g_{u_{m+1},v_{m+1}}, g_{u_0,v_0}, \ldots, \widehat{g_{u_i,v_i}}, \ldots, g_{u_{m},v_{m}}  \right\} \\
%&&- \left\{g_{u_0,v_0}, g_{u_1,v_1}, \ldots, g_{u_m,v_m} \right\},
%\ee
%which means (after multiplying $(-1)^m$) that
%$$
%\sum_{i=0}^{m+1} (-1)^i \left\{g_{u_0,v_0}, \ldots, \widehat{g_{u_i,v_i}}, \ldots, g_{u_{m+1},v_{m+1}} \right\} =0.
%$$
%This proves the assertion for $n=m+1$. The result follows by induction.
%\end{proof}
%

{
%\color{blue}



\begin{proposition}\label{cocyclep}
Let $\ba_j \in A(N)-(0,0)$ for $j=0, \ldots, n$.
If $\ba_0=\ba_1 + \cdots + \ba_n$, then
\bea \label{cocycle}
\sum_{i=0}^n(-1)^i \{g_{\ba_0}, \ldots, \widehat{g_{\ba_i}}, \ldots, g_{\ba_n}  \} =0
\eea
in $\tilde K^M_n(k_\infty)$. 
%$K^M_n(\cM(\overline{\Q}[\frac{1}{\Delta}])^\times)\otimes \Q$.
\end{proposition}

\begin{proof}
The proof goes by induction on $n$. The case $n=1$ is obvious.
%The case $n=2$ follows from \cite[Theorem 1.4]{Bru} or \cite{Gon}: 
The case $n=2$ follows from the Manin relation (Proposition \ref{klem}): if $-\ba_0 + \ba_1 +\ba_2 =0$, then one can easily check that \eqref{cocycle} is equivalent to
Theorem \ref{sres} by using the fact that $\{g_{-\ba_0},g_{\ba_i}\} \equiv \{g_{\ba_0}, g_{\ba_i}\}\equiv - \{g_{\ba_i}, g_{\ba_0} \}$ for $i=1,2$.
% \cite[(7)]{Bru} (or Corollary \ref{sres}).

Now suppose $m\geq 2$ and the result is true for $n = m$. 
We will prove the result for $n=m+1$. 
Set $\ba_0=\ba_1+\cdots+\ba_{m+1}$, which we rewrite as
 $\ba_0'=\ba_1'+\cdots+ \ba_m'$ where
\be
\ba_r':=
\begin{cases}
\ba_0 -\ba_{m+1}, & \text{ if } r=0 \\
\ba_r, &  \text{ if }  1 \leq r \leq m.
\end{cases}
\ee
%with $\ba_0'=\ba_0-\ba_{m+1}$, $\ba_1'=\ba_1, \ba_2'=\ba_2, \cdots, \ba_m'=\ba_m$. 

Then by the induction hypothesis for $n=m$ we have
\begin{eqnarray}\label{ih}
\sum_{i=0}^m (-1)^i \left\{g_{\ba_0'}, \ldots, \widehat{g_{\ba_i'}}, \ldots, g_{\ba_m'} \right\} =0.
\end{eqnarray}
Multiplying \eqref{ih} by $\{g_{\ba_{m+1}}\}$ and $\{g_{\ba_0}\}$, we get
\begin{eqnarray*}
A&:=&\sum_{i=0}^m (-1)^i \left\{g_{\ba_{m+1}}, g_{\ba_0'}, \ldots, \widehat{g_{\ba_i'}}, \ldots, g_{\ba_m'}  \right\} = 0,\\
B&:=&\sum_{i=0}^m (-1)^i \left\{g_{\ba_0}, g_{\ba_0'}, \ldots, \widehat{g_{\ba_i'}}, \ldots, g_{\ba_m'}  \right\} = 0.
\end{eqnarray*}
Using the fact that  $\ba_0=\ba_0'+\ba_{m+1}$ we obtain
\bea\label{stsymbol}
\{g_{\ba_0'}, g_{\ba_{m+1}}\} - \{g_{\ba_0},g_{\ba_{m+1} }\} + \{g_{\ba_0},g_{ \ba_0'}\} =0.
\eea
We now compute $A-B$:
\be
A-B &=&
\left\{g_{\ba_{m+1}}, g_{\ba_1}, \ldots, g_{\ba_m} \right\}  - \left\{g_{\ba_0}, g_{\ba_1}, \ldots, g_{\ba_m} \right\}\\
&&+ \sum_{i=1}^m (-1)^i \left( \{ g_{\ba_{m+1}}, g_{\ba_0'}\} - \{g_{\ba_0},g_{ \ba_0'}\}\right) \cdot \left\{g_{\ba_1}, \ldots, \widehat{g_{\ba_i}}, \ldots, g_{\ba_{m}}  \right\}
\\
&\stackrel{\eqref{stsymbol}}{=}&\left\{g_{\ba_{m+1}}, g_{\ba_1}, \ldots, g_{\ba_m} \right\} - \left\{g_{\ba_0}, g_{\ba_1}, \ldots, g_{\ba_m} \right\} +\sum_{i=1}^m (-1)^i \left\{g_{\ba_{m+1}}, g_{\ba_0}, \ldots, \widehat{g_{\ba_i}}, \ldots, g_{\ba_{m}}  \right\}\\
&=&(-1)^m  \sum_{i=0}^{m+1} (-1)^i \left\{g_{\ba_0}, \ldots, \widehat{g_{\ba_i}}, \ldots, g_{\ba_{m+1}} \right\}
\ee
Since $A-B=0$, we get the desired assertion for $n=m+1$.
Thus the proposition follows by induction.
\end{proof}
%}


We need to prove \eqref{four} of Theorem \ref{mainthm}. 
\begin{definition}
A set $\{v_1, \ldots, v_r\}$ of vectors in a $\Q$-vector space is said to be in general position if
the condition  
\be
c_1 v_1 + \cdots +c_r v_r = 0, \quad c_1, \ldots, c_r \in \Q
\ee
implies that either $c_1=\cdots =c_r =0$ or $c_1 \cdots c_r \neq 0$.
\end{definition}
This is equivalent to saying that any proper subset of $\{v_1, \ldots, v_r\}$ is linearly independent.
%A set of vectors in an $n$-dimensional vector space will be said to be in general position if every subset with no more than $n$ elements is linearly independent. Given $n+1$ vectors $v_0, \ldots, v_n$ in an $n$-dimensional vector space, which are in general position, there are non-zero scalars $u_0, \ldots, u_n \in \Q^\times$ such that $u_0v_0 +\cdots + u_nv_n = 0.$
\begin{lemma}\label{flem}
If $\{\l_0, \ldots, \l_n\} \subset V^*$ is in general position and $\l_0=\l_1 +\cdots + \l_n$, then the cocycle relation \eqref{four} holds.
\end{lemma}
\begin{proof}
For any subset $I:=\{i_1, \ldots, i_n\}\subset \{0, 1, \ldots, n\}$ of cardinality $n$, let $\l_I:=\{\l_{i_1}, \ldots, \l_{i_n}\} \subset V^*$.
Then $\l_I$ is linearly independent.
For any $u, v \in V$ and $r \in \Q^\times$, we have (see Definition \ref{glnm})
\bea \label{computation}
\xi_{\l_I} \left([u+rL_{\l_I}] \otimes [v+rL_{\l_I}]\right)
=\{g_{\ba_{i_1}}, \ldots, g_{\ba_{i_n}} \}, \quad
%=\{g_{(\l_{i_1}(u),\l_{i_1}(v))}, \ldots, g_{(\l_{i_n}(u), \l_{i_n}(v))} \}
\ba_{i}:=(\l_{i}(u)/r,\l_{i}(v)/r) \in A(N), \quad i \in I,
\eea
where we recall $L_{\l_I}=\{ x \in V : \l_{i}(x) \in \Z, \forall i \in I\}$ and $[u+r L_{\l_I}] \otimes [v+r L_{\l_I}] \in \cS(V)\otimes \cS(V)$ is the tensor product of the characteristic functions on the affine lattices $u+r L_{\l_I}$ and $v+r L_{\l_I}$.
%For simplicity of notation, put
%\be
%\ba_{i}:=(\l_{i}(u),\l_{i}(v)) \in A(N), \quad i \in I.
%\ee
Note that $\{ [u+rL_{\l_I}] \otimes [v+rL_{\l_I}] : u, v\in V, r \in \Q^\times\}$ generates $\cS(V)\otimes \cS(V)$ as an abelian group and $\l_0=\l_1 +\cdots + \l_n$ implies that $\ba_0=\ba_1+ \cdots + \ba_n$. Therefore, \eqref{four} reduces to Proposition \ref{cocyclep}.
\end{proof}

%Now, we use the induction.

We need to prove \eqref{four} when $\{\l_0, \ldots, \l_n\}$ is not in general position. To this end, we need the following lemma.
\begin{lemma}
There exists a subset $J =\{j_0, \ldots, j_k\} \subset \{0, \ldots, n\}$ such that $\l_J=\{\l_{j_0}, \ldots, \l_{j_k} \} \subset V^*$ is in general position and linearly dependent.
\end{lemma}
\begin{proof}
Let us consider the set $S$ of all subsets $\l_J$ of $\{\l_0, \ldots, \l_n\}$ which are linearly dependent. Then $S$ is a finite non-empty set: $\{\l_0, \ldots, \l_n\}$ is always linearly dependent. Thus there exists a minimal linearly dependent subset $\l_J=\{\l_{j_0}, \ldots, \l_{j_k} \} \subset \{\l_0, \ldots, \l_n\}$. 
Because of the minimality, every proper subset of $ \l_J$ is linearly independent, which implies that $\l_J$ is in general position.
%if $0 \in \{\l_0, \ldots, \l_n\}$, then $\{0\}$ is a linearly dependent set, and if $0 \notin \{\l_0, \ldots, \l_n\}$, then $\{\l_0, \ldots, \l_n\}$ is a linearly dependent set.
\end{proof}

%For the general case, we may choose a maximal subset $J=\{j_0, \cdots, j_k\}$ of $\{0, \ldots, n\}$ such that $k=\dim_\Q \sum_{j\in J} \Q \l_j$ and $\l_J=\{\l_{j_0}, \ldots, \l_{j_k}\}$ is in general position in the $k$-dimensional vector space $\sum_{j\in J} \Q \l_j \subset V^*$ and, for any subset $I \supsetneq J$, 
%$\l_I$ is not in general position in $\sum_{i\in I} \Q \l_i$.
%By using the homogeneity \eqref{two} and the anti-symmetry \eqref{three}, we may assume that 
%%there exists $k$ such that 
%$k=\dim_\Q \sum_{j=0}^k \Q \l_j $ and $\{\l_{0}, \ldots, \l_{k}\}$ is in general position in the $k$-dimensional vector space $\sum_{j=0}^k \Q \l_j \subset V^*$ with $\l_0=\l_1+ \cdots +\l_k$ and, for any $r \geq 1$, $\{\l_{0}, \ldots, \l_{k+r}\}$ is not in general position in $\sum_{j=0}^{k+r} \Q \l_j \subset V^*$.
Using this lemma, we choose a set $\{\l_{j_0}, \ldots, \l_{j_k} \} $ which is linearly dependent and in general position. Note that this set does not have to be unique and $k$ could be zero (this is the case if $0 \in \{\l_1, \ldots, \l_n\}$, and we get $\{\l_{j_0}\}=\{0\}$). By using the homogeneity \eqref{two} and the anti-symmetry \eqref{three}, we may assume that $\{\l_{0}, \ldots, \l_{k}\}$ is in general position and linearly dependent with relation $\l_0=\l_1+ \cdots +\l_k$ (if $k=0$, then $\l_0=0$).

%For the general case, we choose a subset $J=\{j_1, \ldots, j_k\}$ of $\{0, \ldots, n\}$ such that $\l_J=\{\l_{j_1}, \ldots, \l_{j_k}\}$ is in general position in the $k$-dimensional vector space $\sum_{j\in J} \Q \l_j \subset V^*$.
%
%For the general case, we choose a subset $J=\{0, \ldots, k\}$ of $\{0, \ldots, n\}$ such that $\l_J=\{\l_{0}, \ldots, \l_{k}\}$ is in general position in the $k$-dimensional vector space $\sum_{j\in J} \Q \l_j \subset V^*$ and $\l_{J\sqcup{\l_{k+1}, \ldots, \l_{k+r}}}$ is not in general position in the $k+r$-dimensional vector space $\sum_{j\in \sqcup{\l_{k+1}, \ldots,\l_{k+r}}} \Q \l_j$.
%
%let $k$ be the dimension of the vector space generated by $\l_0, \ldots, \l_n$.
%By using the homogeneity \eqref{two} and the anti-symmetry \eqref{three}, we may assume that $\l_0, \ldots, \l_{k}$ is in general position
%%, i.e. $\{\l_0, \ldots, \hat \l_i, \ldots, \l_{k}\}$ is linearly independent for any $0 \leq i \leq k$, 
%and $\l_0 = \l_1 + \cdots + \l_k$.
If $k=n$, then Lemma \ref{flem} implies \eqref{four}.
%\be
%\sum_{i=0}^n (-1)^k \xi_V (\l_0, \ldots, \hat \l_i, \ldots, \l_{n})=0.
%\ee
If $k < n$, then $\{ \l_0,  \ldots, \hat \l_i, \ldots, \l_{n}\}$ is linearly dependent for any $ i \geq k+1$, because $\{\l_0, \ldots, \l_k\}$ is linearly dependent. Thus \eqref{five} says that
\bea \label{firstn}
\xi_{(\l_0, \ldots, \hat \l_i, \ldots, \l_{n})}:=\xi_V(\l_0,  \ldots, \hat \l_i, \ldots, \l_{n} )=0, \quad \forall i \geq k+1.
\eea
Let $\l^{(i)}=\{\l_0,  \ldots, \hat \l_i, \ldots, \l_{n}  \}$. For $i \leq k$, then (the same notation as \eqref{computation})
\bea \label{anothercon}
\xi_{(\l_0,  \ldots, \hat \l_i, \ldots, \l_{n} )}([u+rL_{ \l^{(i)}}]\otimes [v+rL_{ \l^{(i)}} ])=
\begin{cases}
0, & \text{ if } \l^{(i)} \text { is linearly dependent,} \\
\{g_{\ba_{0}}, \ldots, \hat g_{\ba_i}, \ldots, g_{\ba_n} \} &  \text{ otherwise. }
\end{cases}
\eea
Since the vector space generated by $\{\l_0, \ldots, \hat\l_i, \ldots, \l_k\}$ is independent of $i \leq k$, the vector space generated by $\l^{(i)}$ is also independent of $i \leq k$.
%For $\l_i$ with $i \geq k+1$, either $\l_i \in V_k$ or $\l_i \notin V_k$ where $V_k$ is the vector space generated by $\l_0, \ldots, \l_k$. In other words, if $\l^{(i)}$ is linearly dependent (respectively, linearly independent) for some $0 \leq i \leq k$, then $\l^{(i)}$ is linear dependent (respectively, linearly independent) for any $i=0,1,\ldots, k$.


Since $\{\l_0, \ldots, \l_k\}$ is in general position and $\l_0=\l_1 +\cdots + \l_k$, a similar argument as in Lemma \ref{flem} implies that
\bea \label{secondn}
\sum_{i=0}^k (-1)^i  \{g_{\ba_{0}}, \ldots, \hat g_{\ba_i}, \ldots, g_{\ba_k} \}=0.
\eea

Finally, we compute
\be
&&\sum_{i=0}^n (-1)^i\xi_{(\lambda_0,\ldots,\hat \lambda_i,\ldots,\lambda_n)} ([u+rL_{ \l^{(i)}}]\otimes [v+rL_{ \l^{(i)}} ])\\
%&=&\sum_{i=0}^k (-1)^i \xi_{(\l_0, \ldots, \hat \l_i, \ldots, \l_{n})}([u+rL_{ \l^{(i)}}]\otimes [v+rL_{ \l^{(i)}} ]) + \sum_{i=k+1}^n (-1)^i \xi_{(\l_0, \ldots, \hat \l_i, \ldots, \l_{n})}([u+rL_{ \l^{(i)}}]\otimes [v+rL_{ \l^{(i)}} ]) \\
&\stackrel{\eqref{firstn}}{=}&\sum_{i=0}^k (-1)^i \xi_{(\l_0, \ldots, \hat \l_i, \ldots, \l_{n})}([u+rL_{ \l^{(i)}}]\otimes [v+rL_{ \l^{(i)}} ]) \\
&\stackrel{\eqref{anothercon}}{=}&\sum_{i=0}^k (-1)^i  \{g_{\ba_{0}}, \ldots, \hat g_{\ba_i}, \ldots, g_{\ba_n} \} \quad (\text {by using the comment below \eqref{anothercon}} \\
&=&\left( \sum_{i=0}^k (-1)^i  \{g_{\ba_{0}}, \ldots, \hat g_{\ba_i}, \ldots, g_{\ba_k} \}\right)\cdot\{g_{\ba_{k+1}}, \ldots, g_{\ba_n} \} 
\stackrel{\eqref{secondn}}{=}0.
\ee
From this, we conclude \eqref{four}.

%Since $(\l_0, \ldots, \l_k) \in (\sum_{j=0}^k \Q \l_j)^{k+1} \subset (V^*)^{k+1}$ is in general position in $\sum_{j=0}^k \Q \l_j$ and $\l_0=\l_1 +\cdots + \l_k$, a similar argument as in Lemma \ref{flem} implies that
%\be
%\sum_{i=0}^k (-1)^k \xi_V (\l_0, \ldots, \hat \l_i, \ldots, \l_{n})=0.
%\ee
%Thus we conclude \eqref{four}.
%\be
%\sum_{i=0}^n (-1)^k \xi_V (\l_0, \ldots, \hat \l_i, \ldots, \l_k, \l_{k+1}, \ldots, \l_{n})=0.
%\ee
%Since $\{\l_0, \ldots, \l_k\}$ is linearly dependent, \eqref{five} says that
%\be
%\xi_V (\l_0,  \ldots, \l_k, \l_{k+1},\ldots, \hat \l_j, \ldots, \l_{n})=0, \quad k+1 \leq j \leq n,
%\ee
%Thus we have \eqref{four}.


%Let $k$ be the dimension of the $\Q$-vector subspace in $V^\ast$ generated by $\lambda_0, \lambda_1, \ldots, \lambda_n$.
%%By the $G$-equivariance \eqref{one} and the homogeneity \eqref{two}, 
%%and the anti-symmetry \eqref{three}, 
%By the $G$-equivariance \eqref{one}, the homogeneity \eqref{two}, 
%and the anti-symmetry \eqref{three}, 
%checking the cocycle relation reduces to the case where $(\lambda_0,\lambda_1,\ldots,\lambda_n) \in \left(V^*\right)^{n+1}$ satisfies
%\be
%&&\lambda_1, \ldots, \lambda_k \text{ is linearly independent,}\\
%&&\lambda_0=\lambda_1 + \cdots + \lambda_m,
%\ee
%for some $ m \le k \le n$. 
%The cocycle relation \eqref{four} is equivalent to
%\bea \label{coc}
%%\xi_V(\lambda_1, \ldots, \lambda_n) + 
%\sum_{i=0}^m(-1)^i \xi_V(\lambda_0, \lambda_1, \ldots, \widehat{\lambda_i}, \ldots \lambda_m, \ldots, \lambda_n)=0,
%\eea
%because
%\be
%\sum_{i=m+1}^n (-1)^i \xi_V (\lambda_0, \lambda_1, \ldots \lambda_m, \ldots, \widehat{\lambda_i} \ldots \lambda_n)=0 \quad (\text{since } \lambda_0=\lambda_1 + \cdots + \lambda_m).
%\ee
%If $k$ is strictly smaller than $n$, then each term of the left hand side of \eqref{coc} is automatically zero due to the linear dependence of $\lambda_0, \ldots, \widehat{\lambda_i}, \ldots, \lambda_n$ for $i=0,\ldots, n$
%(since $k=\dim_\Q \sum_{i=0}^n \Q\cdot \lambda_i $).
%So we only have to deal with the case $k=n$.
%If $m=k=n$, then $\lambda_0, \lambda_1, \ldots, \lambda_n$ are in general position, i.e. $\lambda_0, \ldots, \widehat{\lambda_i}, \ldots \lambda_n$ are linearly independent for any $0 \le i \le n$, and we can check that the cocycle relation \eqref{four} reduces to \eqref{cocycle} in Proposition \ref{cocyclep}, using the construction of $\xi_V$ and its $G$-equivariance.
%Finally, if $m < k=n$, then the same proof as for Proposition \ref{cocyclep} gives the cocycle condition \eqref{coc}.


%Similarly, we may assume that $m=k=n$.
%Then, using the construction of $\xi_V$ and its $G$-equivariance, we can check that the cocycle relation \eqref{four} reduces to \eqref{cocycle} in Proposition \ref{cocyclep}.



%
%
%\section{Relationships with elliptic cocycles}
%
%\subsection{Elliptic cocycles}
%
%In \cite{Charo}, P. Charollois defined the following elliptic cocycle (see \cite{Zhang})
%\be
%\cE(\tau, \sigma, x, x')=\frac{1}{\sigma} \sum_{y,y' \in \Z^n/\sigma\Z^n} e(x\cdot y) \cdot \cK^M(\tau, x\sigma, \sigma^{-1} (x'+y\tau+y')).
%\ee
%for $\sigma \in M_n(\Z)$.


%
%\section{A detailed description for $n=2$ - to be deleted} \label{comparison}
%
%Let 
%\be
%\tilde \cK^M_n := \varinjlim_N \cK^M_n(Y(N)), \quad \cK^M_n(Y(N)):= K^M_n\left(Y(N) \otimes \Z[\frac{1}{N}]\right) \otimes \Q.
%\ee
%
%Let $V=\Q^2$ (as column vectors) on which $\GL_2(\Q)$ acts from the left.
%For $\ulambda=(\lambda_1, \lambda_2) \in V^\ast \times V^\ast$.
%We can view $\lambda_i:V\to \Q$ as a row vector $(u_i, v_i)$ in $\Q^2$ so that
%\bea \label{dual}
%\lambda_i (\begin{pmatrix} x \\ y \end{pmatrix}) =\begin{pmatrix} u_i v_i\end{pmatrix} \begin{pmatrix} x \\ y \end{pmatrix}= u_i x + v_i y, \quad i=1,2.
%\eea
%We define the $\GL_2(\Q)$-modular symbol
%\be
%{\bm\xi}:\Q^2 \times \Q^2 \to \Dist(V\times V, \tilde K^M_2) = \Dist(\M_{2\times2}(\Q), \tilde K^M_2)
%\ee
%by
%\be
%{\bm\xi} (\begin{pmatrix} r_1 \\ s_1 \end{pmatrix},\begin{pmatrix} r_2 \\ s_2 \end{pmatrix})=
%\xi_V (\lambda_1, \lambda_2)
%\ee
%where $\lambda_1 = (s_1, -r_1)$ and $\lambda_2= (s_2, -r_2)$ via the identification \eqref{dual}.
%Note that $\Ker(\lambda_i)$ is a one-dimensional line (if $\lambda_i \neq 0$) in $V=\Q^2$, which defines an element in $\BP^1(\Q)$ under the following bijection
%\bea \label{projective}
%\frac{\Q^2-\{(0,0)\}}{\sim} \simeq \BP^1(\Q)=\Q \sqcup \{\infty\}, \quad  [\begin{pmatrix} x \\ y \end{pmatrix}] \mapsto \frac{x}{y} 
%\eea
%where $\sim$ is the usual homothety relation.
%The restriction of ${\bm\xi}$ to $(\Q^2-\{(0,0)\})\times (\Q^2-\{(0,0)\})$ induces a map (with the same notation ${\bm\xi}$)
%\bea \label{mapproj}
%{\bm\xi}: \BP^1(\Q) \times \BP^1(\Q) \to \Dist(\M_{2\times2}(\Q), \tilde \cK_2).
%\eea
%by using the property \eqref{two}. Then \eqref{one} implies that
%\be
%{\bm\xi}(\gamma r, \gamma s)=  \gamma \cdot ( {\bm\xi}(r,s)), \quad r,s \in \BP^1(\Q),
%\ee
%where the action of $\GL_2(\Q)$ is given as follows:
%\bea \label{laction}
%\gamma \cdot r= \frac{ar + b}{cr + d}, \quad  \gamma =\begin{pmatrix} a & b \\ c & d \end{pmatrix} \in \GL_2(\Q), \quad r \in \BP^1(\Q),
%\eea
%\be
%(\gamma \cdot \mu) (f):= (\gamma^{-1\ast} \mu) (f) = \mu(\gamma^{\ast} (f))=\mu(f \circ \gamma), \quad
%\mu \in \Dist(M_{2\times2}(\Q), \tilde \cK_2), \quad f \in \cS(M_{2\times2}(\Q)).
%\ee
%
%
%
%
%Let $\Delta_0$ is the group of degree zero divisors on $\BP^1(\Q)$, where $\GL_2(\Q)$ acts on the left as in \eqref{laction}. 
%Then there is a linear map
%\bea \label{maps}
% \Z[\BP^1(\Q) \times \BP^1(\Q)] \to  \Delta_0, \quad \sum_i a_i [(p_{i1}, p_{i2})] \mapsto \sum_i a_i [ \{p_{i1}\} - \{p_{i2}\}] .
%\eea
%Moreover, it is well-known (\cite{Ma}) that $\Delta_0$ is generated by $\{0\} -\{\infty\}$ as $\Z[\GL_2(\Q)]$-module whose relations are given by the Manin relations.
%In fact, \eqref{three} and \eqref{four} give us the Manin relation, which implies that the map arising from the $\Z$-linear extension of ${\bm\xi}$ factors through \eqref{maps}. In other words, we have the following corollary of Theorem \ref{mainthm}, which is equivalent to Theorem \ref{mainone}.
%\begin{corollary}
%There is a $\GL_2(\Q)$-equivariant map induced from ${\bm\xi}$ (again with the same notation)
%\be
%{\bm\xi}: \Delta_0 \to \Dist(M_{2\times2}(\Q), \tilde \cK_2)
%\ee
%such that
%\be
%{\bm\xi} \left(\{0\}-\{\infty\}\right) \left(\left[\begin{pmatrix} a+r\Z & c+r\Z \\ b+s\Z & d + s\Z\end{pmatrix}\right]\right)
%=\{g_{a/r, c/r}, g_{b/s, d/s} \}
%\ee
%where $a,b,c,d \in \Q$ and $r,s \in \N$.
%\end{corollary}
%
%
%Let
%\be
%\cK_n(Y(m,M)):=K_n\left(Y(m,M)\otimes \Z[\frac{1}{mM}]\right) \otimes \Q 
%\ee
%for $m, M \geq 1$ such that $m+M \geq 5$.
%A simple computation confirms that
%\be
%{\bm\xi}\left(\begin{pmatrix} \frac{-u}{M} \\ \frac{s}{m} \end{pmatrix},  \begin{pmatrix} \frac{-v}{M} \\ \frac{t}{m} \end{pmatrix}\right) \left(\left[\begin{pmatrix} 1+m\Z & m\Z \\ M\Z & 1 + M\Z\end{pmatrix}\right]\right)
%=\{g_{\frac{s}{m}, \frac{u}{M}}, g_{\frac{t}{m},\frac{v}{M}} \}\in \cK_2(Y(m,M))
%\ee
%where $s,t,u,v \in \Z$ ($sv-tu=1$), $m, M \in \N$ such that $m+M \geq 5$ and $(\frac{s}{m},\frac{u}{M}), (\frac{t}{m}, \frac{v}{M})$ are linearly independent in $\Q^2$.
%
%For $r,s \in \BP^1(\Q)$, define the map ${\bm\xi}_{m,M}:\Delta_0 \to \cK_2(Y(m,M))$ by
%\be
%{\bm\xi}_{m,M} ((r)-(s)):= {\bm\xi}\left((r)-(s)\right)\left(\left[\begin{pmatrix} 1+m\Z & m\Z \\ M\Z & 1 + M\Z\end{pmatrix}\right]\right).
%\ee
%Let 
%\be
%\Gamma(m,M):=\{\begin{pmatrix} a & b \\ c & d \end{pmatrix} \in \SL_2(\Z) (or \GL_2(\Z)): a-1\equiv b \mod m,\quad  c \equiv d-1 \mod M\}
%\ee
%Since $\Gamma(m,M)$ fixes the test function $\left[\begin{pmatrix} 1+m\Z & m\Z \\ M\Z & 1 + M\Z\end{pmatrix}\right]$, ${\bm\xi}_{m,M}$ belongs to $\cK_2(Y(m,M))\otimes_\Z \Hom_{\Gamma(m,M)}(\Delta_0,\Z)$ where
%\be
%\Hom_{\Gamma(m,M)}(\Delta_0,\Z) =\{\psi: \Delta_0 \to \Z: \psi(\gamma D)= \psi(D), \gamma \in \Gamma(m,M), D\in \Delta_0\}.
%\ee
%
%%
%%\begin{corollary} (not yet clear, need to understand how to represent the cusps of $X(m,M)$)
%%%The map
%%%\be
%%% (\frac{-um}{sM}) - (\frac{-vm}{tM})=\left(\frac{-vm}{tM} \to \frac{-um}{sM}\right) \mapsto
%%% \{g_{\frac{s}{m}, \frac{u}{M}}, g_{\frac{t}{m},\frac{v}{M}} \}
%%%\ee
%%The map ${\bm\xi}_{m,M}$ induces a map (with the same notation)
%%\be
%%{\bm\xi}_{m,M}: H_1(X(m,M), cusps;\Z) \to \cK_2(Y(m,M)) 
%%\ee
%%\end{corollary}
%
%
%Let us represent the cusp $(r_1/r_2)$ of $\Gamma(m,M)$ where $r_1,r_2$ are integers such that $(r_1, r_2)=1$.
%%by $\begin{pmatrix} r_1 \\ r_2 \end{pmatrix}$ 
%Then the $\Gamma(m,M)$-cusp $(\frac{r_1}{r_2})$ is the same cusp as the $\Gamma(m,M)$-cusp $(\frac{r_1'}{r_2'})$ if and only if (this might not be true; it is definitely true if $m=M$)
%$r_1\equiv \pm r_1' \mod m$ and $r_2 \equiv \pm r_2' \mod M$.
%\begin{corollary}
%The assignment 
%\be
% (\frac{r_1}{r_2}) - (\frac{s_1}{s_2}) &\mapsto&
% \{g_{\frac{r_1}{m}, \frac{r_2}{M}}, g_{\frac{s_1}{m},\frac{s_2}{M}} \}
%\ee
%defines a well-defined map (denoted by $\Phi_{m,M}$)
%\be
%\Phi_{m,M}: H_1(X(m,M), cusps;\Z) \to \cK_2(Y(m,M)).
%\ee
%\end{corollary}
%
%Compare $\Phi_{m,M}$ to ${\bm\xi}_{m,M}$ which is described as
%\be
%{\bm\xi}_{m,M}: \Delta_0&\to& \cK_2(Y(m,M)) \\
% (\frac{-r_2/M}{r_1/m}) - (\frac{-s_2/M}{s_1/m}) &\mapsto&
% \{g_{\frac{r_1}{m}, \frac{r_2}{M}}, g_{\frac{s_1}{m},\frac{s_2}{M}} \}.
%\ee
%%\be
%%{\bm\xi}_{m,M}: \Delta_0 &\to& \cK_2(Y(m,M)) \\
%% (\frac{-um}{sM}) - (\frac{-vm}{tM}) &\mapsto&
%% \{g_{\frac{s}{m}, \frac{u}{M}}, g_{\frac{t}{m},\frac{v}{M}} \}.
%%\ee
%
%
%
%For $r \geq 2$, we have the following commutative diagram
%\[\begin{tikzcd}
%	{H^1(Y(mp^r,Mp^r),\Z)} & {H_1(X(mp^r,Mp^r),cusps;\Z)} & {} & {\cK_2(Y(mp^r,Mp^r))} \\
%	{H^1(Y(mp^{r-1},Mp^{r-1}),\Z)} & {H_1(X(mp^{r-1},Mp^{r-1}),cusps;\Z)} & {} & {\cK_2(Y(mp^{r-1},Mp^{r-1}))}
%	\arrow["cores", from=1-1, to=2-1]
%	\arrow[from=1-2, to=2-2]
%	\arrow["\simeq", from=1-1, to=1-2]
%	\arrow["\simeq", from=2-1, to=2-2]
%	\arrow["{\Phi_{mp^r,Mp^r}}", from=1-2, to=1-4]
%	\arrow["{\Phi_{mp^{r-1},Mp^{r-1}}}", from=2-2, to=2-4]
%	\arrow["Norm", from=1-4, to=2-4]
%\end{tikzcd}\]

%
%
%\section{A detailed analysis for $n=2$ (should be revised)} 
%Let 
%\be
%\tilde \cK_n := \varinjlim_N \cK_n(Y(N)), \quad \cK_n(Y(N)):= K_n\left(Y(N) \otimes \Z[\frac{1}{N}]\right) \otimes \Q.
%\ee
%Let
%\be
%\cK_n(Y(m,M)):=K_n\left(Y(m,M)\otimes \Z[\frac{1}{mM}]\right) \otimes \Q 
%\ee
%for $m, M \geq 1$ such that $m+M \geq 5$.
%Let $V=\Q^2$ and we fix a standard inner product on $V$.
%Then $V^\ast$ is canonically isomorphic to $V$ by this inner product.
%Under this identification $V^\ast \simeq \Q^2=V$, $\xi_V$ provides a $\GL_2$-modular symbol
%\be
%{\bm\xi}: \Q^2 \times \Q^2 \to \Dist(\Q^2 \times \Q^2, \tilde \cK_2).
%\ee
%According to the definition \eqref{keydef}, we have the following formula
%\be
%{\bm\xi} \left((\frac{s}{m},\frac{u}{M}), (\frac{t}{m}, \frac{v}{M})\right) 
%\left([(1+m\Z)\times M\Z] \otimes [m\Z \times (1+M\Z)]  \right)
%=\{ g_{\frac{s}{m}, \frac{u}{M}}, g_{\frac{t}{m}, \frac{v}{M}}\} \in \cK_2(Y(m,M)).
%\ee
%where $s,t,u,v \in \Z$ ($sv-tu=1$), $m, M \in \N$ such that $m+M \geq 5$ and $(\frac{s}{m},\frac{u}{M}), (\frac{t}{m}, \frac{v}{M})$ are linearly independent in $\Q^2$.
%Note that $g_{\frac{s}{m}, \frac{u}{M}}, g_{\frac{t}{m}, \frac{v}{M}} \in \cK_1(Y(m,M))$.
%We fix a bijection
%\bea \label{projective}
%\frac{\Q^2-\{(0,0)\}}{\sim} \simeq \BP^1(\Q)=\Q \sqcup \{\infty\}, \quad  [(x,y)] \mapsto \frac{y}{x} 
%\eea
%where $\sim$ is the usual homothety relation.
%The restriction of ${\bm\xi}$ to $(\Q^2-\{(0,0)\})\times (\Q^2-\{(0,0)\})$ induces a $\GL_2(\Q)$-equivariant map (with the same notation ${\bm\xi}$)
%\bea \label{mapproj}
%{\bm\xi}: \BP^1(\Q) \times \BP^1(\Q) \to \Dist(\Q^2 \times \Q^2, \tilde \cK_2).
%\eea
%by using the property \eqref{two}, where the action of $\GL_2(\Q)$ is translated as follows\footnote{Since $V^\ast$ has a contragradient action of $G \simeq \GL_2(\Q)$, i.e. $\gamma^\ast \lambda = \lambda \circ \gamma$, it is natural to view $\BP^1(\Q)$ as a right $\GL_2(\Q)$-module under our identification $V^\ast \simeq \Q^2$ via the standard inner product and \eqref{projective}.}:
%\bea \label{rightaction}
%\frac{y}{x} \cdot \gamma= \frac{b+d\frac{y}{x}}{a+c\frac{y}{x}}, \quad \gamma^\ast \mu (f) = \mu(\gamma^{-1\ast} (f))=\mu(f \circ \gamma^{-1}), \quad \gamma =\begin{pmatrix} a & b \\ c & d \end{pmatrix} \in \GL_2(\Q), \frac{y}{x} \in \BP^1(\Q),
%\eea
%where $f \in \cS(\Q^2\times \Q^2)$ and $\gamma^{-1}$ acts component-wise on $\Q^2\times \Q^2$ (as the column vector) from the left.
%
%Let $\Delta_0$ is the group of degree zero divisors on $\BP^1(\Q)$, where $\GL_2(\Q)$ acts on the right as in \eqref{rightaction}. 
%Then there is a linear map
%\bea \label{map}
% \Z[\BP^1(\Q) \times \BP^1(\Q)] \to  \Delta_0, \quad \sum_i a_i [(p_{i1}, p_{i2})] \mapsto \sum_i a_i [ \{p_{i2}\} - \{p_{i1}\}] .
%\eea
%Moreover, it is well-known (\cite{Ma}) that $\Delta_0$ is generated by $\{0\} -\{\infty\}$ as $\Z[\GL_2(\Q)]$-module whose relations are given by the Manin relations.
%In fact, \eqref{three} and \eqref{four} give us the Manin relation, which implies that the map arising from the $\Z$-linear extension of ${\bm\xi}$ factors through \eqref{map}. In other words, we have the following corollary of Theorem \ref{mainthm}.
%\begin{corollary}
%There is a $\GL_2(\Q)$-equivariant map induced from ${\bm\xi}$ (again with the same notation)
%\be
%{\bm\xi}: \Delta_0 \to \Dist(\Q^2 \times \Q^2, \tilde \cK_2)
%\ee
%such that
%\be
%{\bm\xi} (\{\infty\}-\{0\}) ([(a+r\Z)\times (b+s\Z)]\otimes [(c+r\Z) \times (d+s\Z)]
%=\{g_{a/r, c/r}, g_{b/s, d/s} \}
%\ee
%where $a,b,c,d \in \Q$ and $r,s \in \N$.
%\end{corollary}

%
%\newpage
%
%
%\section{Kummer's theory and $p$-adic Chern classes}
%
%
%%Let $\cM(\C^n)$ be the field of meromorphic functions on $\C^n$.
%%For each $\ulambda \in \cB$, we define a ring inclusion
%%\be
%%\cR(\ulambda)=\Q^{ab}[V^\ast]_{S(\ulambda)} \hookrightarrow \cM(\C^n), \quad e(t_r\lambda_r) \mapsto ( (z_1,\cdots,z_n) \mapsto e^{2 \pi i t_r z_r}), \quad  t_r \in \Q^\times, \ r=1, \cdots, n,
%%\ee
%%where $(z_1,\cdots, z_n)$ is the standard coordinate on $\C^n$.
%%Let $j_\ulambda:K_n(\cR(V)) \to K_n(\cM(\C^n))$ be the canonical map induced by this inclusion $\cR(\ulambda) \hookrightarrow \cM(\C^n)$.
%
%\subsection{The $p$-adic Chern class associated to the Dedekind-Eisenstein symbol}
%Let $\cF(V)$ be the quotient field of $\Q^{ab}[V^\ast]$. If $\ulambda \in \cB$ is given, there is a ring inclusion 
%\be
%\cF(V) \hookrightarrow \cM(\C^n), \quad e(t_r\lambda_r) \mapsto ( (z_1,\cdots,z_n) \mapsto e^{2 \pi i t_r z_r}), \quad  t_r \in \Q^\times, \ r=1, \cdots, n,
%\ee
%where $(z_1,\cdots, z_n)$ is the standard coordinate on $\C^n$.
%%From now on we assume that $n=1$. Let us fix a non-zero covector $\ulambda \in V^\ast$.
%We study the automorphism group $\Aut(\cF(V))$ in order to associate the $p$-adic Chern class.
%
%\begin{proposition}
%There is a commutative diagram of exact sequences:
%\[\begin{tikzcd}
%	1 & {\Aut(\cF(V)/\Q^{ab})} & {\Aut(\cF(V))} & {G_{\Q}^{ab}} & 1 & {} \\
%	\\
%	1 & {\SL_n(\hat \Z)} & {\GL_n(\hat \Z)} & {\hat \Z^\times} & 1. & {}
%	\arrow[from=1-1, to=1-2]
%	\arrow[from=3-1, to=3-2]
%	\arrow[from=3-2, to=3-3]
%	\arrow["\det", from=3-3, to=3-4]
%	\arrow[from=3-4, to=3-5]
%	\arrow[from=1-4, to=1-5]
%	\arrow[from=1-3, to=1-4]
%	\arrow[from=1-2, to=1-3]
%	\arrow[from=1-3, to=3-3]
%	\arrow[from=1-2, to=3-2]
%	\arrow["{\chi_{cyo}}", from=1-4, to=3-4]
%\end{tikzcd}\]
%%\[\begin{tikzcd}
%%	1 & {\Aut(\cF(V)/\Q^{ab})} & {\Aut(\cF(V))} & {G_{\Q}^{ab}} & 1 & {} \\
%%	\\
%%	1 & {\{\begin{pmatrix}  u & a\\ 0 & u^{-1}\end{pmatrix} : a\in \hat \Z, u\in \hat \Z^\times\}} & {\{ \begin{pmatrix}  u & a\\ 0 & v\end{pmatrix}: a\in \hat \Z, u,v \in \hat \Z^\times\}} & {\hat \Z^\times} & 1 & {}
%%	\arrow[from=1-1, to=1-2]
%%	\arrow[from=3-1, to=3-2]
%%	\arrow[from=3-2, to=3-3]
%%	\arrow["\det", from=3-3, to=3-4]
%%	\arrow[from=3-4, to=3-5]
%%	\arrow[from=1-4, to=1-5]
%%	\arrow[from=1-3, to=1-4]
%%	\arrow[from=1-2, to=1-3]
%%	\arrow[from=1-3, to=3-3]
%%	\arrow[from=1-2, to=3-2]
%%	\arrow["{\chi_{cyo}}", from=1-4, to=3-4]
%%\end{tikzcd}\]
%\end{proposition}
%
%
%
%
%
%Let 
%\be 
%\tilde \cF(V)^\times:=\{(x_n)_{n\geq 0}, x_n \in \cF(V)^\times , x_{n+1}^p=x_n, n \geq 0\}.
%%, \quad Z=\Q \otimes \tilde \cF(V)^\times.
%\ee
%Then $Z$ is equipped with action of $\Aut(\cF(V))$ and consider the $\Aut(\cF(V))$-module exact sequence
%\be
%0 \to \Q_p(1) \to  \tilde \cF(V)^\times &\to& \cF(V)^\times \to 0. \\
%(x_n)_{n \ge 0}&\mapsto& x_0
%\ee
%For each $1 \le i \le n$, let $pr_i:K_n(\cR(\ulambda)) \to \cR(\ulambda)^\times \subset \cF(V)^\times$ be the map defined by
%\be
%\{\alpha_1, \cdots, \alpha_n\} \mapsto \alpha_i.
%\ee
%Let $\tilde\xi_{\ulambda}$ be any lift of $\xi_{\ulambda} \in \Dist(V, \cK_n(V))$ to $\Dist(V, K_n(\cR(V)))$.
%We denote the image of $pr_i(\tilde \xi_{\ulambda}) \in \Dist(V,  \cF(V)^\times)$ under the Kummer map associated to the above sequence by
%\be
%\kappa_{\ulambda}^{(i)} \in H^1(\Aut(\cF(V)), \Dist\left(V, \Q_p(1)\right)).
%\ee
%In other words, choose any lift $\mu \in \Dist(V,\tilde \cF(V)^\times) $ of $pr_i(\tilde \xi_{\ulambda})$ and define
%\be
%\kappa_{\ulambda}^{(i)}(\gamma) =\gamma^\ast \mu -\mu, \quad \gamma \in \Aut(\cF(V)).
%\ee
%Now we define the $p$-adic Chern class associated to $\tilde \xi_{\ulambda}$:
%\be
%\kappa_{\ulambda}=\kappa_{\ulambda}^{(1)} \otimes \cdots \otimes \kappa_{\ulambda}^{(n)} \in H^n(\Aut(\cF(V)), \Dist(V, \Q_p(n))).
%\ee
%
%
%\subsection{The case $n=2$}
%
%We assume that $n=2$. 
%Let $\mathcal G=\Aut(\cF(V))$ and $\mathcal N=\Aut(\cF(V)/\Q^{ab})$ for simplicity so that $\mathcal G/\mathcal N =G_\Q^{ab}$:
%\[\begin{tikzcd}
%	{\cF(V)} \\
%	{\Q^{ab}} \\
%	\Q
%	\arrow[no head, from=1-1, to=2-1]
%	\arrow["{\mathcal G/ \mathcal N}", no head, from=2-1, to=3-1]
%	\arrow["{\small{\mathcal N=\Aut(\cF(V)/\Q^{ab})}}", no head, from=1-1, to=2-1]
%\end{tikzcd}\]
% Then the Hochschild-Serre spectral sequence gives us the exact sequence
%\be
%&0&\to H^1(\mathcal G/\mathcal N, \Dist(V, \Q_p(2))^{\mathcal N} ) \xrightarrow{\mathrm{inf}} H^1(\mathcal G, \Dist(V, \Q_p(2)))
%\xrightarrow{\mathrm{res}} H^1(\mathcal N, \Dist(V, \Q_p(2))) \\
%&\to& H^2 (\mathcal G/\mathcal N, \Dist(V, \Q_p(2))^{\mathcal N}) \to H^2(\mathcal G, \Dist(V, \Q_p(2)))^* \to H^1\left(\mathcal G/\mathcal N, H^1(\mathcal N, \Dist(V, \Q_p(2)))\right)
%\ee
%where 
%\be
%H^2(\mathcal G, \Dist(V, \Q_p(2)))^* =\Ker \left(H^2(\mathcal G, \Dist(V, \Q_p(2))) \xrightarrow{\mathrm{res}} H^2(\mathcal N, \Dist(V, \Q_p(2)))\right).
%\ee
%
%{\color{blue} Question: Does $k_\ulambda$ belong to $H^2(\mathcal G, \Dist(V, \Q_p(2)))^*$ ?
%For the moment, let us assume that it is true.}
%If we apply the map in the Hochschild-Serre spectral sequence to $\kappa_\ulambda \in H^2(\mathcal G, \Dist(V,\Q_p(2)))^*$, then we get an element
%\be
%c_{\ulambda} \in H^1(G_\Q^{ab}, H^1(\mathcal N, \Dist(V, \Q_p(2))).
%\ee
%
%
%\subsection{Towards to real quadratic fields}
%We continue to assume $n=2$. Here let $\cM(V)$ be the fraction field of the integral domain $\overline{\Q}[V^\ast]$ instead of $\Q^{ab}$.
%The exactly same construction would provide us a cohomology class
%\be
%C_\ulambda \in H^1(G_\Q, H^1(\frak N, \Dist(V, \Q_p(2)))).
%\ee
%where $\frak N=\Aut(\cM(V)/\overline{\Q})$.
%Let $F$ be a real quadratic field. Then there is a $\Q$-vector space isomorphism $\iota: F \simeq V$.
%
%{\color{blue} 
%Questions:
%1. We should choose a ``modular symbol " $\ulambda=(\lambda_1, \lambda_2)$ associated to $\{\infty, \gamma \infty\}$ where the matrix $\gamma$ corresponds to the fundamental unit of $F$. 
%
%2. We should make a correct test function in $\cS(V)$ using the information on $F$ via $\iota$. Then evaluate $C_\ulambda$ at it.
%
%3. We should construct a $\G_\Q$-equivariant map $H^1(\frak N, K(2)) \to \Q_p(1)$ where $K$ is some finite extension of $\Q_p$.
%
%4. Project $C_\ulambda$ to $H^1(G_\Q, \Q_p(1))$ and restrict to $G_H$ get a cohomology class $u_\ulambda \in H^1(G_H, \Q_p(1))$ using the above maps, where $H$ is an abelian extension of $F$. Examine how the unit (up to tensoring with $\Q$) associated to $u_\ulambda$ is related to the known units in $(\mathcal O_{H}[\frac{1}{p}])^\times \otimes \Q$, ``the Darmon-Dasgupta units" or ``the Gross-Stark units".
%
%5. Can we apply ``dual exponential map" to $C_\ulambda$ to get some sort of explicit reciprocity law?
%}
%
%
%
%
%\newpage
%



\section{Appendix} \label{appendix}
We provide a short explanation on the algebraicity of $\pi_N:\cE_N \to Y_N$ in \eqref{univell}.
\begin{theorem}
    Let $N \geq 3$. 
    \begin{enumerate}[(1)]
            \item There is a smooth affine curve $Y(N)$ over $\C$ such that $Y(N)(\C)$ is isomorphic to the complex manifold $Y_N$, and $Y_N$ is bijective with the set of isomorphism classes of triples $(E,e_1,e_2)$ where $E$ is an elliptic curve over $\C$ and $(e_1,e_2)$ is a pair of sections of $E$ over $\Spec(\C)$ which forms a $\Z/N\Z$-basis of $\Ker(N:E \to E)$ such that the Weil pairing of $e_1$ and $e_2$ is a fixed $N$-th root of unity.
        The smooth compactification $X(N)$ of $Y(N)$ is a smooth projective curve over $\C$ such that  $X(N)(\C)$ is bijective with $Y_N \sqcup \{\Gamma(N)\text{-cusps}\}$.
        \item There is a smooth affine curve $Y(N)$ over $\C$ such that $Y(N)(\C)$ is isomorphic to the complex manifold $Y_N$. %, and $Y_N$ is bijective with the set of isomorphism classes of triples $(E,e_1,e_2)$ where $E$ is an elliptic curve over $\C$ and $(e_1,e_2)$ is a pair of sections of $E$ over $\Spec(\C)$ which forms a $\Z/N\Z$-basis of $\Ker(N:E \to E)$.
        The smooth compactification $X(N)$ of $Y(N)$ is a smooth projective curve over $\C$ such that  $X(N)(\C)$ is bijective with $Y_N \sqcup \{\Gamma(N)\text{-cusps}\}$.
    \item   There is an elliptic surface $\overline \cE(N)$ over $X(N)$ (in the sense of \cite[Definition on p 202]{Silverman}) such that $\cE_N:=\tilde \fH/\tilde \Gamma(N)$ is isomorphic to the complex manifold $\cE(N)(\C)$ associated to $\cE(N):=Y(N) \times_{X(N)} \overline \cE(N)$.
    \end{enumerate}
\end{theorem}
\begin{proof}
{\it(1)}. Following \cite[Section 1.1]{Kat}, for $N\geq 3$, let $\cM(N)$ be the modular curve over $\Q$ of level $N$ without cusps, which represents the functor sending $S$ to the set of isomorphism classes of triples $(E, e_1,e_2)$ where $E$ is an elliptic curve over $S$
and $(e_1,e_2)$ is a pair of sections of $E$ over $S$ which forms a $\Z/N\Z$-basis of $\Ker (N : E\to E)$; the representability was proven in \cite{DR}; we also refer to \cite[Proposition 1.10.12, Theorem 4.7.0, Corollary 4.7.1, and Corollary 4.7.2]{KatzMazur} for details on representability.
Then $\cM(N)$ is a smooth irreducible affine curve over $\Q$. The total constant field of $\cM(N)$ (the field of all algebraic numbers in the affine ring $\cO(\cM(N))$) is not $\Q$, but is generated over $\Q$ by a primitive $N$-th root $\zeta_N$ of 1. Let $\overline\cM(N)$ be the smooth compactification of $\cM(N)$. By \cite[Section 1.8]{Kat}, the canonical map
\be
\nu:\fH \to \cM(N)(\C); \tau \mapsto (\C/\Z\tau+\Z,\tau/N\! \! \! \mod \Z\tau+\Z, 1/N\!\! \! \mod\Z\tau+\Z)
\ee
induces an isomorphism
\bea \label{cm}
(\Z/N\Z)^\times \times \fH/\Gamma(N) \xrightarrow{\simeq} \cM(N)(\C).
\eea
Note that the complex algebraic variety $\cM(N)_\C=\cM(N)\otimes_\Q \C$ is not irreducible; let $\cM(N)_\C^0$ be an irreducible component such that $\nu:\fH/\Gamma(N) \xrightarrow{\simeq} \cM(N)_\C^0(\C)$ is an isomorphism. 
Let $\overline\cM(N)_\C^{0}$ be the smooth compactification of $\cM(N)_\C^0$ corresponding to the compactification $\overline\cM(N)$ of $\cM(N)$.
We define
\begin{eqnarray*}
    Y(N):=\cM(N)_\C^0, \quad X(N):=\overline\cM(N)_\C^0.
\end{eqnarray*}
Then $X(N)$ and $Y(N)$ satisfy all the desired properties by the construction.

{\it(2)}. The existence of the elliptic surface $\overline \cE(N)$ over $X(N)$ (in the sense of \cite[Definition on p. 202]{Silverman}) was proved in \cite[Section 8, Theorem 11.5]{Kodaira} by Kodaira; this result was explicitly stated in \cite[Definition 4.1 and above]{Shioda} and \cite[Theorem 5.5]{Shioda}. 
In particular, this means that $\overline \cE(N)$ is an algebraic variety over $\C$.
Note that the smooth compactification $X(N)$ is isomorphic to $\Delta(N)$ in \cite[Theorem 5.5]{Shioda}; see \cite[(4.1) and below]{Shioda}.
%$\Delta_{\Gamma(N)}$ in \cite[Definition 4.1]{Shioda} (equivalently $\Delta(N)$ in \cite[Theorem 5.5]{Shioda}); see \cite[(4.1) and below]{Shioda}.
Let $\overline \cE(N)$ be the elliptic surface $B(N)$ over $\Delta(N)$ in \cite[Theorem 5.5]{Shioda} and $\cE(N)$ be $B'(N)$ in the proof of \cite[Theorem 5.5]{Shioda}; we have the following cartesian diagram.
%$B_{\Gamma(N)}$ over $\Delta_{\Gamma(N)}$ in \cite[Definition 4.1]{Shioda} (equivalently .
\[\begin{tikzcd}
	{\cE(N)}=B'(N) & {\overline \cE(N)=B(N)} \\
	{Y(N)} & {X(N)=\Delta(N).}
	\arrow[" ", from=1-1, to=1-2]
	\arrow["{}"', from=1-1, to=2-1]
	\arrow[from=1-2, to=2-2]
	\arrow[" ", from=2-1, to=2-2]
\end{tikzcd}\]

Then it was shown that $\cE(N)(\C) =B'(N)(\C)$ is isomorphic to $\cE_N:=\tilde \fH/\tilde \Gamma(N)$ in \cite[(5.7) and p. 39]{Shioda}.
%Let $(\Ell)$ be the category whose objects are elliptic curves $E/S$ over variable schemes and whose morphisms are cartesian squares of elliptic curves. (\cite[(4.1)]{KatzMazur})
\end{proof}

%For $N \geq 3$, let $\cE \to \cM(N)$ be the universal elliptic curve over $\Q$ that exists by the representability theorem with the full level $N$-structure mentioned in the proof above.
%Let $\cE^*$ denote the N\'eron model of $\cE$ over $\overline \cM(N)$. Then, in \cite[vi) on p.51]{Shioda}, it was stated that $\cE^*\times_{\Spec(\Q)} \Spec(\C)$ is isomorphic to 
%$\overline \cE(N)$ over $X(N)$ by using \cite[Theorem 5.5]{Shioda}: see \cite[A. Algebraic reformulation]{Shioda} for more details.
%$B(N)$ over $\Delta(N)$ by using \cite[Theorem 5.5]{Shioda}: see \cite[A. Algebraic reformulation]{Shioda} for more details.




We recall $_{N}\!\Theta(u,\tau)^{12} :=  \Theta(u,\tau)^{12 N^2}/\Theta(Nu,\tau)^{12}$ from Definition \ref{mero} and the function field $F_N=k_N(E(N))$ in Definition \ref{fn}.
    By Corollary \ref{periodic}, $_{N}\!\Theta(u,\tau)^{12}$ is a meromorphic function on $\cE_N\simeq\cE(N)(\C)$.
    \begin{proposition} \label{rationality}
        For $N\geq 3$, the function $_{N}\!\Theta(u,\tau)^{12}$ extends to a meromorphic function on $\overline \cE(N)(\C)$. Consequently, $_{N}\!\Theta(u,\tau)^{12} \in F_N$ and $_{N}\!\Theta_\ba(u,\tau)^{12} \in F_N$ for $\ba \in A(N)$.
    \end{proposition}
    \begin{proof}
    Let $\overline \pi_N:\overline \cE_N \to X_N$ be the universal generalized elliptic curve over $X_N$, where $\overline \cE_N=\overline \cE(N)(\C)$ and $X_N=X(N)(\C)$. 
    %Let $\overline \pi_N:\overline \cE_N \to X_N$ be the map induced from the map $\overline \cE(N) \to X(N)$ in the proof above, where $\overline \cE_N=\overline \cE(N)(\C)$ and $X_N=X(N)(\C)$. 
    Then $X_N=Y_N \sqcup \{\Gamma(N)-\text{cusps}\}$.
    For a $\Gamma(N)$-cusp $v$, let $C_v:=\overline \pi_N^{-1}(v)$ be the singular fiber, which is the standard N\'eron $N$-gon by \cite{Kodaira}; see \cite[Proposition 4.2, Theorem 5.2 (iii), and Example 5.4]{Shioda} for details. 
    Following \cite{Shioda} and \cite{Kodaira}, let $\theta_{v,i}$ $(0 \leq i \leq N-1)$ be the irreducible components of the divisor $C_v$ and take $\theta_{v,0}$ to be the unique component of $C_v$ containing the identity $o(v)$. 
    Let $C_v^\#$ denote the set of points of multiplicity one on $C_v$, which is shown to be a commutative algebraic group with identity $o(v)$ by \cite[Section 9]{Kodaira}; $\theta_{v,0}^\# =\theta_{v,0}\cap C_v^\#$ is the connected component of the identity. 
    Then $\theta_{v,0}^\#$ is a multiplicative subgroup and $C_v^\#/\theta_{v,0}^\# \simeq \Z/N\Z$.
    
    In order to show the meromorphicity of $_{N}\!\Theta(u,\tau)^{12}$ around $C_v$, we use the local chart of $\overline \cE_N$ around the singular fiber $C_v^\#$, given in the proof of \cite[Theorem 5.5]{Shioda}. We may assume that $v$ is the cusp at infinity $v_0$ without loss of generality, and we use the coordinate on $C_v^\#$ in \cite[p. 39]{Shioda}
    \begin{eqnarray*}
        q_N=e^{2 \pi i \tau/N}, \quad t= e^{2 \pi i u}, \quad (u, \tau) \in \C \times \fH;
    \end{eqnarray*}
    $C_v^\#$ is covered by $N$ open sets $W_i \quad (0 \leq i \leq N-1)$ of $\overline \cE_N$ with coordinates $((q_N, t))_i \quad (0 \leq i \leq N-1)$ (cf. \cite[p. 597-600]{Kodaira}). 
    Recall the following formula
        \begin{eqnarray*}
           _{N}\!\Theta(u,\tau)^{12}=q^{N^2-1} \frac{(t^{1/2} - t^{-1/2})^{12N^2} \prod_{n=1}^\infty \left(1 - q^nt\right)^{12N^2} \left(1 - q^nt^{-1}\right)^{12N^2}}{ (t^{N/2} - t^{-N/2})^{12} \prod_{n=1}^\infty \left(1 - q^n t^{N}\right)^{12}\left(1 - q^nt^{-N}\right)^{12}}
        \end{eqnarray*}
        for $t=e^{2 \pi i u}$ and $q=e^{2 \pi i \tau}$ and $(u,\tau) \in \C\times \fH$. 
         From this formula and the coordinate chart, it is clear that $_{N}\!\Theta(u,\tau)^{12}$ is well defined on $C_v$ where $[\tau] \to [i \infty]=v$ and so $q_N \to 0$, whereas $t\not\to 0, \infty$. Thus, $_{N}\!\Theta(u,\tau)^{12}$ extends to a meromorphic function on $\overline\cE_N=\cE_N \sqcup \bigsqcup_{v}C_v$. 
    \end{proof}
   




\begin{thebibliography}{PTW02} % '2nd argument contains the widest acronym'
\bibitem{AR}
Ash, Avner; Rudolph, Lee: The modular symbol and continued fractions in higher dimensions. Invent. Math. 55 (1979), no. 3, 241-250.


\bibitem{AS}
Ash, Avner; Stevens, Glenn: Modular forms in characteristic $\ell$ and special values of their $L$-functions. Duke Math. J. 53 (1986), no. 3, 849-868.
%"p-Adic deformations of arithmetic cohomology." preprint (2008).  See also 
%"p-adic deformations of cohomologyclasses of subgroups of $GL(n, \Z)$", Collect. Math. 48 (1997) 1-30.
%Or maybe its  "Cohomology of arithmetic groups and
%congruences between systems of Hecke eigenvalues," J. f. d. reine u. angew. Math.,365 (1986) 192-220.

\bibitem{BCG}
Bergeron, Nicolas; Charollois, Pierre; Garcia, Luis: Transgressions of the Euler class and Eisenstein cohomology of $\GL_N(\Z)$. Japanese Journal of Mathematics 15 (2020), pp. 311--379.

\bibitem{Bloch} 
Bloch, Spencer; The moving lemma for higher Chow groups. Journal of  Algebraic Geometry (1994), pp. 537-568.

\bibitem{Bru}
Brunault, Francois: Beilinson-Kato elements in $K_2$ of modular curves. Acta Arith. 134 (2008), no. 3, 283-298.

\bibitem{Bru2}
Brunault, Francois: On the $K_4$ group of modular curves. arXiv:2009.07614 (2022).
%{\color{red}  Note that, following the advice of the referee, I changed ``$K_4$-group" to ``$K_4$ group".}


%\bibitem{Charo} 
%P. Charollois:
% Le cocycle d'Eisenstein-Kronecker sur $\GL_N$, preprint.
\bibitem{Colmez}
Colmez, Pierre: La conjecture de Birch et Swinnerton-dyer $p$-adique. 
S\'eminaire Bourbaki: volume 2002/2003, expos\'es 909-923, Ast\'erisque, no. 294 (2004), Talk no. 919, pp. 251-319.
%S?eminaire BOURBAKI Juin 2003 55e ann?ee, 2002-2003, no 919, p. 251-319
%Asterisque

%\bibitem{delignesing}
%Deligne, P.: Equations differentielles a points singuliers reguliers. Lecture Notes in Mathematics 163, (1970).

\bibitem{DR}
Deligne, P.; Rapopport, M.: Les sch\'emas de modules des courbes elliptiques. Modular functions of one variable II. (Lect. Notes Math. vol. 349, pp. 143--316). Berlin-Heidelberg-New York: Springer 1973.


\bibitem{DS}
%Fred Diamond and Jerry Shurman. 
Diamond, Fred; Shurman, Jerry: 
A first course in modular forms, volume 228
of Graduate Texts in Mathematics. Springer-Verlag, New York, 2005.


\bibitem{FK}
Fukaya, Takako; Kato, Kazuya:
On conjectures of Sharifi
Kyoto J. Math. 64(2): 277--420 (May 2024)

\bibitem{Gon}
Goncharov, Alexander B.: Euler complexes and geometry of modular varieties. Geom. Funct. Anal. 17 (2008), no. 6, 1872-1914.

\bibitem{GL}
Goncharov, A. B.; Levin A. M.: Zagier's conjecture on $L(E,2)$. Inventiones mathematicae, Volume 132, pages 393-432, (1998).

%\bibitem{grothanalg} 
%Grothendieck, A.: On the de Rham cohomology of algebraic varieties. Publications mathématiques de l’I.H.E.S., tome 29 (1966), p. 95-103.

\bibitem{Kat}
Kato, Kazuya:  p-adic Hodge theory and values of zeta functions of modular forms. Ast\'erisque 295 (2004), 117-290.
%\bibitem{Hill}
%Hill, Richard: Shintani cocycles on $\GL_n$. Bull. Lond. Math. Soc. 39 (2007), no. 6, 993-1004.
\bibitem{KatzMazur}
Katz, Nick; Mazur, Barry: Arithmetic moduli of elliptic curves, Annals of Mathematics Studies 108, Princeton university press.

\bibitem{Kodaira}
Kodaira, K.: On compact analytic surfaces II--III, Ann. of Math., 77 (1963), 563--626; 78 (1963), 1--40.

\bibitem{Kriz}
Kriz, Sophia: On Weil reciprocity in motivic cohomology. Available at  krizsophie.github.io/Motives22102.pdf. 

\bibitem{Lang}
Lang, Serge: Elliptic functions. Springer New York, 1987.

\bibitem{KubertLang}
Kubert, Daniel S.; Serge Lang: Modular units. Springer New York, 1981.

\bibitem{Ma}
Manin, Ju. I.:
Parabolic points and zeta functions of modular curves. (Russian)
Izv. Akad. Nauk SSSR Ser. Mat. 36 (1972), 19-66.

\bibitem{Sch}
Scholl, Anthony J: An introduction to Kato's Euler systems. London Mathematical Society Lecture Note Series (1998), 379-460.
%\bibitem{Zhang}
%Zhang, Hao:
%Elliptic cocycles for $\GL_N(\Z)$ and Hecke operators.

%\bibitem{shimura}
%Shimura, Goro: Introduction to the arithmetic theory of automorphic functions. Vol. 1. Princeton university press, 1971.
\bibitem{SV}
Sharifi, Romyar; Venkatesh, Akshay: Eisenstein cocycles in motivic cohomology. Preprint 2023 arXiv:2011.07241.

\bibitem{Shioda}
Shioda, Tetsuji: On elliptic modular surfaces, J. Math. Soc. Japan, Vol. 24, No. 1, 1972.

\bibitem{Silverman}
Silverman, J.: Advanced topics in the arithmetic of elliptic curves, Graduate Texts in Mathematics (GTM, volume 151).

\bibitem{Sus}
Suslin, Andrey Aleksandrovich: Reciprocity laws and the stable rank of polynomial rings. Izvestiya Rossiiskoi Akademii Nauk. Seriya Matematicheskaya 43.6 (1979), 1394-1429.

\bibitem{Totaro} Totaro, Burt: Milnor $K$-theory is the simplest part of algebraic $K$-theory. $K$-theory 6.2 (1992), 177-189.

\bibitem{Weibel} Weibel, Charles: The K-book: an introduction to algebraic K-theory, Graduate Studies in Math. vol. 145, AMS, 2013; 
the online version is available at https://sites.math.rutgers.edu/~weibel/Kbook/Kbook.III.pdf.

\bibitem{Xu} Xu, Peter: Arithmetic Eisenstein theta lifts. PhD thesis, McGill University, 2023.

%\bibitem{BK}
%Barannikov, S., Kontsevich, M.:
%\newblock{\it Frobenius manifolds and formality of Lie algebras of polyvector fields},
%\newblock Internat. Math. Res. Notices {\bf 4} (1998) 201--215.
%
%%\bibitem{AdSp06}
%% Adolphson, A., Sperber S.: \emph{On the Jacobian ring of a complete intersection}, J. Algebra 304 (2006), no. 2, 1193-1227. 
%  
%\bibitem{Dim95}
%Dimca, A.: \emph{Residues and cohomology of complete intersections}, Duke Math. J., 78 (1995) No. 1 89--100.
%%
%%\bibitem{Dw60}
%%Dwork, B.: \emph{On the rationality of the zeta function of an algebraic variety.} Amer. J. Math. 82 (1960) 631--648.
%
%\bibitem {Gr69}
%Griffiths, P. A.: {\it On the periods of certain rational integrals. I, II,}
%Ann. of Math. (2) 90 (1969), 460-495; ibid. (2) 90 (1969), 496-541.
%%
%%\bibitem{Har75}
%%Hartshorne, R.:\emph{On the de Rham cohomology of algebraic varieties}, Pub. Math. L'I.H.E.S., tome 45 (1975) 5--99.
%%
%%\bibitem{Har77}
%%Hartshorne, R.:\emph{Algebraic Geometry}, Graduate Texts in Mathematics 52, Springer-Verlag (1977).
%
%\bibitem{KKP17}
%Kim, KH.; Kim,Y.; Park, J.: 
%{\it Inverse values of the modular $j$-invariant and homotopy Lie theory},  Proc. Amer. Math. Soc. 146 (2018), no. 8, 3295--3305.
%
%\bibitem{KKP}
%Kim, D.; Kim, Y.; Park, J.: {\it Differential Gerstenhaber-Batalin-Vilkovisky algebra for Calabi-Yau hypersurface complements}, Mathematika 64 (2018), no. 3, 637--651.
%
%
%\bibitem{Ko91}
%Konno, K.:  \emph{On the variational Torelli problem for complete intersections}, Comp. Math., 78 (1991), 271-296.
%
%\bibitem{KulWo}
%Kulkarni, R. S.; Wood, J. W.:\emph{Topology of Nonsingular Complex Hypersurfaces}, Adv. in Math., 25 (1980), 239--263.
%
%
%\bibitem{Mnev}
%Mnev, Patel:
%%\bibitem[PP16]{PP16}
%%Park, Jae-Suk; Park, Jeehoon:
%%\emph{Enhanced homotopy theory for period integrals of smooth projective hypersurfaces},
%%%CGP Preprint [CGP12004], 2012 \& 
%%arXiv:1310.6710. 
%
%\bibitem{PP16}
%Park, J.S.; Park, J.: {\it Enhanced homotopy theory for period Integrals of smooth projective hypersurfaces,} Communications
%in Number theory and Physics, Volume 10 (2016), Number 2, Pages 235-337.
%
%%\bibitem{PY17}Park, J.; Yhee, D.: {\it The Koszul-Tate type resolution for Gerstenhaber-Batalin-Vilkovisky algebras}, a preprint available at http://math.postech.ac.kr/~jeehoonpark/papers.html
%
%
%
%%\bibitem{Stacks} Stacks Project Authors: \emph{The Stacks Project}, \url{https://stacks.math.columbia,edu} (2018).
%
%
%
%\bibitem{Ter90}
%Terasoma, T.:\emph{Infinitesimal variation of Hodge structures and the weak global Torelli theorem for complete intersections}, Ann. of Math., 132 (1990), 213-235.
%
%\bibitem{Vois}
%Voisin, C.:\emph{Hodge Theory and Complex Algebraic Geometry I}, Cambridge studies in advanced mathematics 76 (2002).
%
%
%\bibitem{Wa87}
%Wahl, J.:
%\emph{The Jacobian algebra of a graded Gorenstein singularity},
%Duke Math. J. 55 (1987), 843--871.
%
%\bibitem{OW}
%Wells, R. O.:\emph{Differential Analysis on Complex Manifolds 3rd. ed.}, Graduate Texts in Mathematics 65 (2008).
%
%%\bibitem{Wi93} Witten, E.:\emph{Phases of N = 2 theories in two dimensions}, Nuclear Phys. B 403 (1993), no. 1-2, 159--222.

\end{thebibliography}













\end{document}
